\documentclass[UTF8, 12pt, fontset=fandol, oneside]{ctexart}
\usepackage{researchnotes}

\graphicspath{{math/数分习题/}}

\title{第三章\quad 一元微分学}
\author{}
\date{}

\begin{document}

\maketitle

\noindent\textbf{导读}\quad 本章是基础性内容，难度不大，适合各类读者，其中微分中值定理、Taylor 公式、导数应用是重点. 可微性问题主要针对数学院系学生.

本章主要讨论一元函数导数，微分中值定理，Taylor 公式，不等式与凸函数，以及导数的综合应用. 非数学院系的学生主要侧重于计算和应用.

\section*{§3.1\quad 导数}

\subsection*{*一、关于导数的定义与可微性}

\noindent\textbf{要点}\quad 若 $f(x)$ 在 $x_0$ 及其附近，由同一初等函数给出，则 $f(x)$ 在 $x_0$ 处可微，且 $f'(x_0)$ 可由初等函数的导数公式来计算. 若 $f(x)$ 在 $x_0$ 及其附近由不同的初等函数表示，则 $f(x)$ 在 $x_0$ 处的可微性与导数值，必须用导数定义来检验和计算. 因导数是用极限
\begin{equation}
  f'(x_0)=\lim_{x\to x_0}\frac{f(x)-f(x_0)}{x-x_0}
  =\lim_{h\to0}\frac{f(x_0+h)-f(x_0)}h
\tag{A}
\end{equation}
定义的，因此，可微性的证明就是极限 (A) 存在性的证明. 原则上，可以用第一章所介绍的方法. 要证明不可微，我们可以证明某个可微的必要条件不满足，例如 $f(x)$ 在 $x_0$ 处不连续，或 $f'_+(x_0)\ne f'_-(x_0)$，或 $x$ 以不同的方式趋向 $x_0$ 时，极限 (A) 取不同的值.

\noindent\textbf{$\star$例 3.1.1}\quad 1）设 $f(0)=0,\ f'(0)=A$，求极限 $\lim\limits_{n\to\infty}\sum\limits_{k=1}^n f\left(\dfrac{k}{n^2}\right)$；（武汉大学）

2）证明 Riemann 函数
\[
R(x)=
\begin{cases}
\dfrac1q, & \text{当 }x=\dfrac pq\text{ 时 }(p,q\text{ 为正整数，}\dfrac pq\text{ 为既约真分数}),\\
0, & \text{当 }x=0,1\text{ 和无理数时}
\end{cases}
\]
在 $[0,1]$ 上处处不可微.

\noindent\textbf{解}\quad 1）$\lim\limits_{x\to0}\left|\dfrac{f(x)-f(0)}x-f'(0)\right|=0\Rightarrow\dfrac{f(x)-f(0)}x-f'(0)=o(1)\ (x\to0)$
\[
  \Rightarrow f(x)=f(0)+f'(0)x+o(1)\cdot x.
\]
现在 $f(0)=0,\ f'(0)=A$，故
\[
  f\left(\frac{k}{n^2}\right)=f(0)+f'(0)\frac{k}{n^2}+o(1)\cdot\frac{k}{n^2}
  =\frac{k}{n^2}(A+o(1))\quad(n\to\infty).
\]
因此
\[
\sum_{k=1}^n f\left(\frac{k}{n^2}\right)
=\sum_{k=1}^n\frac{k}{n^2}(A+o(1))
=\frac{n+1}{2n}(A+o(1))\quad(n\to\infty).
\]
取极限得
\[
  \lim_{n\to\infty}\sum_{k=1}^n f\left(\frac{k}{n^2}\right)=\frac A2.
\]

\noindent\textbf{注}\quad 补充条件，有另解，见例 3.3.20.

2）因 $R(x)$ 在有理点不连续（例 2.1.1），故 $R(x)$ 在有理点上不可微，现只需证明无理点的情况.

设 $x_0\in(0,1)$ 为任一无理点，则 $x$ 沿无理点的点列 $\{x_n\}$ 趋向 $x_0$ 时，
\[
  \lim_{n\to\infty}\frac{R(x_n)-R(x_0)}{x_n-x_0}
  =\lim_{n\to\infty}\frac0{x_n-x_0}=0.
\]
现只要证明沿某个有理点的点列 $\{x'_n\}$ 趋向 $x_0$ 时，上述极限不为零即可.

因 $x_0$ 为无理数，可用无限不循环小数表示，$x_0=0.\alpha_1\alpha_2\cdots\alpha_n\cdots$. 截取前 $n$ 位小数，令 $x'_n=0.\alpha_1\alpha_2\cdots\alpha_n$，则 $\{x'_n\}$ 是趋向 $x_0$ 的有理数列. 注意 $\{\alpha_i\}_{i=1}^\infty$ 有无穷多项不为零！记第一个不为零的下标为 $N$，按 $R(x)$ 的定义，当 $n>N$ 时，有 $R(x'_n)=R(0.\alpha_1\alpha_2\cdots\alpha_n)>\dfrac1{10^n}$，故
\[
  \left|\frac{R(x'_n)-R(x_0)}{x'_n-x_0}\right|
  =\frac{R(0.\alpha_1\alpha_2\cdots\alpha_n)}
  {0.00\cdots0\alpha_{n+1}\alpha_{n+2}\cdots}\ge1.
\]
即 $\lim\limits_{n\to\infty}\dfrac{R(x'_n)-R(x_0)}{x'_n-x_0}\ne0$. 证毕.

\noindent\textbf{new 练习}\quad 假设函数 $h(x)$ 为处处不可导的连续函数（注\quad 此种函数有经典范例，见例 5.2.49），以此为基础构造连续函数 $f(x)$，使 $f(x)$ 仅仅在两点可导，并说明理由.（浙江大学）

\noindent\textbf{提示}\quad 例如，$f(x)=(x-a)(x-b)h(x)$.（可用导数定义和连续函数的运算性质说明！）

\noindent\textbf{例 3.1.2}\quad 证明：$f(x)=|x|^3$ 在 $x=0$ 处的三阶导数 $f'''(0)$ 不存在.

\noindent\textbf{证}\quad $f(x)=|x|^3=\operatorname{sgn}x\cdot x^3$，因此
\[
  f'(x)=3\operatorname{sgn}x\cdot x^2=3x|x|\quad(x\ne0\text{ 时}),
\]
\[
  f'(0)=\lim_{x\to0}\frac{f(x)-f(0)}{x-0}
  =\lim_{x\to0}\frac{\operatorname{sgn}x\cdot x^3}{x}=0.
\]
\[
  f''(x)=(3\operatorname{sgn}x\cdot x^2)'=6\operatorname{sgn}x\cdot x=6|x|\quad(x\ne0\text{ 时}),
\]
\[
  f''(0)=\lim_{x\to0}\frac{f'(x)-f'(0)}{x-0}
  =\lim_{x\to0}\frac{3x|x|}{x}=0.
\]
\[
  f'''_+(0)=\lim_{x\to0+}\frac{f''(x)-f''(0)}{x-0}
  =\lim_{x\to0+}\frac{6\operatorname{sgn}x\cdot x}{x}=6,
\]
而 $f'''_-(0)=-6$，所以 $f'''(0)$ 不存在.

\noindent\textbf{例 3.1.3}\quad 设 $p(x)=x,\ q(x)=1-x,\ f(x)$ 为多项式，$f(x)\ge p(x),\ f(x)\ge q(x)$（$\forall x\in(-\infty,+\infty)$），试证：$f\left(\dfrac12\right)>\dfrac12$.

\noindent\textbf{证}\quad 已知 $f\left(\dfrac12\right)\ge p\left(\dfrac12\right)=\dfrac12$，现证 $f\left(\dfrac12\right)>\dfrac12$. 事实上，若 $f\left(\dfrac12\right)=\dfrac12$，则 $x>\dfrac12$ 时，
\[
  \frac{f(x)-\dfrac12}{x-\dfrac12}
  \ge\frac{p(x)-\dfrac12}{x-\dfrac12}
  =\frac{x-\dfrac12}{x-\dfrac12}=1,
\]
所以
\[
  f'_+\left(\frac12\right)
  =\lim_{x\to(1/2)+}\frac{f(x)-\dfrac12}{x-\dfrac12}\ge1.
\]
$x<\dfrac12$ 时，
\[
  \frac{f(x)-\dfrac12}{x-\dfrac12}
  \le\frac{q(x)-\dfrac12}{x-\dfrac12}
  =\frac{1-x-\dfrac12}{x-\dfrac12}=-1,
\]
所以
\[
  f'_-\left(\frac12\right)
  =\lim_{x\to(1/2)-}\frac{f(x)-\dfrac12}{x-\dfrac12}\le-1.
\]
故 $f$ 在 $x=\dfrac12$ 处不可微，与 $f(x)$ 为多项式矛盾.

\noindent\textbf{例 3.1.4}\quad 设函数 $f(x)$ 在点 $x_0$ 的邻域 $I$ 内有定义. 证明：导数 $f'(x_0)$ 存在的充分必要条件是：存在这样的函数 $g(x)$，它在 $I$ 内有定义，在点 $x_0$ 连续，且使得在 $I$ 内成立等式
\begin{equation}
  f(x)=f(x_0)+(x-x_0)g(x),
\tag{1}
\end{equation}
又这时还有等式 $f'(x_0)=g(x_0)$.（武汉大学）

\noindent\textbf{证}\quad （必要性）已知 $f'(x_0)=\lim\limits_{x\to x_0}\dfrac{f(x)-f(x_0)}{x-x_0}=A$（存在），表明函数
\[
g(x)=
\begin{cases}
\dfrac{f(x)-f(x_0)}{x-x_0}, & \text{当 }x\ne x_0\text{ 时},\\
A, & \text{当 }x=x_0\text{ 时}
\end{cases}
\]
在 $x=x_0$ 处连续，式 (1) 成立，$f'(x_0)=g(x_0)$.

（充分性）$f'(x_0)=\lim\limits_{x\to x_0}\dfrac{f(x)-f(x_0)}{x-x_0}\stackrel{\text{式}(1)}{=}\lim\limits_{x\to x_0}g(x)\stackrel{g\text{在}x_0\text{处连续}}{=}g(x_0)$，故 $f'(x_0)$ 存在.

\noindent\textbf{例 3.1.5}\quad 设函数 $f(x)$ 在闭区间 $[a,b]$ 上连续，$f(a)=f(b)$，且在开区间 $(a,b)$ 内有连续的右导数：
\[
  f'_+(x)=\lim_{h\to0+}\frac{f(x+h)-f(x)}h\quad(a<x<b),
\]
试证：存在一点 $\xi\in(a,b)$，使得 $f'_+(\xi)=0$.（吉林大学）

\noindent\textbf{证}\quad 若 $f(x)\equiv$ 常数，则 $f'_+(x)\equiv0$，问题自明. 现设 $f(x)\ne$ 常数. 为了证明 $\exists \xi\in(a,b)$，使 $f'_+(\xi)=0$，只要证明 $\exists \alpha,\beta\in(a,b)$，使 $f'_+(\alpha)\le0,\ f'_+(\beta)\ge0$，那么由 $f'_+(x)$ 的连续性，便知存在 $\xi$ 使 $f'_+(\xi)=0$. 事实上，要找这样的 $\alpha,\beta$，只要找最大（小）值点即可. 因 $f(x)$ 在 $[a,b]$ 上连续，故在 $[a,b]$ 上必达最大、最小值. 而 $f(a)=f(b)$，所以最大、最小值至少有一个在内部达到. 设 $\alpha\in(a,b)$ 是 $f$ 的最大值点（内部达最小值类似讨论），于是
\[
  f'_+(\alpha)=\lim_{x\to\alpha+}\frac{f(x)-f(\alpha)}{x-\alpha}\le0.
\]
任取一点 $c:a<c<\alpha$，因 $f$ 在 $[c,\alpha]$ 上连续，$f$ 在 $[c,\alpha]$ 上也必有一点 $\beta<\alpha$ 达到最小值. 于是
\[
  f'_+(\beta)=\lim_{x\to\beta+}\frac{f(x)-f(\beta)}{x-\beta}\ge0.
\]
如此我们即达到了目的.

\noindent\textbf{$\star$例 3.1.6}\quad 设 $f(x)$ 在 $x=x_0$ 处可微，$\alpha_n<x_0<\beta_n\ (n=1,2,\cdots),\ \lim\limits_{n\to\infty}\alpha_n=\lim\limits_{n\to\infty}\beta_n=x_0$，证明：
\[
  \lim_{n\to\infty}\frac{f(\beta_n)-f(\alpha_n)}{\beta_n-\alpha_n}=f'(x_0).
\]
（湖北大学）

\noindent\textbf{证}\quad 首先，容易看到
\begin{equation}
\begin{aligned}
\frac{f(\beta_n)-f(\alpha_n)}{\beta_n-\alpha_n}
={}&\frac{f(\beta_n)-f(x_0)+f(x_0)-f(\alpha_n)}{\beta_n-\alpha_n}\\
={}&\frac{\beta_n-x_0}{\beta_n-\alpha_n}\frac{f(\beta_n)-f(x_0)}{\beta_n-x_0}
  -\frac{\alpha_n-x_0}{\beta_n-\alpha_n}\frac{f(\alpha_n)-f(x_0)}{\alpha_n-x_0}.
\end{aligned}
\tag{1}
\end{equation}
若记 $\lambda_n=\dfrac{\beta_n-x_0}{\beta_n-\alpha_n}$，则 $\dfrac{x_0-\alpha_n}{\beta_n-\alpha_n}=1-\lambda_n$，且 $0<\lambda_n<1,\ 0<1-\lambda_n<1$，式 (1) 可改写成
\[
  \frac{f(\beta_n)-f(\alpha_n)}{\beta_n-\alpha_n}
  =\lambda_n\frac{f(\beta_n)-f(x_0)}{\beta_n-x_0}
  +(1-\lambda_n)\frac{f(\alpha_n)-f(x_0)}{\alpha_n-x_0}.
\]
但 $f'(x_0)=\lambda_n f'(x_0)+(1-\lambda_n)f'(x_0)$，故易知 $\forall\varepsilon>0,\exists N>0$，当 $n>N$ 时，
\[
\begin{aligned}
\left|\frac{f(\beta_n)-f(\alpha_n)}{\beta_n-\alpha_n}-f'(x_0)\right|
\le{}&\lambda_n\left|\frac{f(\beta_n)-f(x_0)}{\beta_n-x_0}-f'(x_0)\right|\\
&+(1-\lambda_n)\left|\frac{f(\alpha_n)-f(x_0)}{\alpha_n-x_0}-f'(x_0)\right|\\
<&\lambda_n\varepsilon+(1-\lambda_n)\varepsilon=\varepsilon.
\end{aligned}
\]
原极限获证.

\noindent\textbf{$\star$例 3.1.7}\quad 设函数 $f(x)$ 在 $x=0$ 处连续，并且 $\lim\limits_{x\to0}\dfrac{f(2x)-f(x)}x=A$，求证：$f'(0)$ 存在，并且 $f'(0)=A$.（中国科学院）

\noindent\textbf{证}\quad （目标在于证 $\lim\limits_{x\to0}\dfrac{f(x)-f(0)}x=A$）因已知 $\lim\limits_{x\to0}\dfrac{f(2x)-f(x)}x=A$，即 $\forall\varepsilon>0,\exists\delta>0$，当 $|x|<\delta$ 时，有
\[
  A-\frac\varepsilon2<\frac{f(2x)-f(x)}x<A+\frac\varepsilon2.
\]
特别，取 $x_k=\dfrac{x}{2^k}\ (k\in\mathbb N)$，上式亦成立. 故有
\[
  \frac1{2^k}\left(A-\frac\varepsilon2\right)
  <\frac{f\left(\dfrac{x}{2^{k-1}}\right)-f\left(\dfrac{x}{2^k}\right)}x
  <\frac1{2^k}\left(A+\frac\varepsilon2\right),
\]
$k=1,2,\cdots,n$. 将此 $n$ 式相加，注意
\[
  \sum_{k=1}^n\left[f\left(\frac{x}{2^{k-1}}\right)-f\left(\frac{x}{2^k}\right)\right]
  =f(x)-f\left(\frac{x}{2^n}\right),\quad
  \sum_{k=1}^n\frac1{2^k}=1-\frac1{2^n},
\]
有
\[
  \left(1-\frac1{2^n}\right)\left(A-\frac\varepsilon2\right)
  <\frac{f(x)-f(x_n)}x
  <\left(1-\frac1{2^n}\right)\left(A+\frac\varepsilon2\right).
\]
再令 $n\to\infty$，取极限，这时 $x_n=\dfrac{x}{2^n}\to0$，而 $f$ 在 $0$ 处连续，$\lim\limits_{n\to\infty}f(x_n)=f(0)$，故
\[
  A-\frac\varepsilon2\le\frac{f(x)-f(0)}x\le A+\frac\varepsilon2.
\]
亦即 $\left|\dfrac{f(x)-f(0)}x-A\right|\le\dfrac\varepsilon2<\varepsilon$，$f'(0)$ 存在且 $f'(0)=A$.

\noindent\textbf{例 3.1.8}\quad 设 $f(x)$ 在 $[a,b]$ 上可微，试证：$f'(x)$ 在 $[a,b]$ 上连续的充要条件是 $f(x)$ 在 $[a,b]$ 上一致可微，即：$\forall\varepsilon>0,\exists\delta>0$，当 $0<|h|<\delta$ 时，有
\[
  \left|\frac{f(x+h)-f(x)}h-f'(x)\right|<\varepsilon
\]
对一切 $x\in[a,b]$ 成立.（北京师范大学）

\noindent\textbf{证}\quad $1^\circ$（必要性）因 $f'(x)$ 在 $[a,b]$ 上连续，故一致连续，即 $\forall\varepsilon>0,\exists\delta>0$，当 $x',x''\in[a,b],\ |x'-x''|<\delta$ 时，便有 $|f'(x')-f'(x'')|<\varepsilon$. 由此当 $0<|h|<\delta$ 时，$\forall x\in[a,b]$，有
\[
  \left|\frac{f(x+h)-f(x)}h-f'(x)\right|
  =|f'(\xi)-f'(x)|\quad(\xi\text{ 在 }x\text{ 与 }x+h\text{ 之间})
  <\varepsilon.
\]

$2^\circ$（充分性）已知 $\forall\varepsilon>0,\exists\delta>0$，当 $0<|h|<\delta$ 时，$\forall x\in[a,b]$，
\[
  \left|\frac{f(x+h)-f(x)}h-f'(x)\right|<\varepsilon.
\]
因此，$\forall x_0\in[a,b]$，当 $0<|h|<\delta$ 时，只要 $x_0+h\in[a,b]$，便有
\[
\begin{aligned}
|f'(x_0+h)-f'(x_0)|
=&\left|f'(x_0+h)-\frac{f(x_0+h)-f(x_0)}h
+\frac{f(x_0+h)-f(x_0)}h-f'(x_0)\right|.
\end{aligned}
\]
\[
\le\left|f'(x_0+h)-\frac{f(x_0+h-h)-f(x_0+h)}{-h}\right|
+\left|\frac{f(x_0+h)-f(x_0)}h-f'(x_0)\right|<2\varepsilon.
\]
所以 $f'(x)$ 在 $x_0$ 处连续. 由 $x_0$ 的任意性，知 $f'(x)$ 在 $[a,b]$ 上连续.

\subsection*{$\star$二、高阶导数与 Leibniz 公式}

\noindent\textbf{a.\quad 先拆项再求导}

\noindent\textbf{要点}\quad 有些式子不易直接求高阶导数，当拆项以后，变成易于求高阶导数的一些基本形式之和，便立即可以直接求导. 基本形式主要有
\[
  (x^k)^{(n)}=k(k-1)\cdots(k-n+1)x^{k-n}\quad(n\le k),
\]
\[
  (e^x)^{(n)}=e^x,\quad (a^x)^{(n)}=a^x(\ln a)^n,
\]
\[
  (\ln x)^{(n)}=(-1)^{n-1}(n-1)!x^{-n},
\]
\[
  (\sin x)^{(n)}=\sin\left(x+\frac{n\pi}2\right),\quad
  (\cos x)^{(n)}=\cos\left(x+\frac{n\pi}2\right).
\]
并特别注意 $[f(ax+b)]^{(n)}=a^n f^{(n)}(ax+b)$，因子 $a^n$ 不要漏掉.

\noindent\textbf{例 3.1.9}\quad 计算 $y^{(n)}$，设

（1）$y=\sin^6x+\cos^6x$；（中南大学）

（2）$y=\sin ax\sin bx$；

（3）$y=\dfrac{x^2+x+1}{x^2-5x+6}$.

\noindent\textbf{提示}\quad （1）$y=\left(\dfrac{1-\cos2x}{2}\right)^3+\left(\dfrac{1+\cos2x}{2}\right)^3=\dfrac58+\dfrac38\cos4x$.

（2）先积化和差.

（3）先分项，$y=1-\dfrac7{x-2}+\dfrac{13}{x-3}$.

\noindent\textbf{b.\quad 直接使用 Leibniz 公式}

\noindent\textbf{要点}\quad 把要求导的函数写成两项相乘，然后直接应用 Leibniz 公式
\[
  (u\cdot v)^{(n)}=\sum_{k=0}^n C_n^k u^{(k)}v^{(n-k)}.
\]

\noindent\textbf{*例 3.1.10}\quad 设 $f_1(x),f_2(x)$ 在 $x_0$ 及其附近有定义，在 $x_0$ 处有直到 $n$ 阶导数，记
\[
  N(f)=\sum_{k=0}^n\frac1{k!}|f^{(k)}(x_0)|.
\]
试证明：$N(f_1f_2)\le N(f_1)N(f_2)$.（北京大学）

\noindent\textbf{证}\quad 按 $N$ 的定义，
\[
\begin{aligned}
N(f_1f_2)
&=\sum_{k=0}^n\frac1{k!}|(f_1f_2)^{(k)}(x_0)|\\
&\le\sum_{k=0}^n\frac1{k!}\sum_{j=0}^k C_k^j
|f_1^{(j)}(x_0)|\,|f_2^{(k-j)}(x_0)|\\
&=\sum_{k=0}^n\sum_{j=0}^k
\frac1{j!(k-j)!}|f_1^{(j)}(x_0)|\,|f_2^{(k-j)}(x_0)|\\
&\le\sum_{j=0}^n\frac1{j!}|f_1^{(j)}(x_0)|
\cdot\sum_{k=0}^n\frac1{k!}|f_2^{(k)}(x_0)|\\
&=N(f_1)\cdot N(f_2).
\end{aligned}
\]
此式（第三行）$\le$（第四行）是因为：第四行乘开后共有 $(n+1)^2$ 项，可排成一方阵，而第三行仅是此方阵左上方三角形内各项之和.

\noindent\textbf{$\star$例 3.1.11}\quad 试证
\[
  (a^2+b^2)^{n/2}e^{ax}\sin(bx+n\varphi)
  =\sum_{i=0}^n C_n^i a^{\,n-i}b^i e^{ax}\sin\left(bx+\frac\pi2 i\right),
\]
其中 $\varphi=\arctan\dfrac ba$.

\noindent\textbf{证}\quad 令 $f(x)=e^{ax}\sin bx$，则
\[
  f'(x)=e^{ax}(a\sin bx+b\cos bx)
  =e^{ax}(a^2+b^2)^{1/2}\sin(bx+\varphi)
  \quad\left(\varphi=\arctan\frac ba\right).
\]
反复这么做，得
\begin{equation}
  f^{(n)}(x)=(a^2+b^2)^{n/2}e^{ax}\sin(bx+n\varphi).
\tag{1}
\end{equation}
另一方面，直接使用 Leibniz 公式，
\begin{equation}
  f^{(n)}(x)=\sum_{i=0}^n C_n^i a^{\,n-i}b^i e^{ax}\sin\left(bx+\frac\pi2 i\right).
\tag{2}
\end{equation}
比较式 (1) 和 (2)，即得所求的等式.

\noindent\textbf{c.\quad 用数学归纳法求高阶导数}

\noindent\textbf{要点}\quad 当高阶导数不能一次求出时，可先求出前几阶导数，总结归纳，找出规律，然后用数学归纳法加以证明.

\noindent\textbf{$\star$例 3.1.12}\quad 证明：$\left(x^{n-1}e^{1/x}\right)^{(n)}=\dfrac{(-1)^n}{x^{n+1}}e^{1/x}$.（同济大学）

\noindent\textbf{证}\quad 当 $n=1,2$ 时，可直接验证其成立. 设当 $n=k-1$ 和 $n=k$ 时欲证等式成立，下证 $n=k+1$ 时也成立. 事实上，
\[
  \left(x^k e^{1/x}\right)^{(k+1)}
  =\left[(x^k e^{1/x})'\right]^{(k)}
  =\left(kx^{k-1}e^{1/x}-x^{k-2}e^{1/x}\right)^{(k)}.
\]
再利用 $n=k$ 与 $n=k-1$ 时的已有结果，代入整理，即可证得.

\noindent\textbf{$\star$例 3.1.13}\quad 证明：若函数
\[
f(x)=
\begin{cases}
e^{-1/x^2}, & x\ne0,\\
0, & x=0,
\end{cases}
\]
则在 $x=0$ 处 $f^{(n)}(0)=0\ (n=1,2,\cdots)$.（浙江大学）

\noindent\textbf{证}\quad
\[
  f'(0)=\lim_{x\to0}\frac{e^{-1/x^2}-0}{x-0}
  =\lim_{x\to0}\frac{1/x}{e^{1/x^2}}
  \stackrel{y=1/x}{=}\lim_{y\to\infty}\frac{y}{e^{y^2}}=0.
\]
设 $f^{(n-1)}(0)=0$，易证
\[
  f^{(n-1)}(x)=p\left(\frac1x\right)e^{-1/x^2}\quad(x\ne0),
\]
其中 $p\left(\dfrac1x\right)$ 表示关于 $\dfrac1x$ 的某个多项式. 因此
\[
  f^{(n)}(0)=\lim_{x\to0}\frac{p(1/x)e^{-1/x^2}-0}{x-0}
  \stackrel{y=1/x}{=}\lim_{y\to\infty}\frac{yp(y)}{e^{y^2}}=0.
\]

\noindent\textbf{例 3.1.14}\quad 已知 $f(x)$ 有 $n$ 阶导数（$n$ 为某奇数），且
\[
  g(x)=|x-a|^n f(x).
\]
试证：$f(a)=0$ 时，$g(x)$ 在 $x=a$ 处有 $n$ 阶导数；$f(a)\ne0$ 时，$g(x)$ 在 $x=a$ 处无 $n$ 阶导数.

\noindent\textbf{提示}\quad 因 $n$ 为奇数，当 $x\ne a$ 时，
\[
  g(x)=(x-a)^n f(x)\operatorname{sgn}(x-a).
\]
利用 Leibniz 公式可直接求出 $g(x)$ 的 1 至 $n$ 阶导数. 进而按导数定义顺次可知：当 $f(a)=0$ 时，$g(x)$ 在 $x=a$ 处 1 至 $n-1$ 阶导数为零，$n$ 阶导数存在；当 $f(a)\ne0$ 时，$g^{(n)}(a)$ 不存在.

\noindent\textbf{例 3.1.15}\quad 设
\[
f(x)=
\begin{cases}
x^n\sin(\ln|x|), & x\ne0,\\
0, & x=0
\end{cases}
\quad(n\text{ 为自然数}),
\]
求证：$f(x)$ 在 $x=0$ 处有 $n-1$ 阶导数，而无 $n$ 阶导数.

\noindent\textbf{提示}\quad 易证
\[
  \frac{d^m}{dx^m}[\sin(\ln|x|)]
  =\frac1{x^m}\sum_{k=1}^m a_k\sin\left(\ln|x|+k\frac\pi2\right),
\]
其中 $a_k\ (k=1,2,\cdots,m)$ 为某些常数.

\noindent\textbf{例 3.1.16}\quad 设 $u_1(x),u_2(x),\cdots,u_k(x)$ 都是 $x$ 的 $n$ 次可微函数. 试证：
\[
  (u_1u_2\cdots u_k)^{(n)}
  =\sum_{r_1+r_2+\cdots+r_k=n}
  \frac{n!}{r_1!r_2!\cdots r_k!}u_1^{(r_1)}\cdots u_k^{(r_k)}.
\]

\noindent\textbf{证}\quad 当 $k=2$ 时，这是熟知的 Leibniz 公式. 设 $k\le m$ 时公式已成立，来证 $k=m+1$ 时成立.

记 $u=u_1u_2\cdots u_m$，则
\[
\begin{aligned}
(u_1u_2\cdots u_m u_{m+1})^{(n)}
&=(u\cdot u_{m+1})^{(n)}
=\sum_{r+r_{m+1}=n}\frac{n!}{r!\,r_{m+1}!}u^{(r)}u_{m+1}^{(r_{m+1})}\\
&=\sum_{r+r_{m+1}=n}\frac{n!}{r!\,r_{m+1}!}
\left(\sum_{r_1+\cdots+r_m=r}
\frac{r!}{r_1!\cdots r_m!}u_1^{(r_1)}\cdots u_m^{(r_m)}\right)
u_{m+1}^{(r_{m+1})}\\
&=\sum_{r_1+r_2+\cdots+r_{m+1}=n}
\frac{n!}{r_1!r_2!\cdots r_{m+1}!}
u_1^{(r_1)}u_2^{(r_2)}\cdots u_{m+1}^{(r_{m+1})}.
\end{aligned}
\]

\noindent\textbf{d.\quad 用递推公式求导}

\noindent\textbf{要点}\quad 当高阶导数无法直接求出时，可考虑先求出导数的递推公式. 方法是先求前几阶的导数关系，然后设法将等式作适当处理，使两端同时求导时能得到一般的递推关系.

\noindent\textbf{$\star$例 3.1.17}\quad 设 $f(x)=(\arcsin x)^2$，求 $f^{(n)}(0)$.

\noindent\textbf{解}\quad 由 $f'(x)=2\dfrac{\arcsin x}{\sqrt{1-x^2}}$，得
\begin{equation}
  (1-x^2)f'^2(x)=4f(x).
\tag{1}
\end{equation}
再求一次导数，整理得
\begin{equation}
  -xf'(x)+(1-x^2)f''(x)=2.
\tag{2}
\end{equation}
应用 Leibniz 公式，对式 (2) 两端同时求 $n$ 阶导数得
\begin{equation}
  -xf^{(n+1)}(x)-nf^{(n)}(x)+(1-x^2)f^{(n+2)}(x)
  -2nxf^{(n+1)}(x)-n(n-1)f^{(n)}(x)=0.
\tag{3}
\end{equation}
在式 (1)、(2)、(3) 中令 $x=0$，得
\[
  f'(0)=0,\quad f''(0)=2,\quad f^{(n+2)}(0)=n^2f^{(n)}(0),
\]
从而
\[
  f^{(2k+1)}(0)=0\quad(k=0,1,2,\cdots),
\]
\[
  f^{(2k)}(0)=(2k-2)^2(2k-4)^2\cdots2^2\cdot2
  =2^{2k-1}[(k-1)!]^2\quad(k=1,2,\cdots).
\]

\noindent\textbf{例 3.1.18}\quad 证明 Legendre 多项式
\[
  P_n(x)=\frac1{2^n n!}\left[(x^2-1)^n\right]^{(n)}\quad(n=0,1,2,\cdots)
\]
满足方程
\[
  (1-x^2)P_n''(x)-2xP_n'(x)+n(n+1)P_n(x)=0.
\]

\noindent\textbf{提示}\quad 令 $u=(x^2-1)^n$，可得 $(x^2-1)u'=2nxu$，再两端同时求 $n+1$ 阶导数.

\noindent\textbf{e.\quad 用 Taylor 展开式求导数}

\noindent\textbf{要点}\quad $f(x)$ 按 $(x-a)$ 的幂展开的幂级数必是 $f(x)$ 的 Taylor 展开式：
\[
  f(x)=\sum_{n=0}^\infty\frac{f^{(n)}(a)}{n!}(x-a)^n.
\]
因此，若一旦得到展开式 $f(x)=\sum\limits_{n=0}^\infty a_n(x-a)^n$，则
\[
  f^{(n)}(a)=a_n n!\quad(n=0,1,2,\cdots).
\]

\noindent\textbf{$\star$例 3.1.19}\quad 求 $f(x)=\arctan x$ 在 $x=0$ 处的各阶导数.

\noindent\textbf{解}\quad
\[
  f'(x)=\frac1{1+x^2}=\sum_{n=0}^\infty(-1)^n x^{2n}\quad(|x|<1).
\]
两端从 $0$ 到 $x$ 积分，可得
\[
  f(x)=\sum_{n=0}^\infty(-1)^n\frac{x^{2n+1}}{2n+1}\quad(|x|<1).
\]
由此
\[
f^{(k)}(0)=
\begin{cases}
(-1)^n(2n)!, & \text{当 }k=2n+1\text{ 时},\\
0, & \text{当 }k=2n\text{ 时}.
\end{cases}
\]

\noindent\textbf{new $\star$例 3.1.20}\quad 设函数 $f(x)$ 在 $(-1,1)$ 上无穷次可微，$f(0)=1,\ |f'(0)|\le2$. 若函数 $g(x)=\dfrac{f'(x)}{f(x)}$，有
\begin{equation}
  |g^{(n)}(0)|\le2\cdot(n!).
\tag{1}
\end{equation}
试证：对所有 $n\in\mathbb N$，皆有
\begin{equation}
  |f^{(n)}(0)|\le(n+1)!.
\tag{2}
\end{equation}
（北京大学）

\noindent\textbf{提示}\quad 对 $f'(x)=f(x)g(x)$ 两端同时求 $n$ 阶导数.（可用 Leibniz 公式和数学归纳法证明式 (2). 因 $n=1,2$ 自明，只需证明：若 $n=1$ 至 $n=m$ 结论都成立，则 $n=m+1$ 也成立.）

\noindent\textbf{再提示}\quad
\[
\begin{aligned}
|f^{(m+1)}(0)|
&=|(f'(x))^{(m)}|_{x=0}
=|(f(x)g(x))^{(m)}|_{x=0}\\
&=\left|\sum_{k=0}^m C_m^k g^{(k)}(x)f^{(m-k)}(x)\right|_{x=0}\\
&\le\sum_{k=0}^m\frac{m!}{k!(m-k)!}2\,k!\,(m-k+1)!\\
&=2\cdot m!\sum_{k=0}^m(m+1-k)
=2\cdot m!\frac{(m+1)(m+2)}2=(m+2)!.
\end{aligned}
\]

\noindent\textbf{new $\star$例 3.1.21}\quad 设
\begin{equation}
u=
\begin{vmatrix}
1&1&1&\cdots&1\\
x_1&x_2&x_3&\cdots&x_n\\
x_1^2&x_2^2&x_3^2&\cdots&x_n^2\\
\vdots&\vdots&\vdots& &\vdots\\
x_1^{n-1}&x_2^{n-1}&x_3^{n-1}&\cdots&x_n^{n-1}
\end{vmatrix},
\tag{1}
\end{equation}
求证：$\sum\limits_{j=1}^n x_j\dfrac{\partial u}{\partial x_j}=\dfrac{n(n-1)}2u$.（南开大学）

\noindent\textbf{证}\quad （按第 $j$ 列展开）$u=A_{1j}+x_jA_{2j}+\cdots+x_j^{n-1}A_{nj}$（其中代数余子式 $A_{ij}$ 里不含 $x_j$），因此
\[
  x_j\frac{\partial u}{\partial x_j}
  =x_jA_{2j}+2x_j^2A_{3j}+\cdots+(n-1)x_j^{n-1}A_{nj},\quad j=1,2,\cdots,n.
\]
从而
\begin{equation}
  \sum_{j=1}^n x_j\frac{\partial u}{\partial x_j}
  =\sum_{j=1}^n x_jA_{2j}
  +2\sum_{j=1}^n x_j^2A_{3j}
  +\cdots+(n-1)\sum_{j=1}^n x_j^{n-1}A_{nj}.
\tag{2}
\end{equation}
注意，将式 (1) 按第 $i$ 行展开 $(i=2,3,\cdots,n)$，则有
\[
  \sum_{j=1}^n x_jA_{2j}
  =\sum_{j=1}^n x_j^2A_{3j}
  =\cdots=\sum_{j=1}^n x_j^{n-1}A_{nj}=u.
\]
此式代入式 (2)，得
\[
  \sum_{j=1}^n x_j\frac{\partial u}{\partial x_j}
  =u+2u+\cdots+(n-1)u=\frac{n(n-1)}2u.
\]
等式获证.

\section*{单元练习 3.1}

\noindent\textbf{导数的计算}

\noindent\textbf{3.1.1}\quad 计算下列函数的指定导数：

1）$f(x)=\sqrt{\dfrac{(1+x)\sqrt{x}}{e^{x-1}}}+\arcsin\dfrac{1-x}{\sqrt{1+x^2}}$，求 $f'(1)$；（中国人民大学）\quad $\left(-\dfrac{\sqrt2}{2}\right)$

2）$f(x)=x^{\sin(\sin x^x)}\ (x>0)$，求 $\dfrac{dy}{dx}$；（复旦大学）
\[
  \left(\left[\cos(\sin x^x)\cdot\cos x^x\cdot x^x(\ln x+\ln^2x)
  +\frac1x\sin(\sin x^x)\right]\cdot x^{\sin(\sin x^x)}\right)
\]

3）
\[
f(x)=
\begin{cases}
\cos x, & x<0,\\
\ln(1+x^2), & x\ge0,
\end{cases}
\]
求 $f'(x)$；（华东师范大学）
\[
\left(f'(x)=
\begin{cases}
-\sin x, & x<0,\\
\dfrac{2x}{1+x^2}, & x>0,
\end{cases}
\quad f'(0)\text{ 不存在}\right)
\]

4）$f(x)=\dfrac{x}{\sqrt{1+x^2}}$，$f_n(x)=f(f(\cdots f(x)))$（$n$ 个 $f$），求 $\dfrac{df_n(x)}{dx}$；（西北工业大学）
\[
  \left(\frac1{\sqrt{(1+nx^2)^3}}\right)
\]

$\star$5）$f''(u)$ 存在，$y=f(x+y)$，求 $\dfrac{dy}{dx},\dfrac{d^2y}{dx^2}$；（中国地质大学）
\[
  \left(\frac{dy}{dx}=\frac{f'(u)}{1-f'(u)},\quad
  \frac{d^2y}{dx^2}=\frac{f''(u)}{(1-f'(u))^3}\right)
\]

6）$y=y(x)$ 为 $y=-ye^x+2e^y\sin x-7x$ 所确定的可微函数，求 $y'(0)$；（浙江大学）\quad $\left(-\dfrac52\right)$.

$\star$7）$f(x)$ 有任意阶导数，$f'(x)=(f(x))^2$，求 $f^{(n)}(x)\ (n>2)$；（数学一）\quad $\langle n!(f(x))^{n+1}\rangle$

$\star$8）$x=a(t-\sin t),\ y=a(1-\cos t)$，求 $\dfrac{dy}{dx},\dfrac{d^2y}{dx^2}$；（中国海洋大学）
\[
  \left(\frac{\sin t}{1-\cos t},\quad -\frac1{a(1-\cos t)^2}\right)
\]

$\star$9）$f^{-1}(x)$ 为 $f(y)$ 的反函数，$f'[f^{-1}(x)],f''(f^{-1}(x))$ 都存在，且 $f'(f^{-1}(x))\ne0$，证明：
\[
  \frac{d^2f^{-1}(x)}{dx^2}
  =-\frac{f''(f^{-1}(x))}{[f'(f^{-1}(x))]^3}.
\]
（湖南大学）

$\star$10）对于 $\mathbb R$ 上的实函数，若所讨论的导数存在，试证有如下结论：奇函数的导数为偶函数，偶函数的导数为奇函数. 如果将此结论简记作 $(\text{奇})'=\text{偶}$，$(\text{偶})'=\text{奇}$，则显然有：
\[
  (\text{奇})^{(2n)}=\text{奇},\quad
  (\text{奇})^{(2n-1)}=\text{偶},\quad
  (\text{偶})^{(2n)}=\text{偶},
\]
\[
  (\text{奇})'(0)=0,\quad
  (\text{奇})^{(2n)}(0)=0,\quad
  (\text{偶})^{(2n-1)}(0)=0\quad(n=1,2,\cdots).
\]

11）设
\[
  f(x)=\frac{x^5}{\sqrt{1+x^2}}\cdot\frac{\sin^4x}{1+\cos^2x},
\]
求 $f^{(6)}(0)$ 及 $\int_{-1}^1 f^{(6)}(x)\,dx$；（中国人民大学）\quad $\langle0,0\rangle$

12）求 $d^n(x^2\ln x)$（$x>0,x$ 为自变量）；（南京大学）
\[
  \left\langle d(x^2\ln x)=(2\ln x+1)x\,dx,\quad
  d^2(x^2\ln x)=(2\ln x+3)dx^2,\right.
\]
\[
  \left. d^n(x^2\ln x)=(-1)^{n-1}2(n-3)!x^{2-n}dx^n\ (n>2)\right\rangle
\]

13）$g(x)$ 在 $[-1,1]$ 上无穷次可微，且存在 $M>0$ 使得 $|g^{(n)}(x)|\le n!M$，$g\left(\dfrac1n\right)=\ln(1+2n)-\ln n\ (n=1,2,3,\cdots)$，求 $g^{(k)}(0)\ (k=0,1,2,\cdots)$；（中国科学院）
\[
  \left\langle g(0)=\ln2,\quad
  g^{(k)}(0)=(-1)^{k-1}\frac{(k-1)!}{2^k}\ (k=1,2,\cdots)\right\rangle
\]

14）$f(x)=x\sin\omega x$，证明
\[
  f^{(2n)}(x)=(-1)^n(\omega^{2n}x\sin\omega x-2n\omega^{2n-1}\cos\omega x).
\]
（北京理工大学）

$\star$15）
\[
f(x)=
\begin{cases}
\dfrac{\sin x}x, & x\ne0,\\
1, & x=0,
\end{cases}
\]
求 $f^{(k)}(0)$.（华东师范大学）
\[
  \left\langle f^{(2n)}(0)=(-1)^n\frac1{2n+1},\quad
  f^{(2n-1)}(0)=0\ (n=1,2,\cdots)\right\rangle
\]

\noindent\textbf{提示}\quad 2）$(x^{g(x)})'=(e^{g(x)\ln x})'$；

3）分段求导.

4）先写出 $f_1,f_2,f_3$，再用数学归纳法求出 $f_n$：
\[
  f_n(x)=\frac{x}{\sqrt{1+nx^2}},
\]
然后求导.

5），6）隐函数求导.

7）复合函数求导.

8）参数函数求导.

9）在 $\dfrac{dy}{dx}=\dfrac1{f'(y)}$ 两边对 $x$ 求导：
\[
  \frac{d^2f^{-1}(x)}{dx^2}=\frac{d^2y}{dx^2}
  =\frac{-f''(y)\dfrac{dy}{dx}}{(f'(y))^2}
  =-\frac{f''(f^{-1}(x))}{[f'(f^{-1}(x))]^3}.
\]

11）利用 10）的结果.

12）求出 $d^3(x^2\ln x)=\dfrac2x dx^3$ 之后再用公式直接求.

13）$g\left(\dfrac1n\right)=\ln(1+2n)-\ln n=\ln\left(2+\dfrac1n\right)$（$n=1,2,\cdots$），由此 $g(x)=\ln(2+x)$，很容易求出上述结果. 但这里隐藏着一个重大理论问题，即：$g(x)$ 与 $\ln(2+x)$ 仅在点列 $\left\{\dfrac1n\right\}$ 上相等，为什么在 $[-1,1]$ 上处处相等？

这是因为此时 $f(x)\equiv g(x)-\ln(2+x)$ 在 $[-1,1]$ 上满足
\[
  \left|\frac{f^{(n)}(x)}{n!}\right|\le M,\quad f\left(\frac1n\right)=0\quad(n=1,2,\cdots),
\]
从而可证 $f(x)\equiv0$（于 $[-1,1]$ 上）. 证明方法可参看第五章例 5.3.35.

14）用数学归纳法.

15）可沿用例 3.1.19 的方法. 注意，由 $\sin x$ 的展开式，有
\[
  f(x)=\sum_{n=1}^\infty(-1)^{n-1}\frac{x^{2n-2}}{(2n-1)!}.
\]

\noindent\textbf{导读}\quad （1）以上各题主要是考查复合函数、显式、隐式、参数式、反函数、分段函数、抽象函数等各种类型的求导问题. 至于积分、级数表示的函数如何求导，请参看第四、五章的相关内容.

（2）导数计算是微积分中最基本、最常用的运算，相对也最容易. 这类考题一般考生都会做，有一定的“送分性质”. 要想不丢分，必须做到熟练、准确、快捷. 还不太会的读者，请先复习一下教材的相关内容，此处不重复.

\noindent\textbf{导数定义及可微性质}

\noindent\textbf{3.1.2}\quad 讨论
\[
f(x)=
\begin{cases}
\dfrac1x-\dfrac1{e^x-1}, & x\ne0,\\
\dfrac12, & x=0
\end{cases}
\]
在 $x=0$ 处的连续性与可微性.（东北大学）

\noindent\textbf{提示}\quad
\[
  f'(0)=\lim_{x\to0}\frac{\dfrac1x-\dfrac1{e^x-1}-\dfrac12}{x}
  =\lim_{x\to0}\frac{[2(e^x-1)-2x-x(e^x-1)]'}{[2x^2(e^x-1)]'}
  =-\frac1{12}
\]
（可反复使用 L'Hospital 法则）.

\noindent\textbf{$\star$3.1.3}\quad 设
\[
f(x)=
\begin{cases}
x^2\sin\dfrac\pi x, & x<0,\\
A, & x=0,\\
ax^2+b, & x>0,
\end{cases}
\]
其中 $A,a,b$ 为常数，试问 $A,a,b$ 为何值时 $f(x)$ 在 $x=0$ 处可导，为什么？并求 $f'(0)$.（郑州大学）

\noindent\textbf{提示}\quad
\[
  f'_-(0)=\lim_{x\to0-}\left(x\sin\frac\pi x-\frac Ax\right)
\]
欲存在，必须 $A=0$；
\[
  f'_+(0)=\lim_{x\to0+}\left(ax+\frac bx\right)
\]
欲存在，必须 $b=0$. 此时 $f'(0)=0$.

\noindent\textbf{$\star$3.1.4}\quad 设 $f(x)$ 在 $x=0$ 处可导，$f(0)\ne0,\ f'(0)\ne0$，
\[
  af(h)+bf(2h)-f(0)=o(h)\quad(\text{当 }h\to0\text{ 时}),
\]
求 $a,b$.（数学一）

\noindent\textbf{提示}\quad
\[
o(1)=\frac{af(h)+bf(2h)-f(0)}h
=a\frac{f(h)-f(0)}h+2b\frac{f(2h)-f(0)}{2h}
\frac{(a+b-1)f(0)}h.
\]
令 $h\to0$，取极限得 $a=2,\ b=-1$.

\noindent\textbf{$\star$3.1.5}\quad 设函数 $f(x)$ 在闭区间 $[0,1]$ 上四次连续可微，$f(0)=0,\ f'(0)=0$. 证明：函数
\[
F(x)=
\begin{cases}
\dfrac{f(x)}{x^2}, & \text{当 }0<x\le1\text{ 时},\\
\dfrac{f''(0)}2, & \text{当 }x=0\text{ 时}
\end{cases}
\]
在 $[0,1]$ 上二次连续可微.（吉林大学）

\noindent\textbf{提示}\quad 利用导数定义，反复使用 L'Hospital 法则. 不难求得
\[
F'(x)=
\begin{cases}
\dfrac{xf'(x)-2f(x)}{x^3}, & 0<x\le1,\\
\dfrac{f'''(0)}6, & x=0,
\end{cases}
\]
\[
F''(x)=
\begin{cases}
\dfrac{x^2f''(x)-4xf'(x)+6f(x)}{x^4}, & 0<x\le1,\\
\dfrac{f^{(4)}(0)}{12}, & x=0,
\end{cases}
\]
且能验证 $\lim\limits_{x\to0+}F''(x)=F''(0)$. 从而 $F(x)$ 在 $[0,1]$ 上二次连续可微.

\noindent\textbf{注}\quad 要保证 $F(x)$ 在原点右连续，条件 $f'(0)=0$ 是必要的：
\[
  \lim_{x\to0+}F(x)=\lim_{x\to0+}\frac{f(x)}{x^2}
  =\lim_{x\to0+}\frac{f(0)+f'(0)x+\frac12f''(0)x^2+o(x^3)}{x^2}=F(0).
\]

\noindent\textbf{3.1.6}\quad 设函数 $f(x)$ 在区间 $[a,b]$ 上满足
\[
  |f(x)-f(y)|\le M|x-y|^\alpha,\quad \forall x,y\in[a,b],
\]
其中 $M>0,\alpha>1$ 为常数，证明：$f(x)$ 在 $[a,b]$ 上恒为常数.

\noindent\textbf{提示}\quad
\[
  0<\left|\frac{f(x)-f(y)}{x-y}\right|
  \le M|x-y|^{\alpha-1}\to0\quad(x\to y)\Rightarrow f'(x)\equiv0\Rightarrow f(x)=\text{常数}.
\]

\noindent\textbf{$\star$3.1.7}\quad 设 $f(0)=0$，则 $f(x)$ 在点 $x=0$ 可导的充要条件为 $\lim\limits_{h\to0}\dfrac1h f(1-e^h)$ 存在.（数学一）

\noindent\textbf{提示}\quad
\[
  \lim_{h\to0}\frac{f(1-e^h)-f(0)}h
  \stackrel{1-e^h=x}{=}\lim_{x\to0}\frac{f(x)-f(0)}x\cdot\frac{x}{\ln(1-x)}.
\]

\noindent\textbf{注}\quad 不妨尝试对该题作一小小的推广.

设 $x=g(h)$ 为：具有反函数 $g^{-1}$，且使得 $\lim\limits_{x\to0}\dfrac{x}{g^{-1}(x)}$ 存在的某一函数（例如 $x=1-e^h$），那么对任一函数 $f(x)$，若 $f(0)=0$，则 $f$ 在 $x=0$ 处可导的充要条件是：$\lim\limits_{h\to0}\dfrac1h f(g(h))$ 存在.

\noindent\textbf{$\star$3.1.8}\quad 设 $f(x)$ 在 $x_0$ 的某邻域内有定义.

1）若 $f(x)$ 在 $x_0$ 处可导，试证：
\[
  \lim_{h\to0}\frac{f(x_0+h)-f(x_0-h)}{2h}=f'(x_0);
\]

2）反之，若上式左端之极限存在，是否能推出 $f'(x_0)$ 存在？若结论成立，请证明；不成立请给出反例.（哈尔滨工业大学）

\noindent\textbf{提示}\quad 1）左端极限
\[
=\frac12\lim_{h\to0}\frac{f(x_0+h)-f(x_0)}h
+\frac12\lim_{h\to0}\frac{f(x_0-h)-f(x_0)}{-h}.
\]

2）不能！任何以 $x=x_0$ 为对称轴的函数左端极限均存在且极限为 0，但例如函数 $y=|x-x_0|$，在 $x_0$ 处就不可导.

\noindent\textbf{$\star$3.1.9}\quad 在什么条件下，函数
\[
f(x)=
\begin{cases}
x^n\sin\dfrac1x, & x\ne0,\\
0, & x=0
\end{cases}
\quad(n\text{ 为自然数})
\]
1）在点 $x=0$ 处连续；\quad $\langle n>0\rangle$

2）在点 $x=0$ 处可导；\quad $\langle n>1\rangle$

3）在点 $x=0$ 处导函数连续.（中国科学院）\quad $\langle n>2\rangle$

\noindent\textbf{提示}\quad 1）利用 $0\le\left|x^n\sin\dfrac1x\right|\le|x|^n\to0\ (n>0)$（$x\to0$）.

2）
\[
  \frac{f(x)-f(0)}x=x^{n-1}\sin\frac1x\stackrel{n>1\text{时}}{\longrightarrow}0=f'(0)\quad(x\to0).
\]

3）当 $x\ne0$ 时，
\[
  f'(x)=nx^{n-1}\sin\frac1x-x^{n-2}\cos\frac1x
  \stackrel{n>2\text{时}}{\longrightarrow}0=f'(0)\quad(x\to0).
\]

\noindent\textbf{3.1.10}\quad 试作一函数 $f(x)$：在 $(-\infty,+\infty)$ 内二阶可微，且 $f''(x)$ 在 $x=0$ 处不连续，其余处处连续.（山东大学）
\[
  \left(f(x)=
\begin{cases}
x^4\sin\dfrac1x, & x\ne0,\\
0, & x=0
\end{cases}\right)
\]

\noindent\textbf{$\star$3.1.11}\quad 对于函数 $f(x)=|\sin x|^3,\ x\in(-1,1)$.

（1）证明：$f''(0)$ 不存在；

（2）说明点 $x=0$ 是不是 $f''(x)$ 的可去间断点.（武汉大学）

\noindent\textbf{提示}\quad
\[
  f(x)=(\operatorname{sgn}x)\sin^3x,\quad x\in(-1,1),
\]
\[
  f'(x)=(\operatorname{sgn}x)3\sin^2x\cos x,\quad x\in(-1,1)\ (x=0\text{ 处验证后成立}),
\]
\[
  f''(x)=(\operatorname{sgn}x)(6\sin x\cos^2x-3\sin^3x),\quad x\in(-1,1),\ x\ne0,
\]
$f''_+(0)=6,\ f''_-(0)=-6\Rightarrow f''(0)$ 不存在，且 $x=0$ 不是可去间断点.

\noindent\textbf{3.1.12}\quad 设函数 $f(x)$ 在点 $a$ 处连续，且 $|f(x)|$ 在 $a$ 处也可导，证明：$f(x)$ 在 $a$ 处也可导.（中南大学）

\noindent\textbf{证}\quad 若 $g(x)=|f(x)|$ 在 $a$ 处可导.

$1^\circ$ 当 $f(a)>0$ 时，因 $f(x)$ 连续，根据保号性，在 $a$ 点的某邻域内 $f(x)$ 恒为正，因此 $f(x)=|f(x)|=g(x)$，故 $f(x)$ 在点 $a$ 也可导，且 $f'(a)=g'(a)$.

$2^\circ$ 当 $f(a)<0$ 时，同理可证.

$3^\circ$ 当 $f(a)=0$ 时，
\[
  g'_+(a)=\lim_{x\to a+}\frac{|f(x)|}{x-a}\stackrel{\text{记}}{=}A\ge0,\quad
  g'_-(a)=\lim_{x\to a-}\frac{|f(x)|}{x-a}\stackrel{\text{记}}{=}B\le0.
\]
因 $g$ 可导，故 $A=B=0$. 于是
\[
  \lim_{x\to a}\left|\frac{f(x)}{x-a}\right|=0,
\]
即 $f'(a)=\lim\limits_{x\to a}\dfrac{f(x)}{x-a}=0$，因此 $f(x)$ 在 $a$ 处也可导.

\noindent\textbf{3.1.13}\quad 函数 $f(x)=(x^2-x-2)|x^3-x|$ 的不可导点的个数是多少？（数学一）\quad $\langle2\rangle$

\noindent\textbf{提示}\quad $f(x)=(x-2)(x+1)|x+1|\,|x|\,|x-1|$，注意函数 $u|u|$ 在 $u=0$ 处可导，$|u|$ 在 $u=0$ 处不可导，因此 $f(x)$ 的不可导点为 $0,1$.

\noindent\textbf{3.1.14}\quad 设 $\varphi(x)$ 在 $x=a$ 处连续，分别讨论下面函数在 $x=a$ 处是否可导：

1）$f(x)=(x-a)\varphi(x)$；

2）$f(x)=|x-a|\varphi(x)$；

3）$f(x)=(x-a)|\varphi(x)|$.（武汉大学）

\noindent\textbf{提示}\quad 直接由定义 $f'(a)=\lim\limits_{x\to a}\dfrac{f(x)-f(a)}{x-a}$ 可知，

1）$f'(a)$ 存在 $=\varphi(a)$.

2）当且仅当 $\varphi(a)=0$ 时，$f$ 在 $a$ 点可导，此时 $f'(a)=0$.

3）$f'(a)$ 存在 $=|\varphi(a)|$.

\noindent\textbf{$\star$3.1.15}\quad 设 $f(x)$ 可导，$F(x)=f(x)(1+|\sin x|)$，证明：$f(0)=0$ 是 $F(x)$ 在 $x=0$ 处可导的充要条件.（数学一）

\noindent\textbf{提示}\quad
\[
  F'_+(0)=f'_+(0)+\lim_{x\to0+}\frac{f(x)|\sin x|}{x}=f'(0)+f(0),
\]
\[
  F'_-(0)=f'_-(0)+\lim_{x\to0-}\frac{f(x)|\sin x|}{x}=f'(0)-f(0).
\]

\noindent\textbf{3.1.16}\quad 求 $f(x)=[x]\sin\pi x$ 的单侧导数，并讨论可微性（$[x]$ 表示不超过 $x$ 的最大整数）.
\[
  \left\langle
  \begin{array}{l}
  \text{非整数上可导，}f'(x)=\pi[x]\cos\pi x;\\
  f'_+(k)=(-1)^kk\pi,\quad f'_-(k)=(-1)^k(k-1)\pi,\\
  \text{故整数点处不可导}
  \end{array}
  \right\rangle
\]

\noindent\textbf{提示}\quad 非整数点处，$[x]$ 在该点的某邻域为常数，于是 $f'(x)=\pi[x]\cos\pi x$.

整数 $k$ 处右导数
\[
  f'_+(k)=\lim_{h\to0+}\frac{[k+h]\sin\pi(k+h)-[k]\sin\pi k}{h}
  =k(-1)^k\pi.
\]
左导数
\[
  f'_-(k)=\lim_{h\to0-}\frac{[k+h]\sin\pi(k+h)-[k]\sin\pi k}{h}
  =(k-1)(-1)^k\pi.
\]

\noindent\textbf{3.1.17}\quad 证明：
\[
  f(x)=
  \begin{cases}
  x^2\left|\cos\dfrac{\pi}{x}\right|, & x\ne0,\\
  0, & x=0
  \end{cases}
\]
在 $x=0$ 的任何邻域内都有不可微点，但在 $x=0$ 处可微.

\noindent\textbf{提示}\quad $x_n=\left(2n+\dfrac12\right)^{-1}$ 处，$f'_+(x_n)\ne f'_-(x_n)$，不可微.

\noindent\textbf{再提示}\quad
\[
  f'_+(x_n)=\lim_{h\to0+}\frac{(x_n+h)^2\left|\cos\dfrac{\pi}{x_n+h}\right|}{h}=-\pi,
\]
同理 $f'_-(x_n)=\pi$. 而
\[
  f'(0)=\lim_{x\to0}\frac{x^2\left|\cos\dfrac{\pi}{x}\right|}{x}=0.
\]

\noindent\textbf{3.1.18}\quad 证明：切比雪夫多项式
\[
  T_m(x)=\frac{1}{2^{m-1}}\cos(m\arccos x),\quad m=0,1,2,\cdots
\]
满足
\[
  (1-x^2)T_m''(x)-xT_m'(x)+m^2T_m(x)=0.
\]

\noindent\textbf{提示}\quad 可以应用复合函数微分法直接验证.

\noindent\textbf{3.1.19}\quad 证明：切比雪夫-Laguerre 多项式
\[
  L_m(x)=e^x(x^me^{-x})^{(m)},\quad m=0,1,2,\cdots
\]
满足方程
\begin{equation}
  xL_m''(x)+(1-x)L_m'(x)+mL_m(x)=0.
\tag{1}
\end{equation}

\noindent\textbf{提示}\quad 可利用乘积求导的 Leibniz 公式：
\[
  L_m(x)=(-1)^m\sum_{k=0}^{m}(-1)^k C_m^k(x^m)^{(k)},
\]
代入方程（1），验证系数为零：如 $x^m$ 项系数 $m-m=0$（不计公因子 $(-1)^m$），$x^{m-1}$ 项系数
\[
  -m^3+m+m^3-m^2+m^2-m=0,\ \cdots
\]

\noindent\textbf{※3.1.20}\quad 设
\[
  f(x)=
  \begin{vmatrix}
  u_{11}(x)&u_{12}(x)&\cdots&u_{1k}(x)\\
  u_{21}(x)&u_{22}(x)&\cdots&u_{2k}(x)\\
  \vdots&\vdots&&\vdots\\
  u_{k1}(x)&u_{k2}(x)&\cdots&u_{kk}(x)
  \end{vmatrix}
\]
（$u_{ij}(x)$ 为 $n$ 次可微函数），试证
\[
  f^{(n)}(x)=\sum_{r_1+\cdots+r_k=n}\frac{n!}{r_1!\cdots r_k!}
  \begin{vmatrix}
  u_{11}^{(r_1)}&u_{12}^{(r_1)}&\cdots&u_{1k}^{(r_1)}\\
  u_{21}^{(r_2)}&u_{22}^{(r_2)}&\cdots&u_{2k}^{(r_2)}\\
  \vdots&\vdots&&\vdots\\
  u_{k1}^{(r_k)}&u_{k2}^{(r_k)}&\cdots&u_{kk}^{(r_k)}
  \end{vmatrix}.
\]

\noindent\textbf{提示}\quad 方法 I：可将行列式展开成
\[
  f(x)=\sum(-1)^\lambda u_{1,i_1}u_{2,i_2}\cdots u_{k,i_k},
\]
其中 $\lambda$ 为排列 $i_1,i_2,\cdots,i_k$ 的逆序数，然后逐项求导，利用例 3.1.16 即可.

方法 II：可利用数学归纳法.

\noindent\textbf{*3.1.21}\quad 设
\[
  x=a\cos t+b\sin t,\quad y=a\sin t-b\cos t,
\]
求证：
\[
  \frac{d^m x}{dt^m}\frac{d^n y}{dt^n}
  -\frac{d^n x}{dt^n}\frac{d^m y}{dt^m}
  =(a^2+b^2)\sin\frac{n-m}{2}\pi.
\]

\noindent\textbf{提示}\quad 令
\[
  x=\sqrt{a^2+b^2}\left(\frac{a}{\sqrt{a^2+b^2}}\cos t+
  \frac{b}{\sqrt{a^2+b^2}}\sin t\right)
  =\sqrt{a^2+b^2}\cos(t-\alpha).
\]
同理
\[
  y=\sqrt{a^2+b^2}\sin(t-\alpha),
\]
这时 $\tan\alpha=b/a$.

\noindent\textbf{留念题}\quad

\noindent\textbf{※3.1.22}\quad 对例 3.1.7 如下的证法给出评论，认为正确请说明理由，认为不正确请给出反例.

\noindent\textbf{证}\quad
\[
  \lim_{x\to0}\frac{f(2x)-f(x)}{x}=A
\]
于是
\[
  f(2x)-f(x)=Ax+o(x).
\]
从而
\[
  f\left(\frac{x}{2^{k-1}}\right)-f\left(\frac{x}{2^k}\right)
  =A\frac{x}{2^k}+o\left(\frac{x}{2^k}\right).
\]
故
\[
  \sum_{k=1}^n\left[
  f\left(\frac{x}{2^{k-1}}\right)-f\left(\frac{x}{2^k}\right)\right]
  =A\sum_{k=1}^n\frac{x}{2^k}+o\left(\sum_{k=1}^n\frac{x}{2^k}\right),
\]
即
\[
  f(x)-f\left(\frac{x}{2^n}\right)
  =Ax(1-2^{-n})+o((1-2^{-n})x).
\]
令 $n\to\infty$，得
\[
  f(x)-f(0)=Ax+o(x),
\]
因此
\[
  f'(0)=\lim_{x\to0}\frac{f(x)-f(0)}{x}=A.
\]

\section*{§3.2\quad 微分中值定理}

本节主要讨论微分中值定理的应用，要指出的是，本节各部分都广泛使用辅助函数法，应该说它是应用微分中值的基本方法. 实际上，辅助函数法是转化问题的一种重要手段，它在数学分析的其他地方也时常用到. 我们不准备在下面每一部分里列入辅助函数法的条目，关于如何选用和构造辅助函数的问题，将在例题里穿插讲解.

\subsection*{一、Rolle 定理}

\noindent\textbf{a. 函数零（值）点问题}

\noindent\textbf{要点}\quad 零点存在性问题.

（1）借助介值性求解（连续函数有介值性，导函数也有介值性（见例 3.2.24））.

（2）借助 Rolle 定理求解：即若 $f(x)$ 在 $[a,b]$ 上连续，在 $(a,b)$ 内可导且 $f(a)=f(b)$，则 $\exists\xi\in(a,b)$，使得 $f'(\xi)=0$.

此外，Fermat 定理告诉我们：“函数在其极值点处如有导数存在，此点必是导数的零点.”

\noindent\textbf{零点的唯一性问题}\quad 这里主要通过单调性来确定零点的唯一性，若 $f'>0$，则 $f$ 严增（或 $f'<0$，则 $f$ 严减），这时 $f$ 的零点若存在则必唯一.

\noindent\textbf{$\star$例 3.2.1}\quad 设 $(a,b)$ 为有限或无穷区间，$f(x)$ 在 $(a,b)$ 内可微，且
\[
  \lim_{x\to a+}f(x)=\lim_{x\to b-}f(x)=A
\]
（有限或 $\pm\infty$），试证：$\exists\xi\in(a,b)$，使得 $f'(\xi)=0$.（北京师范大学）

\noindent\textbf{证}\quad 若 $f(x)\equiv A$（有限数），则 $f'(x)\equiv0$，问题自明.

若 $f(x)\not\equiv A$，则 $\exists x_0\in(a,b)$，使得 $f(x_0)\ne A$，下设 $f(x_0)>A$（对 $f(x_0)<A$ 类似可证）. 因为
\[
  \lim_{x\to a+}f(x)=\lim_{x\to b-}f(x)=A,
\]
函数 $f(x)$ 在 $(a,b)$ 内连续，所以对于任意取定的数 $\mu\ (A<\mu<f(x_0))$，$\exists x_1\in(a,x_0)$，$x_2\in(x_0,b)$，使得 $f(x_1)=f(x_2)=\mu$. 从而由 Rolle 定理知，$\exists\xi\in(x_1,x_2)\subset(a,b)$，使得 $f'(\xi)=0$. 若 $A=+\infty$（或 $-\infty$），则 $(a,b)$ 内任取一点作 $x_0$，上面推理保持有效.

\noindent\textbf{new 练习 1}\quad 设 $f(x)$ 在 $[0,1]$ 上可微，$f(0)=0$，$\forall x\in(0,1)$，$f(x)\ne0$. 试证：$\exists\xi\in(0,1)$ 使得
\[
  2\frac{f'(\xi)}{f(\xi)}=\frac{f'(1-\xi)}{f(1-\xi)}.
\]
（四川大学，南开大学）

\noindent\textbf{提示}\quad
\[
  2\frac{f'(\xi)}{f(\xi)}=\frac{f'(1-\xi)}{f(1-\xi)}
  \Longleftrightarrow
  2f'(\xi)f(1-\xi)-f(\xi)f'(1-\xi)=0
\]
（同乘 $f(\xi)$）
\[
  \Longleftrightarrow
  \left.(f^2(x)f(1-x))'\right|_{x=\xi}=0.
\]
作辅助函数 $F(x)=f^2(x)f(1-x)$，应用 Rolle 定理，$F(x)$ 可在 $[0,1]$ 上找到满足条件的 $\xi$.

\noindent\textbf{new 练习 2}\quad 设函数 $f(x)$ 在闭区间 $[-1,1]$ 上具有三阶连续导数，且 $f(-1)=0,\ f(1)=1,\ f'(0)=0$，求证：$\exists\xi\in(-1,1)$，使得 $f'''(\xi)=3$.（华中师范大学）

\noindent\textbf{提示}\quad （待定系数法）若令
\begin{equation}
  F(x)=f(x)-\frac12x^3+ax^2+bx+c,
\tag{1}
\end{equation}
则 $f'''(\xi)=3$ 等价于 $F'''(\xi)=0$.（转化为求 $F'''$ 的零点问题）适当挑选待定系数 $a,b,c$，使 $F(x)$ 有足够多的零点，就可确保 $F'''$ 至少存在一个零点.

\noindent\textbf{再提示}\quad 1）在式（1）里令 $x=0$，知当 $c=-f(0)$ 时，$F(0)=0$.

2）因 $f(-1)=0,\ f(1)=1$，式（1）中令 $a=f(0)-\dfrac12,\ b=0$，可使 $F(-1)=F(1)=0$. 至此 $F$ 共有三个零点：$-1,0,1$. 利用 Rolle 定理，可找出 $F'$ 的两个（互异）零点，分别在 $(-1,0)$ 和 $(0,1)$ 内.

3）对式（1）求导，题设有 $f'(0)=0$，且令 $b=0$，故 $F'(0)=0$，于是 $F'$ 有三个（互异）零点.

4）利用 Rolle 定理，可找出 $F''$ 的两个（互异）零点，在这两个零点之间，$\exists\xi\in(-1,1)$，使得 $F'''(\xi)=0$，亦即 $f'''(\xi)=3$.

\noindent\textbf{例 3.2.2}\quad 若 $P_n(x)=a_0x^n+a_1x^{n-1}+\cdots+a_{n-1}x+a_n\ (a_0\ne0)$ 为实系数多项式，且 $P_n(x)=0$ 的一切根皆为实数，试证：$P_n'(x)=0,\ P_n''(x)=0,\cdots,P_n^{(n-1)}(x)=0$ 也仅有实根.

\noindent\textbf{证}\quad 设
\begin{equation}
  P_n(x)=a_0(x-\alpha_1)^{k_1}\cdots(x-\alpha_m)^{k_m},
\tag{1}
\end{equation}
其中 $\alpha_1<\alpha_2<\cdots<\alpha_m$ 分别为 $P_n(x)=0$ 的 $k_1,k_2,\cdots,k_m$ 重根，$k_1+k_2+\cdots+k_m=n$. 由 Rolle 定理，在相邻两异根之间，存在 $P_n'(x)=0$ 的一个根. 因此，$P_n'(x)=0$ 在 $P_n(x)=0$ 的 $m$ 个根的间隔里共有 $m-1$ 个根. 又从式（1）可知，当 $\alpha_i$ 是 $k_i>1$ 重根时，则 $\alpha_i$ 必是 $P_n'(x)=0$ 的 $k_i-1$ 重根. 因此，$P_n'(x)=0$ 共有
\[
  (m-1)+(k_1-1)+(k_2-1)+\cdots+(k_m-1)
  =k_1+k_2+\cdots+k_m-1=n-1
\]
个根. 但 $P_n'(x)$ 是 $n-1$ 次多项式，故 $P_n'(x)=0$ 也仅有 $n-1$ 个根. 所以 $P_n'(x)=0$ 的根全为实的. 反复这样做 $n-1$ 次，知 $P_n'(x)=0,\cdots,P_n^{(n-1)}(x)=0$ 的根都是实的.

\noindent\textbf{例 3.2.3}\quad 证明：Legendre 多项式
\[
  P_n(x)=\frac{1}{2^nn!}\frac{d^n}{dx^n}\left[(x^2-1)^n\right]
\]
的一切零点在 $(-1,1)$ 内.

\noindent\textbf{证}\quad $f(x)=(x^2-1)^n=(x+1)^n(x-1)^n$ 为 $2n$ 次多项式，$x=\pm1$ 分别为 $f(x)=0$ 的 $n$ 重根. 跟上例同样的道理知，$f'(x)=0$ 以 $x=\pm1$ 为 $n-1$ 重根，且 $\exists\xi_1^{(1)}\in(-1,1)$ 为单根；进而 $f''(x)=0$ 以 $x=\pm1$ 为 $n-2$ 重根，且 $\exists\xi_1^{(2)},\xi_2^{(2)}\in(-1,1)$ 为单根；反复 $n$ 次，可知 $f^{(n)}(x)=0$ 不再以 $x=\pm1$ 为根，在 $(-1,1)$ 有 $n$ 个互异实根，但
\[
  P_n(x)=\frac{1}{2^nn!}\frac{d^n}{dx^n}\left[(x^2-1)^n\right]
  =\frac{1}{2^nn!}f^{(n)}(x)
\]
为 $n$ 次多项式，只有 $n$ 个零点，所以 $P_n(x)$ 的一切根在 $(-1,1)$ 内.

\noindent\textbf{例 3.2.4}\quad 设 $f(x)$ 在 $[a,b]$ 上连续，在 $(a,b)$ 内可导，$f(a)=f(b)=0$，试证：
\[
  \forall\alpha\in\mathbb{R},\ \exists\xi\in(a,b)\text{ 使得 }\alpha f(\xi)=f'(\xi).
\]
（广西师范大学）

\noindent\textbf{分析}\quad 要 $\alpha f(\xi)=f'(\xi)$，即要 $f'(\xi)-\alpha f(\xi)=0$，亦即要 $\xi$ 为函数 $f'(x)-\alpha f(x)$ 的零点. 注意到
\[
  [f(x)e^{-\alpha x}]'=[f'(x)-\alpha f(x)]e^{-\alpha x}.
\]
因此，只要对函数 $F(x)=f(x)e^{-\alpha x}$ 检验 Rolle 定理条件. 但这是明显满足的. 此问题可以推广到一般情况.

\noindent\textbf{*例 3.2.5}\quad 设 $f(x)$ 在 $[a,b]$ 上有 $n$ 阶连续导数，且 $f(x)=0$ 在 $[a,b]$ 上至少有 $n+1$ 个不同实根，
\[
  P_n(x)=x^n+c_{n-1}x^{n-1}+c_{n-2}x^{n-2}+\cdots+c_1x+c_0
\]
是仅有实根的实系数多项式. 引入记号 $D^n=\dfrac{d^n}{dx^n}$，试证：
\[
  P_n(D)f(x)\equiv(D^n+c_{n-1}D^{n-1}+c_{n-2}D^{n-2}+\cdots+c_1D+c_0)f(x)
\]
在 $[a,b]$ 上至少有一个零点.

\noindent\textbf{分析}\quad 设多项式 $P_n(x)=0$ 的 $n$ 个不同实根为 $\alpha_1,\alpha_2,\cdots,\alpha_n$，则由代数知识，$P_n(x)$ 可分解为
\[
  P_n(x)=(x-\alpha_1)(x-\alpha_2)\cdots(x-\alpha_n).
\]
将 $(D-\alpha_1)(D-\alpha_2)\cdots(D-\alpha_n)$ 展开即知
\[
  P_n(D)f(x)=(D-\alpha_1)(D-\alpha_2)\cdots(D-\alpha_n)f(x).
\]
利用上例的结果可知：在 $f(x)=0$ 的两个不同实根之间必有
\[
  (D-\alpha_n)f(x)=f'(x)-\alpha_nf(x)=0
\]
的一个实根. 但已知 $f(x)=0$ 有 $n+1$ 个不同实根，故 $(D-\alpha_n)f(x)=0$ 至少有 $n$ 个不同实根. 进而依次可知 $(D-\alpha_{n-1})(D-\alpha_n)f(x)=0$ 至少有 $n-1$ 个不同实根，$\cdots$，
\[
  P_n(D)f(x)=(D-\alpha_1)(D-\alpha_2)\cdots(D-\alpha_n)f(x)=0
\]
至少有 1 个实根.

试用数学归纳法写出简洁的证明.

\noindent\textbf{b. 证明中值公式}

\noindent\textbf{要点}\quad 构造不同的辅助函数，应用 Rolle 定理，可以导出不同的中值公式.

\noindent\textbf{例 3.2.6}\quad 设 $f(x),g(x),h(x)$ 在 $[a,b]$ 上连续，在 $(a,b)$ 内可导，试证存在 $\xi\in(a,b)$，使得
\begin{equation}
  \begin{vmatrix}
  f(a)&g(a)&h(a)\\
  f(b)&g(b)&h(b)\\
  f'(\xi)&g'(\xi)&h'(\xi)
  \end{vmatrix}=0.
\tag{1}
\end{equation}

\noindent\textbf{证}\quad 记
\[
  F(x)=
  \begin{vmatrix}
  f(a)&g(a)&h(a)\\
  f(b)&g(b)&h(b)\\
  f(x)&g(x)&h(x)
  \end{vmatrix},
\]
则 $F(x)$ 在 $[a,b]$ 上连续，在 $(a,b)$ 内可导，$F(a)=F(b)=0$. 应用 Rolle 定理可知，$\exists\xi\in(a,b)$，使得 $F'(\xi)=0$. 根据行列式性质，$F'(\xi)=0$，即是式（1）.

\noindent\textbf{注}\quad 由本例可以推出 Cauchy 中值定理（只要令 $h(x)\equiv1$）和 Lagrange 定理（令 $h(x)\equiv1,\ g(x)\equiv x$）.

\noindent\textbf{例 3.2.7}\quad 设 $f(x),g(x)$ 在 $[a,b]$ 上连续，在 $(a,b)$ 内可导，$\forall x\in(a,b)$，$g'(x)\ne0$，试证：$\exists\xi\in(a,b)$，使得
\[
  \frac{f'(\xi)}{g'(\xi)}=\frac{f(\xi)-f(a)}{g(b)-g(\xi)}.
\]

\noindent\textbf{提示}\quad 将目标等式变形，找出辅助函数 $F(x)=[f(x)-f(a)][g(b)-g(x)]$.

\noindent\textbf{例 3.2.8}\quad 设 $f(x)$ 在 $[0,+\infty)$ 上可导，且 $0\le f(x)\le\dfrac{x}{1+x^2}$. 证明：$\exists\xi>0$，使得
\[
  f'(\xi)=\frac{1-\xi^2}{(1+\xi^2)^2}.
\]

\noindent\textbf{证}\quad 问题相当于要找出 $\xi>0$，使得
\begin{equation}
  \left.\left(f(x)-\frac{x}{1+x^2}\right)'\right|_{\xi}=0.
\tag{1}
\end{equation}
因函数 $F(x)=f(x)-\dfrac{x}{1+x^2}$ 在 $[0,+\infty)$ 上连续，在 $(0,+\infty)$ 内可导，$F(0)=F(+\infty)=0$，所以由推广了的 Rolle 定理（见例 3.2.1），知 $\exists\xi>0$，使得 $F'(\xi)=0$，此即式 (1) 成立.

\noindent\textbf{*例 3.2.9}\quad 设 $f(x)$ 在包含 $x_0$ 的区间 $I$ 上二次可微，$x_0+h\in I,\lambda\in(0,1)$，试证：$\exists\theta\in(0,1)$，使得
\begin{equation}
  f(x_0+\lambda h)=\lambda f(x_0+h)+(1-\lambda)f(x_0)+\frac{\lambda}{2}(\lambda-1)h^2f''(x_0+\theta h).
\tag{1}
\end{equation}

\noindent\textbf{分析}\quad 因 $0<\lambda<1$，可取数 $M$，使得
\begin{equation}
  f(x_0+\lambda h)-\lambda f(x_0+h)-(1-\lambda)f(x_0)-\frac{\lambda}{2}(\lambda-1)h^2\cdot M=0.
\tag{2}
\end{equation}
故只要证明：$\exists\theta\in(0,1)$，使得 $M=f''(x_0+\theta h)$. 取 $\lambda$ 作为变量，记为 $t$. 令
\begin{equation}
  F(t)=f(x_0+th)-tf(x_0+h)-(1-t)f(x_0)-\frac{t}{2}(t-1)h^2M,
\tag{3}
\end{equation}
则 $F(t)$ 在 $[0,1]$ 上二次可微，且有三个零点：
\[
  F(0)=F(1)=F(\lambda)=0.
\]
两次应用 Rolle 定理，可知 $\exists\theta\in(0,1)$ 使得 $F''(\theta)=0$.

依据式 (3)，此即 $M=f''(x_0+\theta h)$. 代回式 (2)，移项，即得欲证的式 (1).

\subsection*{二、Lagrange 定理}

\noindent\textbf{a. 利用几何意义（弦线法）}

\noindent\textbf{要点}\quad 由 Lagrange 定理知，若 $f(x)$ 在 $[a,b]$ 上连续，在 $(a,b)$ 内可导，则 $\forall x_1,x_2\in[a,b]$，$\exists\xi\in(x_1,x_2)$，使得
\[
  \frac{f(x_2)-f(x_1)}{x_2-x_1}=f'(\xi).
\]
即是说：曲线上任意两点的弦，必与两点间某点的切线平行. 我们可以用这种几何解释进行思考和解题.

\noindent\textbf{$\star$例 3.2.10}\quad 设 $f(x)$ 是可微函数，导函数 $f'(x)$ 严格单调递增，若 $f(a)=f(b)\ (a<b)$，试证：对一切 $x\in(a,b)$，有 $f(x)<f(a)=f(b)$（不得直接利用凸函数的性质）.（华东师范大学）

\noindent\textbf{分析}\quad 任意取一点 $x\in(a,b)$，要证明 $f(x)<f(a)=f(b)$. 如图 3.2.1 作弦 $AC,BC$. 应用 Lagrange 定理，$\exists\xi\in(a,x),\eta\in(x,b)$，使得导数 $f'(\xi),f'(\eta)$ 分别等于弦 $AC,BC$ 的斜率. 但因 $f'$ 严格递增，所以 $f'(\xi)<f'(\eta)$，这就得到（弦 $AC$ 的斜率）$<$（弦 $BC$ 的斜率）：
\[
  \frac{f(x)-f(a)}{x-a}<\frac{f(b)-f(x)}{b-x}.
\]
这便得到关于函数值的不等式. 注意到 $f(a)=f(b)$，移项即得 $f(x)<f(a)=f(b)$.

\begin{center}
  \includegraphics[width=0.38\textwidth]{3-2-1.png}

  图 3.2.1
\end{center}

\noindent\textbf{*例 3.2.11}\quad 若函数 $f(x)$ 在 $(a,b)$ 内可微，且 $\lim\limits_{x\to a+}f(x)=\infty$，试证 $\overline{\lim\limits_{x\to a+}}f'(x)=\infty$.

\noindent\textbf{分析}\quad 要证 $\overline{\lim\limits_{x\to a+}}f'(x)=\infty$，只要证明存在序列 $\{\xi_n\}\to a+$，使得 $\lim\limits_{n\to\infty}|f'(\xi_n)|=+\infty$. 利用 Lagrange 定理，只要找到一串弦，从右侧趋向 $a$，使弦的斜率趋向 $\infty$ 即可.

\noindent\textbf{证}\quad 记 $l=b-a$，又 $c_n=a+\dfrac{l}{n}=a+\dfrac{b-a}{n}$，则当 $n\to\infty$ 时，$c_n\to a+$.

由于 $\lim\limits_{x\to a+}f(x)=\infty$，可取 $x_1\in(a,c_1)$，使得 $x_1$ 与 $a$ 充分接近，以致 $|f(x_1)-f(c_1)|>l$. 从而应用 Lagrange 定理，$\exists\xi_1\in(x_1,c_1)$，使得
\[
  |f'(\xi_1)|=\left|\frac{f(x_1)-f(c_1)}{x_1-c_1}\right|>
  \frac{|f(x_1)-f(c_1)|}{l}>1.
\]
同理，对于任意给定的 $n\in\mathbb N$，$\exists x_n\in(a,c_n)$，使得 $x_n$ 与 $a$ 充分接近，以致 $|f(x_n)-f(c_n)|>nl$. 从而 $\exists\xi_n\in(x_n,c_n)$，使得
\[
  |f'(\xi_n)|=\left|\frac{f(x_n)-f(c_n)}{x_n-c_n}\right|>
  \frac{|f(x_n)-f(c_n)|}{l}>n.
\]
如此，我们得到序列 $\{\xi_n\}:a<x_n<\xi_n<c_n\to a,\ |f'(\xi_n)|\ge n\to\infty$，故 $\overline{\lim\limits_{x\to a+}}f'(x)=\infty$.

\noindent\textbf{注}\quad $f(x)$ 可微，$\lim\limits_{x\to a+}f(x)=\infty$，一般不能推出 $\lim\limits_{x\to a+}f'(x)=\infty$. 读者不妨考虑函数 $f(x)=\dfrac1x+\cos\dfrac1x$ 当 $x\to0+$ 时的情况.

\noindent\textbf{*例 3.2.12}\quad 设 $f(x)$ 在区间 $[0,1]$ 上可微，$f(0)=0,\ f(1)=1,\ k_1,k_2,\cdots,k_n$ 为 $n$ 个正数. 证明在区间 $[0,1]$ 上存在一组互不相等的数 $x_1,x_2,\cdots,x_n$，使得
\[
  \sum_{i=1}^n\frac{k_i}{f'(x_i)}=\sum_{i=1}^n k_i.
\]
（中国科学技术大学）

\noindent\textbf{分析}\quad （将等式变形.）记右端 $\sum\limits_{i=1}^n k_i=m$. 用此常数同除等式两端，得
\begin{equation}
  \frac{\lambda_1}{f'(x_1)}+\frac{\lambda_2}{f'(x_2)}+\cdots+\frac{\lambda_n}{f'(x_n)}=1,
\tag{1}
\end{equation}
其中 $\lambda_i=\dfrac{k_i}{m},0<\lambda_i<1,\lambda_1+\lambda_2+\cdots+\lambda_n=1$. 问题是：寻求 $n$ 个不同的数 $x_1,x_2,\cdots,x_n$，使得相应的导数 $f'(x_i)\ (i=1,2,\cdots,n)$ 满足方程 (1). 根据 Lagrange 定理，这等于说：要找 $n$ 个不同的弦，使相应的斜率 $\tan\alpha_i\ (i=1,2,\cdots,n)$ 满足 $\sum\limits_{i=1}^n\dfrac{\lambda_i}{\tan\alpha_i}=1$. 若把 $\lambda_i$ 看成弦在 $y$ 轴上的投影，则 $\dfrac{\lambda_i}{\tan\alpha_i}$ 等于该弦在 $x$ 轴上的投影. 因此，问题在于找 $n$ 个不同的弦，使它们在 $y$ 轴上的投影分别为 $\lambda_1,\lambda_2,\cdots,\lambda_n$，在 $x$ 轴上的投影之和为 1.

注意到 $f(0)=0,\ f(1)=1$. 为此，我们在 $y$ 轴上对区间 $[0,1]$ 作分划：$0=A_0<A_1<\cdots<A_n=1$，使得 $A_i-A_{i-1}=\lambda_i\ (i=1,2,\cdots,n)$（如图 3.2.2）. 首先过 $A_1$ 作水平直线与曲线交于点 $B_1$（由于 $f(x)$ 连续，这种交点一定存在）. 于是 $OB_1$ 为第一个所需要的弦. 同理，逐次过 $A_2,A_3,\cdots,A_n$ 作水平线，相继可得到交点 $B_2,B_3,\cdots,B_n$. 这时 $OB_1,B_1B_2,B_2B_3,\cdots,B_{n-1}B_n$ 正是要找的弦，它们在 $y$ 轴上的投影为 $\lambda_1,\lambda_2,\cdots,\lambda_n$，在 $x$ 轴上的投影和为 1.

\begin{center}
  \includegraphics[width=0.42\textwidth]{3-2-2.png}

  图 3.2.2
\end{center}

\noindent\textbf{证}\quad 记 $m=\sum\limits_{i=1}^n k_i,\lambda_i=\dfrac{k_i}{m}\ (i=1,2,\cdots,n)$，则 $0<\lambda_i<1,\lambda_1+\lambda_2+\cdots+\lambda_n=1$. 因 $f(0)=0,\ f(1)=1,\ f(x)$ 在 $[0,1]$ 上连续，故由介值性，$\exists c_1\in(0,1)$ 使得 $f(c_1)=\lambda_1$. 又由 $\lambda_1<\lambda_1+\lambda_2<1$，知 $\exists c_2\in(c_1,1)$ 使得 $f(c_2)=\lambda_1+\lambda_2$. 如此继续下去，顺次找到
\[
  0<c_1<c_2<\cdots<c_{n-1}<c_n=1,
\]
使得
\[
  f(c_i)=\sum_{k=1}^i\lambda_k\quad(i=1,2,\cdots,n).
\]
应用 Lagrange 定理，$\exists x_i\in(c_{i-1},c_i)\ (c_0=0)$，使得
\[
  f'(x_i)=\frac{f(c_i)-f(c_{i-1})}{c_i-c_{i-1}}=
  \frac{\lambda_i}{c_i-c_{i-1}},
\]
即 $\dfrac{\lambda_i}{f'(x_i)}=c_i-c_{i-1}$. 从而
\[
  \sum_{i=1}^n\frac{\lambda_i}{f'(x_i)}
  =\sum_{i=1}^n(c_i-c_{i-1})=c_n-c_0=1.
\]
将 $\lambda_i=\dfrac{k_i}{m}$ 代入，即
\[
  \sum_{i=1}^n\frac{k_i}{f'(x_i)}=\sum_{i=1}^n k_i.
\]

\noindent\textbf{$\star$例 3.2.13}\quad 设函数 $f(x)$ 在闭区间 $[a,b]$ 上连续，在开区间 $(a,b)$ 内可导，又 $f(x)$ 不是线性函数，且 $f(b)>f(a)$，试证 $\exists\xi\in(a,b)$，使得
\[
  f'(\xi)>\frac{f(b)-f(a)}{b-a}.
\]
（上海交通大学）

\noindent\textbf{证}\quad 过点 $(a,f(a))$ 与 $(b,f(b))$ 的线性函数为
\[
  y=f(a)+\frac{f(b)-f(a)}{b-a}(x-a).
\]
因 $f(x)$ 不是线性函数，所以
\begin{equation}
  F(x)\equiv f(x)-f(a)-\frac{f(b)-f(a)}{b-a}(x-a)\ne0.
\tag{1}
\end{equation}
我们的问题是：要证明 $\exists\xi\in(a,b)$，使得
\begin{equation}
  F'(\xi)=f'(\xi)-\frac{f(b)-f(a)}{b-a}>0.
\tag{2}
\end{equation}
（按弦线法，这就等于要找曲线 $y=F(x)$ 的一根弦，使其斜率大于 0.）由已知条件可知：函数 $F(x)$ 在 $[a,b]$ 上连续，在 $(a,b)$ 内可导，$F(a)=F(b)$，满足 Lagrange 定理的条件. 由 (1) 知：$\exists x_0\in(a,b)$，使得 $F(x_0)\ne0$. 亦即：要么 $F(x_0)>0$，要么 $F(x_0)<0$. 这两种情况分别对应图 3.2.3(a) 和 (b). 要让 $F'(\xi)>0$，$\xi$ 应选在上升段. 即：当 $F(x_0)>0$ 时，$\xi$ 应在 $x_0$ 的左侧（如图 3.2.3(a)）；当 $F(x_0)<0$ 时，$\xi$ 应在 $x_0$ 的右侧（如图 3.2.3(b)）.

例如，$F(x_0)>0$，在 $[a,x_0]$ 上应用 Lagrange 定理，$\exists\xi\in(a,x_0)\subset(a,b)$，使得
\[
  F'(\xi)=\frac{F(x_0)-F(a)}{x_0-a}=\frac{F(x_0)}{x_0-a}>0.
\]
证毕.

\begin{center}
  \includegraphics[width=0.62\textwidth]{3-2-3.png}

  图 3.2.3
\end{center}

\noindent\textbf{**练习}\quad 设 $f(x)$ 在 $[a,b]$ 上连续，在 $(a,b)$ 内可导，$f(x)$ 既不是常数也不是一次函数. 试证：
\[
  \exists\xi\in(a,b),\text{ 使得 }|f'(\xi)|>\left|\frac{f(b)-f(a)}{b-a}\right|.
\]
（华中师范大学）

\noindent\textbf{提示}\quad 1° 若 $f(b)=f(a)$，这样的 $\xi$ 当然存在（否则 $f'(x)\equiv0,\ f(x)=$ 常数，矛盾）；

2° 若 $f(b)>f(a)$，重复例 3.2.13 证法可得.

3° 若 $f(b)<f(a)$，只需在例 3.2.13 的式 (2) 里将“$>0$”改为“$<0$”，相应地，图 3.2.3 中 $\xi$ 取在 $F(x)$ 下降段即可.

\noindent\textbf{*例 3.2.14}\quad 证明：若函数 $f(x)$ 在 $(0,+\infty)$ 内可微，且当 $x\to+\infty$ 时有 $f(x)=o(x)$，则 $\underline{\lim\limits_{x\to+\infty}}|f'(x)|=0$.

\noindent\textbf{证}\quad 要证明 $\underline{\lim\limits_{x\to+\infty}}|f'(x)|=0$，只要证明存在序列 $\{\xi_n\}\to+\infty$，使得 $\lim\limits_{n\to\infty}f'(\xi_n)=0$.

为此我们要在任意大的 $\Delta_n>0$ 之右侧，找一根弦，使其斜率
\[
  \left|\frac{f(\beta_n)-f(\alpha_n)}{\beta_n-\alpha_n}\right|<\varepsilon_n,
\]
其中 $\varepsilon_n$ 是任意小量. 事实上，已知 $x\to+\infty$ 时 $f(x)=o(x)$，即 $\lim\limits_{x\to+\infty}\dfrac{f(x)}x=0$，从而 $\forall\alpha>0$，
\[
  \lim_{x\to+\infty}\frac{f(x)-f(\alpha)}{x-\alpha}
  =\lim_{x\to+\infty}\left[\frac{f(x)}x\left(1+\frac{\alpha}{x-\alpha}\right)-\frac{f(\alpha)}{x-\alpha}\right]=0.
\]
故对每个 $\varepsilon_n=\dfrac1n,\alpha_n=n$，总存在 $\beta_n>\alpha_n$，使得
\[
  \left|\frac{f(\beta_n)-f(\alpha_n)}{\beta_n-\alpha_n}\right|<\varepsilon_n=\frac1n.
\]
进而由 Lagrange 定理，$\exists\xi_n\in(\alpha_n,\beta_n)$ 使得
\[
  |f'(\xi_n)|=\left|\frac{f(\beta_n)-f(\alpha_n)}{\beta_n-\alpha_n}\right|<\frac1n.
\]
问题获证.

\noindent\textbf{b. 利用有限增量公式导出新的中值公式}

\noindent\textbf{要点}\quad 借助不同的辅助函数，可由有限增量公式
\[
  f(b)-f(a)=f'(\xi)(b-a),\quad \xi\in(a,b),
\]
导出新的中值公式.

\noindent\textbf{$\star$例 3.2.15}\quad 设 $f(x)$ 在 $[a,b]$ 上连续，在 $(a,b)$ 内有二阶导数，试证存在 $c\in(a,b)$ 使得
\begin{equation}
  f(b)-2f\left(\frac{a+b}{2}\right)+f(a)=\frac{(b-a)^2}{4}f''(c).
\tag{1}
\end{equation}
（南开大学）

\noindent\textbf{证}\quad （式 (1) 左端）
\[
\begin{aligned}
& f(b)-2f\left(\frac{a+b}{2}\right)+f(a)\\
={}&\left[f(b)-f\left(\frac{a+b}{2}\right)\right]
-\left[f\left(\frac{a+b}{2}\right)-f(a)\right]\\
={}&\left[f\left(\frac{a+b}{2}+\frac{b-a}{2}\right)-f\left(\frac{a+b}{2}\right)\right]
-\left[f\left(a+\frac{b-a}{2}\right)-f(a)\right].
\end{aligned}
\]
如图 3.2.4，作辅助函数 $\varphi(x)=f\left(x+\dfrac{b-a}{2}\right)-f(x)$，则
\[
\begin{aligned}
\text{上式}
={}&\varphi\left(\frac{a+b}{2}\right)-\varphi(a)
=\varphi'(\xi)\cdot\left(\frac{a+b}{2}-a\right)
\quad\left(\xi\in\left(a,\frac{a+b}{2}\right)\right)\\
={}&\varphi'(\xi)\frac{b-a}{2}
=\left[f'\left(\xi+\frac{b-a}{2}\right)-f'(\xi)\right]\frac{b-a}{2}\\
={}&f''\left(\xi+\theta\frac{b-a}{2}\right)\cdot\frac{b-a}{2}\cdot\frac{b-a}{2}\quad(\theta\in(0,1))\\
={}&f''(c)\cdot\frac{(b-a)^2}{4}
\quad\left(c=\xi+\theta\frac{b-a}{2}\in(a,b)\right).
\end{aligned}
\]

\begin{center}
  \includegraphics[width=0.38\textwidth]{3-2-4.png}

  图 3.2.4
\end{center}

\noindent\textbf{例 3.2.16}\quad 设函数 $f(x)$ 在 $[0,+\infty)$ 上连续，而且有连续的右导数 $f'_+(x)$，试证明导数 $f'(x)$ 存在，且 $f'(x)=f'_+(x)$.（河南师范大学）

\noindent\textbf{分析}\quad $\forall x_0\in(0,+\infty)$，要证 $f'(x_0)$ 存在且等于 $f'_+(x_0)$，只要证明 $f'_-(x_0)=f'_+(x_0)$. 即要证明 $\forall\varepsilon>0,\exists\delta>0$，当 $0<h<\delta$ 时，有
\begin{equation}
  \left|\frac{f(x_0-h)-f(x_0)}{-h}-f'_+(x_0)\right|<\varepsilon.
\tag{1}
\end{equation}
可见，若能证明 $f'_+(x)$ 也有类似的增量公式
\begin{equation}
  f(x-h)-f(x)=f'_+(x-\theta h)\cdot(-h)\quad(0<\theta<1),
\tag{2}
\end{equation}
那么式 (1) 可变成
\begin{equation}
  |f'_+(x_0-\theta h)-f'_+(x_0)|<\varepsilon.
\tag{3}
\end{equation}
根据已知条件：$f'_+(x)$ 连续，当 $h$ 充分小时，式 (3) 确能成立，从而式 (1) 获证. 因此问题归结为证明右导数也有增量公式 (2) 成立. 为此令
\[
  \varphi(t)=f(x_0-th)-f(x_0)-[f(x_0-h)-f(x_0)]t,
\]
则 $\varphi(0)=\varphi(1)=0$. 利用例 3.1.5，可知 $\exists\theta\in(0,1)$ 使得 $\varphi'_+(\theta)=0$，此即式 (2) 成立.

将有限增量公式应用到具体的函数，也可得到具体中值公式.

\noindent\textbf{例 3.2.17}\quad 设 $a,b>0$，试证 $\exists\xi\in(a,b)$，使得
\begin{equation}
  ae^b-be^a=(1-\xi)e^\xi(a-b).
\tag{1}
\end{equation}

\noindent\textbf{提示}\quad 将式 (1) 改写成
\begin{equation}
  \frac1b e^b-\frac1a e^a=
  \left(1-\frac1\eta\right)e^{1/\eta}\left(\frac1b-\frac1a\right)
  \quad\left(\text{其中 }\frac1\eta=\xi\right),
\tag{2}
\end{equation}
令 $f(x)=xe^{1/x}$，在 $\left[\dfrac1a,\dfrac1b\right]$ 上应用有限增量公式.

\noindent\textbf{**练习}\quad 设 $f(x)$ 在 $[a,b]$ 上连续，在 $(a,b)$ 有二阶导数，在平面上两点 $M(a,f(a))$ 和 $P(b,f(b))$ 的连线 $\overline{MP}$ 与曲线 $y=f(x)$ 相交于点 $R(c,f(c))\ (a<c<b)$. 试证：$\exists\xi\in(a,b)$，使得 $f''(\xi)=0$.（华中师范大学）

\noindent\textbf{提示}\quad 如图 3.2.5，依据 Lagrange 定理，$\exists\xi_1\in(a,c)$ 和 $\xi_2\in(c,b)$，使得
\[
  f'(\xi_1)=\frac{f(c)-f(a)}{c-a}=\frac{f(b)-f(c)}{b-c}=f'(\xi_2);
\]
再在 $[\xi_1,\xi_2]$ 上用 Rolle 定理.

\begin{center}
  \includegraphics[width=0.38\textwidth]{3-2-5.png}

  图 3.2.5
\end{center}

\noindent\textbf{c. 作为函数的变形}

\noindent\textbf{要点}\quad 若 $f(x)$ 在 $[a,b]$ 上连续，$(a,b)$ 内可微，则在 $[a,b]$ 上
\[
  f(x)=f(x_0)+f'(\xi)(x-x_0)\quad(\xi\text{ 在 }x\text{ 与 }x_0\text{ 之间}).
\]
这可视为函数 $f(x)$ 的一种变形，它给出了函数与导数的一种关系. 我们可以用它来研究函数的性质. 此式相当于 $f(x)$ 的 Taylor 展开式至 0 次项.

\noindent\textbf{$\star$例 3.2.18}\quad 设 $f(x)$ 在 $[0,+\infty)$ 上可微，$f(0)=0$，并设有实数 $A>0$，使得 $|f'(x)|\le A|f(x)|$ 在 $[0,+\infty)$ 上成立，试证明在 $[0,+\infty)$ 上 $f(x)\equiv0$.（广西大学）

\noindent\textbf{证 I}\quad 因 $f(x)$ 在 $[0,+\infty)$ 上可微，$f(0)=0$，利用 Lagrange 定理，
\[
  |f(x)|=|f(0)+f'(\xi_1)(x-0)|=|f'(\xi_1)x|\le A|f(\xi_1)|x.
\]
当限制 $x\in\left(0,\dfrac1{2A}\right]$ 时，则得
\[
  |f(x)|\le\frac12|f(\xi_1)|,\quad 0<\xi_1<x.
\]
重复使用此式可得
\[
  |f(x)|\le\frac12|f(\xi_1)|\le\frac14|f(\xi_2)|\le\cdots\le\frac1{2^n}|f(\xi_n)|,
\]
这里 $0<\xi_n<\xi_{n-1}<\cdots<\xi_1<x\le\dfrac1{2A}$.

由 $f(x)$ 的连续性，$\exists M>0$，使得在 $\left[0,\dfrac1{2A}\right]$ 上 $|f(x)|\le M$. 故
\[
  |f(x)|\le\frac{M}{2^n}\quad(n=1,2,\cdots).
\]
从而在 $\left[0,\dfrac1{2A}\right]$ 上 $f(x)\equiv0$. 用数学归纳法，可证在一切 $\left[\dfrac{i-1}{2A},\dfrac{i}{2A}\right]\ (i=1,2,\cdots)$ 上恒有 $f(x)\equiv0$. 所以 $[0,+\infty)$ 上 $f(x)\equiv0$.

\noindent\textbf{证 II}\quad 因 $|f(x)|$ 在 $\left[0,\dfrac1{2A}\right]$ 上连续，故 $\exists x_1\in\left[0,\dfrac1{2A}\right]$ 使得
\[
  |f(x_1)|=\max_{0\le x\le1/(2A)}|f(x)|=M.
\]
于是
\[
  M=|f(x_1)|=|f(0)+f'(\xi)(x_1-0)|=|f'(\xi)x_1|
  \le A|f(\xi)|\cdot x_1\le\frac12|f(\xi)|\le\frac12M.
\]
所以 $M=0,\ f(x)\equiv0\left(\text{于 }\left[0,\dfrac1{2A}\right]\text{ 上}\right)$. 以下同证 I.

\noindent\textbf{证 III}\quad （反证法）若 $f(x)\not\equiv0$，则 $\exists x_0\in(0,+\infty)$，使得 $f(x_0)\ne0$，不妨设 $f(x_0)>0$（$f(x_0)<0$ 类似可证）. 记 $x_1\equiv\inf\{x\mid (x,x_0)\text{ 上 }f(x)>0\}$，由连续函数局部保号性，只能 $f(x_1)=0,(x_1,x_0)$ 内 $f(x)>0$. 令
\begin{equation}
  g(x)=\ln f(x)\quad(\text{当 }x\in(x_1,x_0)\text{ 时}),
\tag{1}
\end{equation}
则
\[
  |g'(x)|=\left|\frac{f'(x)}{f(x)}\right|\le A.
\]
故 $g(x)$ 在有限区间 $(x_1,x_0)$ 上有界. 但 $\lim\limits_{x\to x_1+}f(x)=f(x_1)=0$. 由 (1) 知 $\lim\limits_{x\to x_1+}g(x)=-\infty$. 矛盾.

\noindent\textbf{练习}\quad 设 $f(x),g(x)$ 在 $[a,b]$ 上连续，$g(x)$ 在 $(a,b)$ 内可微，且 $g(a)=0$. 若有实数 $\lambda\ne0$，使得
\[
  |g(x)-f(x)+\lambda g'(x)|\le |g(x)|,\quad x\in(a,b)
\tag{1}
\]
成立，试证：$g(x)\equiv0$.（浙江大学）

\noindent\textbf{提示}\quad 可利用上例.

\noindent\textbf{例 3.2.19}\quad 设 $[0,a]$ 上 $|f''(x)|\le M,\ f(x)$ 在 $(0,a)$ 内取最大值，试证
\[
  |f'(0)|+|f'(a)|\le Ma.
\]

\noindent\textbf{提示}\quad 对 $f'(x)$ 应用 Lagrange 定理（0 次 Taylor 展开式），将 $f'(0),f'(a)$ 在 $f(x)$ 的最大值点处展开.

\noindent\textbf{例 3.2.20}\quad 证明：若函数 $f(x)$ 在 $(0,+\infty)$ 内可微，且 $\lim\limits_{x\to+\infty}f'(x)=0$，则
\[
  \lim_{x\to+\infty}\frac{f(x)}x=0.
\]

\noindent\textbf{证}\quad 要证 $\lim\limits_{x\to+\infty}\dfrac{f(x)}x=0$，即要证：$\forall\varepsilon>0,\exists\Delta>0$，使得 $x>\Delta$ 时有
\begin{equation}
  \left|\frac{f(x)}x\right|<\varepsilon.
\tag{1}
\end{equation}
已知 $\lim\limits_{x\to+\infty}f'(x)=0$. 所以对此 $\varepsilon>0$，$\exists A>0$，当 $x>A$ 时有
\begin{equation}
  |f'(x)|<\frac{\varepsilon}{2}.
\tag{2}
\end{equation}
由 Lagrange 定理，对 $x>A$，$\exists\xi:A<\xi<x$ 使得 $f(x)=f(A)+f'(\xi)(x-A)$，故
\begin{equation}
  \left|\frac{f(x)}x\right|\le\left|\frac{f(A)}x\right|+|f'(\xi)|\cdot\frac{x-A}{x}
  <\frac{|f(A)|}{x}+\frac{\varepsilon}{2}.
\tag{3}
\end{equation}
但 $x\to+\infty$ 时 $\dfrac{|f(A)|}{x}\to0$. 可见要式 (1) 成立，只要取 $\Delta=\max\left\{A,\dfrac{2|f(A)|}{\varepsilon}\right\}$，则当 $x>\Delta$ 时，从 (3) 可推得 (1) 成立.

\noindent\textbf{注}\quad 本例如应用 L'Hospital 法则（例 3.2.30），结果十分明显.

\noindent\textbf{**练习 1}\quad 已知函数 $f(x)$ 在 $[a,+\infty)$ 上有界、可导，试证：在 $[a,+\infty)$ 上存在数列 $\{x_n\}\to+\infty$，使得 $\lim\limits_{n\to\infty}f'(x_n)=0$.（南开大学）

\noindent\textbf{提示}\quad 利用 Lagrange 定理，$\exists x_n\in(n,2n)$，使得
\[
  |f'(x_n)|=\frac{|f(2n)-f(n)|}{n}\le\frac{2M}{n}\quad(\forall n\in\mathbb N)\quad(M\text{ 代表 }f(x)\text{ 的界}).
\]
因此 $\lim\limits_{n\to\infty}f'(x_n)=0$.

\noindent\textbf{***练习 2}\quad 设函数 $f(x)$ 在 $[0,1]$ 上连续，在 $(0,1)$ 内可导，且 $\exists M>0$ 使得
\begin{equation}
  |xf'(x)-f(x)|\le x^2M.
\tag{1}
\end{equation}
试证：1）$\dfrac{f(x)}x$ 在 $[0,1]$ 上一致连续；2）极限 $\lim\limits_{x\to0+}\dfrac{f(x)}x$ 存在.（华中师范大学）

\noindent\textbf{提示}\quad 记 $F(x)=\dfrac{f(x)}x$，则 $|F'(x)|\le M$. 对 $F(x)$ 应用 Lagrange 公式，可知 $F(x)=\dfrac{f(x)}x$ 满足 Lipschitz 条件，故 $\dfrac{f(x)}x$ 在 $[0,1]$ 上一致连续. 根据 Cauchy 准则，$\lim\limits_{x\to0+}\dfrac{f(x)}x$ 存在（见例 2.2.6）.

\noindent\textbf{new 练习 3}\quad 设 $f(x)$ 在 $[a,+\infty)$ 上可导，且 $f'(x)\ge c>0$（$c$ 为常数）. 试证：

1）$\lim\limits_{x\to+\infty}f(x)=+\infty$；\quad 2）$f(x)$ 在 $[a,+\infty)$ 上有最小值.（南开大学）

\noindent\textbf{提示}\quad 1）任意取定 $x_0>a$，那么当 $x_0<x$ 时，有 $\xi_x\in(x_0,x)$ 使得
\begin{equation}
  f(x)=f(x_0)+f'(\xi_x)(x-x_0)\ge f(x_0)+c(x-x_0)\to+\infty\quad(x\to+\infty).
\tag{1}
\end{equation}

2）记 $m=\min\{f(x)\mid x\in[a,a+1]\}$，由式 (1)，$\exists\Delta>0$，使得 $x>\Delta$ 时，恒有 $f(x)>m$. 这时 $f(x)$ 在 $[a,\Delta]$ 上有最小值，也就是 $f(x)$ 在 $[a,+\infty)$ 上有最小值.

\noindent\textbf{new *练习 4}\quad 设 $f(x)$ 在区间 $(0,1]$ 上有连续导数，且 $\lim\limits_{x\to0+}\sqrt{x}f'(x)=0$，试证：$f(x)$ 在 $(0,1]$ 上一致连续.（仿华中科技大学）

\noindent\textbf{证 I}\quad 已知 $f(x)$ 在区间 $(0,1]$ 上可导，当然连续，因此要证明 $f(x)$ 在 $(0,1]$ 上一致连续，只需证明极限 $\lim\limits_{x\to0+}f(x)$ 存在（例 2.2.6）. 根据 Cauchy 准则，只需证明：$\forall\varepsilon>0,\exists\delta:0<\delta<1$，当 $0<x_1<x_2<\delta$ 时，有
\begin{equation}
  |f(x_1)-f(x_2)|<\varepsilon.
\tag{1}
\end{equation}
根据题设，$\lim\limits_{x\to0+}\sqrt{x}f'(x)=0$，即 $\forall\varepsilon>0,\exists\delta_1\in(0,1)$，当 $0<x<\delta_1$ 时，有
\begin{equation}
  2\sqrt{x}|f'(x)|<\varepsilon.
\tag{2}
\end{equation}
在 $(0,\delta_1)$ 内任取两点 $x_1,x_2:0<x_1<x_2<\delta_1$，令 $x=t^2$. 记 $t_1=\sqrt{x_1},t_2=\sqrt{x_2}$，则式 (1) 左端
\begin{equation}
  |f(x_1)-f(x_2)|=|f(t_1^2)-f(t_2^2)|.
\tag{3}
\end{equation}
应用 Lagrange 定理，$\exists\xi\in(t_1,t_2)\subset(0,1)$，使得
\begin{equation}
\begin{aligned}
|f(x_1)-f(x_2)|
={}& |f'(\xi^2)2\xi|\,|t_1-t_2|\\
<{}& 2\sqrt{\eta}|f'(\eta)|<\varepsilon
\end{aligned}
\tag{4}
\end{equation}
（令 $\eta=\xi^2$）.
可见，只要先用式 (2) 找出 $\delta_1$，然后取 $\delta$，使得 $\delta^2\le\delta_1<1$，式 (1) 就成立. 表明 $\lim\limits_{x\to0+}f(x)$ 存在，故 $f(x)$ 在 $(0,1]$ 上一致连续.

\noindent\textbf{证 II}\quad 记 $F(t)=f(t^2)$，则
\[
  F'(t)=(f(t^2))'=f'(t^2)\cdot2t=f'(x)\cdot2\sqrt{x}\to0\quad(\text{当 }x\to0+\text{ 时}).
\]
可见 $F'(t)=2tf'(t^2)$ 不仅在 $(0,1]$ 上连续，而且在原点有右极限，故 $F'(t)$ 在 $(0,1]$ 上有界，说明 $F(t)$ 在 $(0,1]$ 上满足 Lipschitz 条件，从而 $F(t)$ 在 $(0,1]$ 上一致连续（见例 2.2.1 前的要点）.

另一方面，$t=\sqrt{x}$ 在 $(0,1]$ 上连续，由 Cantor 定理知 $t=\sqrt{x}$ 在 $[0,1]$（当然在 $(0,1]$）上一致连续，于是复合函数 $f(x)=F(\sqrt{x})$ 也在 $(0,1]$ 上一致连续.

\noindent\textbf{$\star$例 3.2.21}\quad 设 $\varphi(x)=\int_0^x\dfrac{\ln(1-t)}t\,dt$ 在 $-1<x<1$ 有意义，证明：$\varphi(x)+\varphi(-x)=\dfrac12\varphi(x^2)$.（北京航空航天大学）

\noindent\textbf{证}\quad 问题等价于要证明函数
\[
  f(x)\equiv\varphi(x)+\varphi(-x)-\frac12\varphi(x^2)\equiv0.
\]
事实上，
\[
  f'(x)=\varphi'(x)-\varphi'(-x)-x\varphi'(x^2),
\]
而 $\varphi'(x)=\dfrac{\ln(1-x)}x$，故
\[
  f'(x)=\frac{\ln(1-x)}x+\frac{\ln(1+x)}x-\frac{\ln(1-x^2)}{x^2}\cdot x=0.
\]
因此，$f(x)\equiv c$. 但由 $\varphi(0)=0$ 知 $f(0)=0$，所以 $c=0,\ f(x)\equiv0$.

\subsection*{三、导数的两大特性}

导数有如下两个重要特性：1° 导数无第一类间断点；2° 导数的介值性. 下面具体介绍这两个特性的含义、证明和应用.

\noindent\textbf{a. 导数无第一类间断点}

\noindent\textbf{$\star$例 3.2.22（导数无第一类间断点）}\quad 设函数 $f(x)$ 在 $(a,b)$ 内处处有导数 $f'(x)$，证明 $(a,b)$ 中的点或者为 $f'(x)$ 的连续点，或者为 $f'(x)$ 的第二类间断点.（南京大学）

\noindent\textbf{证}\quad 因为 $f(x)$ 在 $(a,b)$ 内处处可导，所以 $\forall x_0\in(a,b)$，
\begin{equation}
\begin{aligned}
f'(x_0)=f'_+(x_0)
={}&\lim_{x\to x_0+}\frac{f(x)-f(x_0)}{x-x_0}\quad(\text{应用 Lagrange 定理})\\
={}&\lim_{x\to x_0+}\frac{f'(\xi)(x-x_0)}{x-x_0}
=\lim_{x\to x_0+}f'(\xi)\quad(x_0<\xi<x).
\end{aligned}
\tag{1}
\end{equation}
故 $f'(x)$ 在 $x_0$ 处有右极限时，必有
\begin{equation}
  f'(x_0)=\lim_{\xi\to x_0+}f'(\xi)=f'(x_0+0).
\tag{2}
\end{equation}
同理可证，若 $f'(x)$ 在 $x_0$ 处有左极限，则有 $f'(x_0)=f'(x_0-0)$.

因此在 $(a,b)$ 内任一点处，除非至少有一侧 $f'(x)$ 无极限（这时 $f'(x)$ 为第二类间断），不然 $f'(x)$ 在此处连续，即 $f'(x_0-0)=f'(x_0)=f'(x_0+0)$.

\noindent\textbf{注}\quad 1° 所谓导数无第一类间断点，是在导数处处存在的前提下讲的. 例如 $f(x)=|x|$，当 $x>0$ 时 $f'(x)=1$；当 $x<0$ 时 $f'(x)=-1$；$x=0$ 处便是第一类间断点. 这与我们的结论不矛盾，因 $f(x)$ 在 $x=0$ 处不可导.

2° 从式 (1) 到式 (2)，用到复合函数求极限的知识.

\noindent\textbf{复合函数的极限}

\noindent\textbf{命题}\quad 设 $y=f(u),u=g(x)$ 可以组成复合函数，已知 $\lim\limits_{x\to x_0}g(x)=u_0,\lim\limits_{u\to u_0}f(u)=A$，若 $x\to x_0$ 的过程中 $g(x)$ 始终保持 $g(x)\ne u_0$，则复合函数的极限 $\lim\limits_{x\to x_0}f(g(x))=A$；否则命题不成立.

例如，
\[
y=f(u)=
\begin{cases}
1, & u\ne0,\\
0, & u=0,
\end{cases}
\quad\text{及}\quad
u=g(x)=
\begin{cases}
x, & x\text{ 为无理点},\\
0, & x\text{ 为有理点},
\end{cases}
\]
$x\to0$ 时 $u\to0$，$u\to0$ 时 $y=f(u)\to1$，但是复合后的极限 $\lim\limits_{x\to0}f(g(x))$ 不存在，因为
\[
f(g(x))=
\begin{cases}
1, & x\text{ 为无理点},\\
0, & x\text{ 为有理点}.
\end{cases}
\]
条件 $g(x)\ne u_0$ 之所以如此重要，是因为：内层函数 $g(x)$ 在求极限的过程中，$g(x)$ 有可能无穷次等于 $u_0$，但外层极限 $\lim\limits_{u\to u_0}f(u)=A$ 在过程中 $u$ 不允许等于 $u_0$.

本例式 (1) 中的极限 $\lim\limits_{x\to x_0+}f'(\xi)$ 实为复合函数的极限，$\xi$ 依赖于 $x$. 因 $x_0<\xi<x$，故当 $x\to x_0+$ 时，$x_0\ne\xi\to x_0$，符合上面命题条件，故式 (2) 成立.

假如外层函数连续，则 $\lim\limits_{x\to x_0}f(g(x))=f\left(\lim\limits_{x\to x_0}g(x)\right)$，就不必假定极限过程中 $g(x)\ne u_0$，也是成立的.

\noindent\textbf{例 3.2.23}\quad 若 $f(x)$ 在 $(a,b)$ 内可导，导函数 $f'(x)$ 在 $(a,b)$ 内单调，则 $f'(x)$ 在 $(a,b)$ 内连续.（武汉大学）

\noindent\textbf{提示}\quad 见例 3.2.22.

\noindent\textbf{b. 导数的介值性}

\noindent\textbf{$\star$例 3.2.24（Darboux 定理）}\quad 若函数 $f(x)$ 在区间 $[a,b]$ 上处处可导（端点处指单侧导数），$f'(a)<f'(b)$，则 $\forall c:f'(a)<c<f'(b)$，$\exists\xi\in(a,b)$，使得 $f'(\xi)=c$.（武汉大学，北京师范大学）

\noindent\textbf{证}\quad 1° 作辅助函数
\[
  g(x)=f(x)-cx,
\]
则 $g(x)$ 在 $[a,b]$ 上处处可导，$g'(a)=f'(a)-c<0,\ g'(b)=f'(b)-c>0$，只要能证明 $\exists\xi\in(a,b)$，使得 $g'(\xi)=0$，则 $f'(\xi)=c$.

2° 由于
\[
  \lim_{x\to a+}\frac{g(x)-g(a)}{x-a}=g'(a)<0,
\]
所以 $x>a,\ x$ 与 $a$ 充分接近时，有 $g(x)<g(a)$.

同理，由 $g'(b)>0$，知 $x<b$，且 $x$ 与 $b$ 充分接近时有 $g(x)<g(b)$.
故 $g(x)$ 在端点 $a,b$ 处不取最小值. 但 $g(x)$ 连续，它在闭区间 $[a,b]$ 上有最小值. 所以 $\exists\xi\in(a,b)$，使得 $g(\xi)=\min\limits_{a\le x\le b}g(x)$. 由 Fermat 定理，$g'(\xi)=0$. 证毕.

\noindent\textbf{例 3.2.25}\quad 设函数 $f(x)$ 在区间 $(-\infty,+\infty)$ 上二次可微，且有界，试证存在点 $x_0\in(-\infty,+\infty)$，使得 $f''(x_0)=0$.（武汉大学）

\noindent\textbf{证 I}\quad 若 $f''(x)$ 变号，则由导数的介值性，$\exists\xi\in(-\infty,+\infty)$，使得 $f''(\xi)=0$（下面证明 $f''$ 不会不变号）.

若 $f''(x)$ 不变号，例如 $f''(x)>0$（$f''(x)<0$ 类似可证），则 $f'(x)$ 严增. 取 $x_0$ 使得 $f'(x_0)\ne0$，假如 $f'(x_0)>0$，则当 $x>x_0$，并令 $x\to+\infty$ 时，
\[
  f(x)=f(x_0)+f'(\xi)(x-x_0)>f(x_0)+f'(x_0)(x-x_0)\to+\infty.
\]
若 $f'(x_0)<0$，则当 $x<x_0$，并令 $x\to-\infty$ 时，
\[
  f(x)=f(x_0)+f'(\xi)(x-x_0)>f(x_0)+f'(x_0)(x-x_0)\to+\infty.
\]
与 $f(x)$ 有界矛盾.

\noindent\textbf{证 II}\quad 若 $\exists a,b\ (a\ne b)$，使得 $f'(a)=f'(b)$，则由 Rolle 定理，$\exists x_0$ 位于 $a,b$ 之间，使得 $f''(x_0)=0$. 现证明相反的情况不可能. 事实上，若 $\forall a,b\ (a\ne b)$，恒有 $f'(a)\ne f'(b)$，则由此易证 $f'(x)$ 严格单调.（因为不然的话，必存在 $x_1<x_2<x_3$ 使得
\[
  f'(x_1)<f'(x_2)>f'(x_3)\quad\text{或}\quad f'(x_1)>f'(x_2)<f'(x_3).
\]
由 $f'(x)$ 的介值性，便知：$\exists a\in(x_1,x_2),b\in(x_2,x_3)$ 使得 $f'(a)=f'(b)$. 与前提矛盾.）重复证 I 中后半部分，推知 $f(x)$ 无界，与题设矛盾. 证毕.

\noindent\textbf{注}\quad 证 I 是利用导数的介值性来证明 $f''(x)$ 有零点，虽然题中并未假定 $f''(x)$ 连续.

\noindent\textbf{new 练习 1}\quad 设函数 $f$ 在 $(0,1)$ 上有界、可导，但 $\lim\limits_{x\to0+}f(x)$ 不存在，试证存在数列 $x_n\to0+\ (n\to\infty)$，使得 $\lim\limits_{n\to\infty}f'(x_n)=0$.（北京大学）

\noindent\textbf{提示}\quad 1）在 $x\to0+$ 过程里，$f'(x)$ 无穷多次变号（否则：当 $x$ 充分接近 $0+$ 后，$f'(x)$ 恒正（或恒负），那么 $f$ 严格单调，加之 $f$ 有界，因此 $\lim\limits_{x\to0+}f(x)$ 存在，矛盾）.

2）因为 $f'$ 有介值性，$f'(x)$ 无穷多次变号，故 $f'(x)$ 有零点，其零点组成的数列 $\{x_n\}$ 满足题中的要求.

\noindent\textbf{new 练习 2}\quad 设 $f(x)$ 在 $[a,b]$ 上可导，试证：假若无 $(\alpha,\beta)\subset[a,b]$ 使得 $f'(x)$ 在 $(\alpha,\beta)$ 上有界，则每个 $x\in[a,b]$ 都必是 $f'(x)$ 的第二类间断点.

\noindent\textbf{证}\quad $\forall x\in(a,b)$，$\left\{\left[x-\dfrac1n,x+\dfrac1n\right]\right\}(n=1,2,\cdots)$ 能构成一闭区间套，且 $x$ 是它们的唯一公共点. 而当 $n$ 充分大时，$\left(x-\dfrac1n,x+\dfrac1n\right)\subset[a,b]$. 按假设，$f'(x)$ 在 $\left(x-\dfrac1n,x+\dfrac1n\right)\subset[a,b]$ 上无界，因此 $\exists x_n\in\left(x-\dfrac1n,x+\dfrac1n\right)$ 使得 $|f'(x_n)|>n$. 即 $\lim\limits_{n\to\infty}x_n=x$，但 $\lim\limits_{n\to\infty}|f'(x_n)|=+\infty$. 可见每个 $x\in(a,b)$ 皆是 $f'(x)$ 的第二类间断点. 端点 $a$ 和 $b$ 同样如此.

\noindent\textbf{练习 3}\quad 如下论断是否成立：若 $f(x)$ 在 $[a,b]$ 上可导，则必存在 $(\alpha,\beta)\subset[a,b]$，使得 $f'(x)$ 在 $(\alpha,\beta)$ 上有界（即上题的情况不会发生）.（若回答是，给出证明；否，给出反例.）（作为互动题.）

\subsection*{四、Cauchy 中值定理及 L'Hospital 法则}

\noindent\textbf{a. 推导中值公式}

\noindent\textbf{要点}\quad Cauchy 中值定理：若 $F(x),G(x)$ 在 $[a,b]$ 上连续，在 $(a,b)$ 内可导，$G'(x)\ne0$，则 $\exists\xi\in(a,b)$，使得
\[
  \frac{F(b)-F(a)}{G(b)-G(a)}=\frac{F'(\xi)}{G'(\xi)}.
\]
我们适当选取函数 $F(x),G(x)$，就可以得到新的中值公式.

\noindent\textbf{$\star$例 3.2.26}\quad 设函数 $f(x)$ 在 $(a,b)$ 内可微，$a,b>0$，且 $f(a+0),f(b-0)$ 均存在（为有限数）. 试证明：$\exists\xi\in(a,b)$ 使得
\begin{equation}
  \frac1{a-b}
  \begin{vmatrix}
  a&b\\
  f(a+0)&f(b-0)
  \end{vmatrix}
  =f(\xi)-\xi f'(\xi).
\tag{1}
\end{equation}
（四川师范大学）

\noindent\textbf{分析}\quad 令 $f(a)=f(a+0),f(b)=f(b-0)$，则 $f(x)$ 在 $[a,b]$ 上连续，在 $(a,b)$ 内可导，欲证明的式 (1) 可改写成
\[
  \frac{af(b)-bf(a)}{a-b}=f(\xi)-\xi f'(\xi),
\]
亦即
\begin{equation}
  \frac{\dfrac{f(b)}b-\dfrac{f(a)}a}{\dfrac1b-\dfrac1a}
  =
  \left.\frac{\left(\dfrac{f(x)}x\right)'}{\left(\dfrac1x\right)'}\right|_{x=\xi}.
\tag{2}
\end{equation}
因此，对函数 $F(x)=\dfrac{f(x)}x,\ G(x)=\dfrac1x$，在 $[a,b]$ 上应用 Cauchy 中值定理即得.

\noindent\textbf{注}\quad 把目标式作适当改写，如式 (1) 改写成 (2)，是寻找辅助函数的关键步骤.

\noindent\textbf{例 3.2.27}\quad 设 $f(x)$ 在 $(a,b)$ 内二次可微，试用 Cauchy 中值定理证明：$\forall x,x_0\in(a,b)$，$\exists\xi$ 在 $x$ 与 $x_0$ 之间，使得
\begin{equation}
  f(x)=f(x_0)+f'(x_0)(x-x_0)+\frac12f''(\xi)(x-x_0)^2
\tag{1}
\end{equation}
成立（此即展开到一次幂的 Taylor 公式）.

\noindent\textbf{证}\quad 只证明 $x>x_0$ 的情况（$x<x_0$ 的情况类似可证，$x=x_0$ 的情况显然）. 式 (1) 可改写成
\begin{equation}
  \frac{f(x)-f(x_0)-f'(x_0)(x-x_0)}
  {\dfrac12(x-x_0)^2}=f''(\xi).
\tag{2}
\end{equation}
为了证明式 (2)，只要令
\[
  F(x)=f(x)-f(x_0)-f'(x_0)(x-x_0),\quad
  G(x)=\frac12(x-x_0)^2,
\]
则 $F'(x)=f'(x)-f'(x_0),\ G'(x)=x-x_0$.
注意到 $F(x_0)=G(x_0)=0,\ F'(x_0)=G'(x_0)=0$，两次应用 Cauchy 中值定理，得
\[
\begin{aligned}
\frac{f(x)-f(x_0)-f'(x_0)(x-x_0)}
{\dfrac12(x-x_0)^2}
&=\frac{F(x)}{G(x)}
=\frac{F(x)-F(x_0)}{G(x)-G(x_0)}
=\frac{F'(\eta)}{G'(\eta)}\quad(\eta\in(x_0,x))\\
&=\frac{F'(\eta)-F'(x_0)}{G'(\eta)-G'(x_0)}
=\frac{F''(\xi)}{G''(\xi)}
=f''(\xi)\quad(\xi\in(x_0,\eta)).
\end{aligned}
\]
证毕.

将 Cauchy 中值公式改写成
\[
  F(b)-F(a)=\frac{F'(\xi)}{G'(\xi)}[G(b)-G(a)].
\]
我们看到，若选取不同的函数 $G$，便可把 $F(b)-F(a)$ 表示成不同的形式. 如另取 $G_1(x)$，则 $\exists\xi,\eta\in(a,b)$，使得
\[
  \frac{F'(\xi)}{G'(\xi)}[G(b)-G(a)]
  =
  \frac{F'(\eta)}{G'_1(\eta)}[G_1(b)-G_1(a)].
\]
一般来说，当 $G(x)$ 取 $n$ 个不同的函数（只要满足 Cauchy 中值定理的条件）时，便可得到含 $n$ 个中值的 $n-1$ 个等式.

\noindent\textbf{$\star$例 3.2.28}\quad 设 $f(x)$ 在 $[a,b]$ 上连续，在 $(a,b)$ 内可导 $(0\le a<b)$，$f(a)\ne f(b)$，证明：$\exists\xi,\eta\in(a,b)$ 使得
\begin{equation}
  f'(\xi)=\frac{a+b}{2\eta}f'(\eta).
\tag{1}
\end{equation}
（华中师范大学，吉林大学）

\noindent\textbf{证}\quad （用 $(b-a)$ 乘 (1) 式两端.）式 (1) 等价于
\begin{equation}
  \frac{f'(\xi)}1(b-a)=\frac{f'(\eta)}{2\eta}(b^2-a^2).
\tag{2}
\end{equation}
为证此式，只要取 $F(x)=f(x)$，取 $G(x)=x$ 和 $x^2$，在 $[a,b]$ 上分别应用 Cauchy 中值定理，则
\[
  f(b)-f(a)=\frac{f'(\xi)}1\cdot(b-a)=\frac{f'(\eta)}{2\eta}(b^2-a^2),\quad \xi,\eta\in(a,b).
\]

\noindent\textbf{new 练习}\quad 设 $f(x)$ 在 $[a,b]\ (b>a>0)$ 上连续，在 $(a,b)$ 内可导，证明：在 $(a,b)$ 内存在两点 $\xi,\eta$，使得
\[
  f'(\xi)=\eta^2\frac{f'(\eta)}{ab}.
\]
（华中师范大学）

\noindent\textbf{提示}\quad
\[
  f'(\xi)\stackrel{\text{Lagrange 定理}}{=}\frac{f(b)-f(a)}{b-a}
  \stackrel{\text{Cauchy 中值定理}}{=}
  \frac{f'(\eta)}{\left(-\dfrac{ab}{x}\right)'\big|_{x=\eta}}
  =\eta^2\frac{f'(\eta)}{ab}.
\]

\noindent\textbf{$\star$例 3.2.29}\quad 设 $f(x)$ 在 $[a,b]$ 上连续，在 $(a,b)$ 内可导，$0<a<b$. 证明：在 $(a,b)$ 内存在 $x_1,x_2,x_3$，使得
\begin{equation}
  \frac{f'(x_1)}{2x_1}
  =(b^2+a^2)\frac{f'(x_2)}{4x_2^3}
  =\frac{\ln\dfrac ba}{b^2-a^2}x_3f'(x_3).
\tag{1}
\end{equation}
（四川大学）

\noindent\textbf{提示}\quad 式 (1) 可改写成
\[
  \frac{f'(x_1)}{2x_1}(b^2-a^2)
  =
  \frac{f'(x_2)}{4x_2^3}(b^4-a^4)
  =
  \frac{f'(x_3)}{1/x_3}(\ln b-\ln a).
\]

\noindent\textbf{b. 作为函数与导数的关系}

\noindent\textbf{要点}\quad 由 Cauchy 中值定理可知，若 $F(x),G(x)$ 在某区间 $I$ 内可导，则 $\forall x_1,x_2\in I$，$\exists\xi$ 使得
\[
  \frac{F(x_2)-F(x_1)}{G(x_2)-G(x_1)}=\frac{F'(\xi)}{G'(\xi)}
  \quad(\xi\text{ 在 }x_1\text{ 与 }x_2\text{ 之间}).
\]
即 Cauchy 中值公式给出了函数差分比与导数比的一种关系. 利用 $\xi$ 在 $x_1$ 与 $x_2$ 之间，我们能求解不少问题（虽然 $\xi$ 在 $x_1,x_2$ 之间什么位置不能肯定）.

作为一个典型例子，我们来看 $\dfrac{\infty}{\infty}$ 型 L'Hospital 法则的证明.

\noindent\textbf{$\star$例 3.2.30}\quad 设 $f(x),g(x)$ 在 $(a,b)$ 内可微，且 $\forall x\in(a,b),g'(x)\ne0$，当 $x\to a+$ 时 $g(x)\to\infty$，且 $\lim\limits_{x\to a+}\dfrac{f'(x)}{g'(x)}=A$（有限数，或 $\infty$），证明：$\lim\limits_{x\to a+}\dfrac{f(x)}{g(x)}=A$.

\noindent\textbf{分析}\quad 已知 $\lim\limits_{x\to a+}\dfrac{f'(x)}{g'(x)}=A$，因此保持取 $a<x<x_1$，应用 Cauchy 中值定理，然后令 $x_1\to a+$，知函数差分比
\begin{equation}
  \frac{f(x)-f(x_1)}{g(x)-g(x_1)}=\frac{f'(\xi)}{g'(\xi)}\to A\quad(\text{当 }x_1\to a+\text{ 时}).
\tag{1}
\end{equation}
剩下的问题在于根据 $g(x)\to\infty$（当 $x\to a+$ 时），由差分比 $\dfrac{f(x)-f(x_1)}{g(x)-g(x_1)}\to A$，推出 $\dfrac{f(x)}{g(x)}\to A$（当 $x\to a+$ 时）.

事实上，$f(x)$ 可以改写成
\[
  f(x)=\frac{f(x)-f(x_1)}{g(x)-g(x_1)}(g(x)-g(x_1))+f(x_1),
\]
因此
\begin{equation}
  \frac{f(x)}{g(x)}
  =
  \frac{f(x)-f(x_1)}{g(x)-g(x_1)}
  \left(1-\frac{g(x_1)}{g(x)}\right)
  +\frac{f(x_1)}{g(x)}.
\tag{2}
\end{equation}
1° 若 $A=$ 有限数. 由 (2) 可得
\begin{equation}
  \frac{f(x)}{g(x)}-A
  =
  \left(\frac{f(x)-f(x_1)}{g(x)-g(x_1)}-A\right)
  \left(1-\frac{g(x_1)}{g(x)}\right)
  +\frac{f(x_1)-Ag(x_1)}{g(x)}.
\tag{3}
\end{equation}
保持 $a<x<x_1$，令 $x_1\to a+$，则 $\forall\varepsilon>0,\exists\delta_1>0$，使得当 $a<x<x_1<a+\delta_1$ 时，有
\[
  \left|\frac{f(x)-f(x_1)}{g(x)-g(x_1)}-A\right|<\frac{\varepsilon}{4}.
\]
再将 $x_1$ 固定，令 $x$ 继续趋向 $a+$. 根据 $g(x)\to\infty$（当 $x\to a+$ 时），知 $\exists\delta>0\ (\delta<x_1-a)$，使得 $a<x<\delta$ 时，有
\[
  \left|\frac{g(x_1)}{g(x)}\right|<1,\quad
  \left|\frac{f(x_1)-Ag(x_1)}{g(x)}\right|<\frac{\varepsilon}{2}.
\]
于是由 (3)，
\[
  \left|\frac{f(x)}{g(x)}-A\right|
  \le
  \left|\frac{f(x)-f(x_1)}{g(x)-g(x_1)}-A\right|
  \left(1+\left|\frac{g(x_1)}{g(x)}\right|\right)
  +\left|\frac{f(x_1)-Ag(x_1)}{g(x)}\right|
  <\frac{\varepsilon}{4}\cdot2+\frac{\varepsilon}{2}=\varepsilon.
\]
2° 若 $A=\infty$，则 $x$ 充分接近 $a+$ 时 $f'(x)\ne0$. 并且对 $M=1$，$\exists\delta>0$，当 $a<x<x_1<a+\delta$ 时，有
\[
  \left|\frac{f(x)-f(x_1)}{g(x)-g(x_1)}\right|
  =
  \left|\frac{f'(\xi)}{g'(\xi)}\right|>1
  \quad(x<\xi<x_1),
\]
从而
\[
  |f(x)-f(x_1)|>|g(x)-g(x_1)|\ge |g(x)|-|g(x_1)|\to+\infty
  \quad(x_1\text{固定且令 }x\to a+).
\]
可见
\[
  |f(x)|\to+\infty,\quad f(x)\to\infty\quad(\text{当 }x\to a+\text{ 时}).
\]
由此可利用 1° 中结果，由 $\lim\limits_{x\to a+}\dfrac{g'(x)}{f'(x)}=0$ 得出 $\lim\limits_{x\to a+}\dfrac{g(x)}{f(x)}=0$，从而 $\lim\limits_{x\to a+}\dfrac{f(x)}{g(x)}=\infty$.

\noindent\textbf{注}\quad 1° 这里是 $x\to a+$ 的情况，$x\to a-$ 以及 $x\to+\infty,x\to-\infty$ 的情况亦有类似的结论和证法；由 $x\to a+$ 及 $x\to a-$ 的结论，可知 $x\to a$ 时结论也成立.

2° 本例虽称为 $\dfrac{\infty}{\infty}$ 型的 L'Hospital 法则，实际上对于 $\lim f(x)/g(x)$，条件只要求分母 $g(x)\to\infty$，并不一定要求分子 $f(x)\to\infty$. 这一点与 $\dfrac{0}{0}$ 型的 L'Hospital 法则不同.

关于 L'Hospital 法则的应用另见例 1.3.12.

\noindent\textbf{$\star$例 3.2.31}\quad 设 $f(x)$ 在 $(-\infty,+\infty)$ 上连续可导，且
\[
  \sup_{-\infty<x<+\infty}|e^{-x^2}f'(x)|<+\infty,
\]
证明：
\[
  \sup_{-\infty<x<+\infty}|xe^{-x^2}f(x)|<+\infty.
\]
（北京大学）

\noindent\textbf{注}\quad 任意函数 $F(x)$，条件 $\sup |F(x)|<+\infty$ 等价于 $F(x)$ 在 $(-\infty,+\infty)$ 上有界. 因此问题等价于已知函数 $e^{-x^2}f'(x)$ 有界，证明 $xe^{-x^2}f(x)$ 有界.

\noindent\textbf{证}\quad 因为 $xe^{-x^2}f(x)$ 在 $[-1,1]$ 上连续，所以在 $[-1,1]$ 上有界，剩下只要证明在 $(1,+\infty)$ 与 $(-\infty,-1)$ 上有界. 以 $(1,+\infty)$ 为例进行证明，$(-\infty,-1)$ 的情况类似. 设 $x>1$ 为任意数，则
\[
\begin{aligned}
  \left|\frac{xf(x)}{e^{x^2}}\right|
  &=
  \left|\frac{xf(x)-f(1)}{e^{x^2}}+\frac{f(1)}{e^{x^2}}\right|  \\
  &\le
  \left|\frac{xf(x)-1\cdot f(1)}{e^{x^2}-e^{1^2}}\right|+\frac{|f(1)|}{e}  \\
  &=
  \left.\left|\frac{(xf(x))'}{(e^{x^2})'}\right|\right|_{x=\xi}+\frac{|f(1)|}{e}
  \quad(1<\xi<x),
\end{aligned}
\tag{1}
\]
其中右端第一项
\[
\begin{aligned}
  \left.\left|\frac{(xf(x))'}{(e^{x^2})'}\right|\right|_{x=\xi}
  &\le \frac12|e^{-\xi^2}f'(\xi)|
  +\frac12\left|e^{-\xi^2}\frac{f(\xi)-f(0)}{\xi-0}\right|
  +\frac12\left|\frac{e^{-\xi^2}}{\xi}f(0)\right| \\
  &\le \frac12|e^{-\xi^2}f'(\xi)|
  +\frac12|e^{-\eta^2}f'(\eta)|+\frac12|f(0)|
  \quad(0<\eta<\xi).
\end{aligned}
\tag{2}
\]
因 $e^{-x^2}f'(x)$ 有界，由 (1)、(2) 知 $xe^{-x^2}f(x)$ 亦有界.

\noindent\textbf{$\star$例 3.2.32}\quad 设 $f(x)$ 在 $\mathbb{R}$ 上可导，试证：若 $\lim\limits_{x\to+\infty}[f(x)+f'(x)]=A$，则 $\lim\limits_{x\to+\infty}f(x)=A,\ \lim\limits_{x\to+\infty}f'(x)=0$.（清华大学）

\noindent\textbf{证}\quad （$e^x$ 的妙用）（可参看例 1.3.21）如上面的注 2° 所述，能用 L'Hospital 法则：
\[
  \lim_{x\to+\infty}f(x)
  =
  \lim_{x\to+\infty}\frac{e^xf(x)}{e^x}
  =
  \lim_{x\to+\infty}(f(x)+f'(x))=A.
\]
因此
\[
  \lim_{x\to+\infty}f'(x)
  =
  \lim_{x\to+\infty}(f(x)+f'(x)-f(x))=A-A=0.
\]

\noindent\textbf{注}\quad 本题说明：如果 $f(x)$ 可导，对 $\lim\limits_{x\to+\infty}\dfrac{e^xf(x)}{e^x}$ 用 L'Hospital 法则，可将 $\lim\limits_{x\to+\infty}f(x)$ 转化为 $\lim\limits_{x\to+\infty}[f(x)+f'(x)]$ 的极限，只要后者极限存在，则
\[
  \lim_{x\to+\infty}f(x)=\lim_{x\to+\infty}[f(x)+f'(x)]=A<+\infty,
\]
且 $\lim\limits_{x\to+\infty}f'(x)=0$；倘若极限 $A$ 不存在，原极限待定.

\noindent\textbf{练习 1}\quad 设 $f(x)$ 在 $\mathbb{R}$ 上有 $n+1$ 阶导数，且 $\lim\limits_{x\to+\infty}f(x)$ 和 $\lim\limits_{x\to+\infty}f^{(n+1)}(x)$ 都存在（有限）（$f^{(0)}(x)\equiv f(x)$），试证：
\[
  \lim_{x\to+\infty}\sum_{k=0}^{n}(-1)^kf^{(k)}(x)
  =
  \lim_{x\to+\infty}[f(x)+(-1)^nf^{(n+1)}(x)],
\]
且
\[
  \lim_{x\to+\infty}\sum_{k=0}^{n}(-1)^kf^{(k+1)}(x)=0.
\]

\noindent\textbf{提示}\quad 记 $F(x)=\sum\limits_{k=0}^{n}(-1)^kf^{(k)}(x)$，对 $\lim\limits_{x\to+\infty}\dfrac{e^xF(x)}{e^x}$ 用 L'Hospital 法则：
\[
\begin{aligned}
  \lim_{x\to+\infty}F(x)
  &=\lim_{x\to+\infty}[F(x)+F'(x)] \\
  &=\lim_{x\to+\infty}
  \sum_{k=0}^{n}(-1)^k[f^{(k)}(x)+f^{(k+1)}(x)] \\
  &=\lim_{x\to+\infty}[f(x)+(-1)^nf^{(n+1)}(x)].
\end{aligned}
\]
且
\[
  \lim_{x\to+\infty}F'(x)
  =
  \lim_{x\to+\infty}\left(\sum_{k=0}^{n}(-1)^kf^{(k)}(x)\right)'
  =
  \lim_{x\to+\infty}\sum_{k=0}^{n}(-1)^kf^{(k+1)}(x)=0.
\]

\noindent\textbf{练习 2}\quad 设 $q>0$，若极限 $\lim\limits_{x\to+\infty}\left[f(x)+\dfrac{1}{q}f'(x)\right]$ 存在，证明 $\lim\limits_{x\to+\infty}f'(x)=0$.

\noindent\textbf{证}\quad
\[
  \lim_{x\to+\infty}f(x)
  =
  \lim_{x\to+\infty}\frac{e^{qx}f(x)}{e^{qx}}
  =
  \lim_{x\to+\infty}\left[f(x)+\frac{1}{q}f'(x)\right]\text{存在},
\]
故知 $\lim\limits_{x\to+\infty}f'(x)=0$.

\noindent\textbf{$\star\star$练习 3}\quad 设 $f(x)$ 在 $\mathbb{R}$ 上有三阶导数，如果 $\lim\limits_{x\to+\infty}f(x)$ 和 $\lim\limits_{x\to+\infty}f'''(x)$ 都存在（有限）. 试证：
\[
  \lim_{x\to+\infty}f'(x)=\lim_{x\to+\infty}f''(x)=\lim_{x\to+\infty}f'''(x)=0.
\]
（中国科学技术大学）

\noindent\textbf{证}\quad 1°
\[
\begin{aligned}
  \lim_{x\to+\infty}[f(x)-f'(x)+f''(x)]
  &=
  \lim_{x\to+\infty}\frac{e^x}{e^x}[f(x)-f'(x)+f''(x)] \\
  &=
  \lim_{x\to+\infty}\{[f(x)-f'(x)+f''(x)]+[f(x)-f'(x)+f''(x)]'\} \\
  &=\lim_{x\to+\infty}[f(x)+f'''(x)].
\end{aligned}
\tag{1}
\]
极限存在. 再利用例 3.2.32，得
\[
  \lim_{x\to+\infty}[f'(x)-f''(x)+f'''(x)]=0.
\tag{2}
\]
故
\[
  \lim_{x\to+\infty}[f'(x)-f''(x)]
  =
  -\lim_{x\to+\infty}f'''(x).
\tag{3}
\]

2° 对 $f'(x)-\dfrac1p f''(x)$ 利用 1° 中的方法，则 $\forall p>0$，
\[
\begin{aligned}
  \lim_{x\to+\infty}\left[f'(x)-\frac1p f''(x)\right]
  &=
  \lim_{x\to+\infty}\frac{e^{qx}}{e^{qx}}\left[f'(x)-\frac1p f''(x)\right] \\
  &=
  \lim_{x\to+\infty}\left\{
  \left[f'(x)-\frac1p f''(x)\right]
  +\frac1q\left[f'(x)-\frac1p f''(x)\right]'\right\} \\
  &=
  \lim_{x\to+\infty}\left[f'(x)+\left(\frac1q-\frac1p\right)f''(x)-\frac1{pq}f'''(x)\right].
\end{aligned}
\]
令 $pq=1,\ p-q=-1$，即 $p,(-q)$ 是方程 $x^2+x-1=0$ 的根：$p=\dfrac{\sqrt5-1}{2},q=\dfrac{\sqrt5+1}{2}$（例 1.1.9 有类似的做法），即
\[
  \lim_{x\to+\infty}\left[f'(x)-\frac1p f''(x)\right]
  =
  \lim_{x\to+\infty}[f'(x)-f''(x)-f'''(x)]
  =
  -2\lim_{x\to+\infty}f'''(x)
\tag{4}
\]
存在.

再利用练习 2，得
\[
  \lim_{x\to+\infty}\left[f'(x)-\frac1p f''(x)\right]'
  =
  \lim_{x\to+\infty}\left[f''(x)-\frac1p f'''(x)\right]=0,
\]
故
\[
  \lim_{x\to+\infty}f''(x)=\lim_{x\to+\infty}\frac1p f'''(x).
\tag{5}
\]
结合式 (3) 知 $\lim\limits_{x\to+\infty}f'(x)$ 也存在. 于是，$\lim\limits_{x\to+\infty}[f'(x)+f''(x)]$ 存在.

再由
\[
  \lim_{x\to+\infty}f'(x)
  =
  \lim_{x\to+\infty}\frac{e^x}{e^x}f'(x)
  =
  \lim_{x\to+\infty}[f'(x)+f''(x)]
\]
存在，利用例 3.2.32 得 $\lim\limits_{x\to+\infty}f''(x)=0$. 由式 (5) 得 $\lim\limits_{x\to+\infty}f'''(x)=0$. 回到式 (3)，知 $\lim\limits_{x\to+\infty}f'(x)=0$. 证毕.

\noindent\textbf{注}\quad 练习 3 采用了学者张相锦和“龙凤呈祥”在网上发布的思想方法，谢谢两位老师.

\noindent\textbf{$\star$例 3.2.33}\quad $f(x)$ 在 $[a,+\infty)$ 上连续可微，$\limsup\limits_{t\to+\infty,\ x>t}|f(x)+f'(x)|\le M<+\infty$. 求证：$\limsup\limits_{t\to+\infty,\ x>t}|f(x)|\le M<+\infty$.（中国科学技术大学）

\noindent\textbf{注}\quad 看起来，本题似乎可用上例的方法同样处理，实际上这里是上极限，不是普通极限，不能像上面那样应用 L'Hospital 法则. 但是，L'Hospital 法则是从 Cauchy 中值定理推得的，本题是否可借助于 Cauchy 中值定理？

\noindent\textbf{证}\quad
\[
  \overline{\lim}_{t\to+\infty}|f(x)+f'(x)|
  =
  \limsup_{\substack{t\to+\infty\\ x>t}}|f(x)+f'(x)|\le M<+\infty.
\]
因此，对 $\varepsilon>0$，$\exists A>0$，当 $x>A$ 时，有
\[
  -\varepsilon-M\le f(x)+f'(x)\le M+\varepsilon.
\]
对函数 $e^xf(x)$ 和 $e^x$ 在 $[A,x]$ 上应用 Cauchy 中值定理：存在 $\xi\in(A,x)$ 使得
\[
  \frac{e^xf(x)-e^Af(A)}{e^x-e^A}
  =
  f(\xi)+f'(\xi)\in(-M-\varepsilon,M+\varepsilon).
\tag{1}
\]
另一方面，利用上、下极限的已知公式（例 1.7.2），有
\[
\begin{aligned}
  \overline{\lim}_{x\to+\infty}\frac{e^xf(x)-e^Af(A)}{e^x-e^A}
  &=
  \overline{\lim}_{x\to+\infty}\frac{e^xf(x)}{e^x-e^A}
  +\lim_{x\to+\infty}\left(-\frac{e^Af(A)}{e^x-e^A}\right) \\
  &=
  \lim_{x\to+\infty}\frac{e^x}{e^x-e^A}\cdot
  \overline{\lim}_{x\to+\infty}f(x)
  =
  \overline{\lim}_{x\to+\infty}f(x).
\end{aligned}
\tag{2}
\]
以上过程将 $f(x)$ 改为 $-f(x)$ 也有效，即
\[
  \overline{\lim}_{x\to+\infty}(-f(x))
  =
  \overline{\lim}_{x\to+\infty}
  \frac{e^x(-f(x))-e^A(-f(A))}{e^x-e^A}.
\]
于是，
\[
  \overline{\lim}_{x\to+\infty}|f(x)|\in[-M-\varepsilon,M+\varepsilon].
\]
令 $\varepsilon\to0$，有
\[
  \limsup_{\substack{t\to+\infty\\ x>t}}|f(x)|
  =
  \overline{\lim}_{x\to+\infty}|f(x)|\le M<+\infty.
\]

\noindent\textbf{$\star$例 3.2.34}\quad 1）设 $F(x)$ 在 $\mathbb{R}=(-\infty,+\infty)$ 上可导，若存在 $|x_n|\to+\infty,\ |y_n|\to-\infty$，使得 $\lim\limits_{n\to\infty}F(x_n)=\lim\limits_{n\to\infty}F(y_n)=c\in\mathbb{R}$. 证明：存在 $\xi\in\mathbb{R}$，使得 $F'(\xi)=0$；

2）设 $f(x),g(x)$ 在 $\mathbb{R}$ 上可导，且 $f'(x)\ne0,\ g'(x)\ne0(\forall x\in\mathbb{R})$. 若 $|x_n|\to+\infty,\ |y_n|\to-\infty,\ |x'_n|\to+\infty,\ |y'_n|\to-\infty$，使得
\[
  \lim_{n\to\infty}f(x_n)=B\in\mathbb{R},\quad
  \lim_{n\to\infty}f(y_n)=A\in\mathbb{R},
\]
\[
  \lim_{n\to\infty}g(x'_n)=b\in\mathbb{R},\quad
  \lim_{n\to\infty}g(y'_n)=a\in\mathbb{R}.
\]
试证：存在 $\xi\in\mathbb{R}$，使得
\[
  \frac{f'(\xi)}{g'(\xi)}=\frac{B-A}{b-a}.
\]
（华东师范大学）

\noindent\textbf{提示}\quad 1）可用反证法（注意导函数有介值性）；

2）作辅助函数 $F(x)=f(x)-kg(x)$，其中 $k=\dfrac{B-A}{b-a}$.

\noindent\textbf{再提示}\quad 1）假设 $\forall x\in\mathbb{R},F'(x)\ne0$. 由导函数的介值性，则 $F'(x)$ 要么恒正要么恒负. 因此，$F$ 严格单调（严增或严减），故 $\lim\limits_{n\to\infty}F(x_n)\ne\lim\limits_{n\to\infty}F(y_n)$，与已知条件矛盾.

2）分析 为证存在 $\xi\in\mathbb{R}$，使得 $\dfrac{f'(\xi)}{g'(\xi)}=\dfrac{B-A}{b-a}\equiv k$，那么 $f'(x)=kg'(x)$，亦即 $f'(x)-kg'(x)=0$. 记 $F(x)=f(x)-kg(x)$，则问题转化为存在 $\xi\in\mathbb{R}$，使得 $F'(\xi)=0$.

已知 $g'(x)\ne0(\forall x\in\mathbb{R})$，因导数有介值性，故 $g'(x)$ 不能变号，即：$g'(x)$ 恒正或恒负，因而 $g(x)$ 严格单调. 当 $x\to+\infty$ 时，
\[
  b=\lim_{n\to\infty}g(x'_n)=\lim_{x\to+\infty}g(x)=\lim_{n\to\infty}g(x_n).
\]
故
\[
  F(x_n)=f(x_n)-kg(x_n)\to B-kb
  =
  B-\frac{B-A}{b-a}b
  =
  \frac{Ab-Ba}{b-a}\equiv c.
\]
当 $y\to-\infty$ 时，同理有
\[
  F(y_n)=f(y_n)-kg(y_n)\to A-ka
  =
  A-\frac{B-A}{b-a}a
  =
  \frac{Ab-Ba}{b-a}=c.
\]
可见，$F(x)=f(x)-kg(x)$ 满足第 1）小题的条件. 因此，存在 $\xi\in\mathbb{R}$，使得 $F'(\xi)=0$，亦即 $\dfrac{f'(\xi)}{g'(\xi)}=\dfrac{B-A}{b-a}$. 结论 2）成立.

\noindent\textbf{另证}\quad 如上所述，因为导数有介值性，已知 $g'(x)\ne0$，则 $g'(x)$ 不变号，$g(x)$ 严格单调，$\lim\limits_{x\to+\infty}g(x)$ 以有限数或无穷大作为极限. 但已知存在 $|x'_n|\to+\infty,\lim\limits_{n\to\infty}g(x'_n)=b$，因此，$\lim\limits_{x\to+\infty}g(x)=\lim\limits_{n\to\infty}g(x'_n)=b$. 同理，因为 $f'(x)\ne0$，故对 $f(x)$ 也有类似的结论：

至此，在 $\left(-\dfrac{\pi}{2},\dfrac{\pi}{2}\right)$ 内对 $f(\tan x)$ 和 $g(\tan x)$ 可应用 Cauchy 中值定理：因为，若
\[
  f(\tan x)\bigg|_{x=\frac{\pi}{2}}
  =
  \lim_{x\to\frac{\pi}{2}}f(\tan x)=B
  \quad\text{和}\quad
  f(\tan x)\bigg|_{x=-\frac{\pi}{2}}
  =
  \lim_{x\to-\frac{\pi}{2}}f(\tan x)=A,
\]
\[
  g(\tan x)\bigg|_{x=\frac{\pi}{2}}
  =
  \lim_{x\to\frac{\pi}{2}}g(\tan x)=b
  \quad\text{和}\quad
  g(\tan x)\bigg|_{x=-\frac{\pi}{2}}
  =
  \lim_{x\to-\frac{\pi}{2}}g(\tan x)=a,
\]
则 $f(\tan x)$ 和 $g(\tan x)$ 在 $\left[-\dfrac{\pi}{2},\dfrac{\pi}{2}\right]$ 上连续，在 $\left(-\dfrac{\pi}{2},\dfrac{\pi}{2}\right)$ 内可导，并能应用 Cauchy 中值定理.

如此，可直接得到：存在 $\xi=\tan\eta\in\mathbb{R}$，使得
\[
  \frac{f'(\xi)}{g'(\xi)}
  =
  \frac{f'(\tan\eta)}{g'(\tan\eta)}
  =
  \frac{
    f(\tan x)\big|_{x=\frac{\pi}{2}}^{x=-\frac{\pi}{2}}
  }{
    g(\tan x)\big|_{x=\frac{\pi}{2}}^{x=-\frac{\pi}{2}}
  }
  =
  \frac{B-A}{b-a}.
\]

\noindent\textbf{※附：导数的推广——广义导数}

\noindent\textbf{要点}\quad 定义 1\quad $\Delta_h f(x)=f(x+h)-f(x-h)$ 称为 $f(x)$ 在点 $x$ 处的对称差，如下极限（若存在）称为 $f(x)$ 在点 $x$ 处的广义导数或对称导数，记作
\[
  f^{[1]}(x)
  =
  \lim_{h\to0+}\frac{\Delta_h f(x)}{2h}
  =
  \lim_{h\to0+}\frac{f(x+h)-f(x-h)}{2h}.
\tag{1}
\]

\noindent\textbf{推论 1}\quad 若 $f(x)$ 在点 $x$ 处可导，则 $f(x)$ 在点 $x$ 处有广义导数，且它们的值相等：
\[
  f^{[1]}(x)
  =
  \lim_{h\to0+}\frac12\left[
  \frac{f(x+h)-f(x)}{h}
  +\frac{f(x-h)-f(x)}{-h}\right]
  =
  f'(x).
\]
逆定理不成立，因为：$f^{[1]}(x)$ 只与 $f$ 在对称差有关，与中心点的值无关. 广义可导未必真可导，例如，
\[
  f(x)=
  \begin{cases}
    x\sin\dfrac1x,& x\ne0,\\
    0,& x=0,
  \end{cases}
  \quad\text{或}\quad f(x)=|x|,
\]
在 $x=0$ 处广义可导，$f^{[1]}(0)=0$，但 $f$ 在原点不可导.

\noindent\textbf{注}\quad 若 $f$ 只在 $[a,b]$ 上有定义，那么 $f$ 在 $[a,b]$ 的端点处无“对称导数”可言. 另外，$f^{[1]}(x_0)$ 存在不能推出 $f$ 在 $x_0$ 处连续（这些都跟导数不一样）.

\noindent\textbf{定义 2}\quad 如下的极限存在，称为 $f(x)$ 在点 $x$ 处有二阶广义导数（或称二阶对称导数），记作
\begin{equation}
  f^{[2]}(x)
  =
  \lim_{h\to0+}\frac{\Delta_h\Delta_h f(x)}{(2h)^2}
  =
  \lim_{h\to0+}\frac{\Delta_h[f(x+h)-f(x-h)]}{(2h)^2}
  =
  \lim_{h\to0+}\frac{f(x+2h)-2f(x)+f(x-2h)}{4h^2}.
\tag{2}
\end{equation}
如此递推下去，可以定义任意 $n$ 阶广义导数 $f^{[n]}(x)\ (n=2,3,\cdots)$.

\noindent\textbf{推论 2}\quad 在某点 $x$，若 $f(x)$ 有二阶导数，则二阶广义导数也存在，且值相等（见练习 3.2.32 题）：$f''(x)=f^{[2]}$.

就线性函数而言，有
\[
  (\alpha x+\beta)^{[1]}=(\alpha x+\beta)'=\alpha,
  \qquad
  (\alpha x^2+\beta x+\gamma)^{[2]}=(\alpha x^2+\beta x+\gamma)''=2\alpha.
\]
同样，逆定理不成立. 例如，
\[
  f(x)=
  \begin{cases}
    x^2\sin\dfrac1x,& x\ne0,\\
    0,& x=0
  \end{cases}
  \quad\text{或 }f(x)=x|x|,
\]
此时 $f(x)$ 在 $x=0$ 处无二阶导数，但 $f^{[2]}(0)$ 存在，$f^{[2]}(0)=0$.

类似地，若 $n$ 阶导数存在，则 $n$ 阶广义导数必然存在，且 $f^{[n]}(x)=f^{(n)}(x)$.

\noindent\textbf{推论 3}\quad 从式 (2) 明显看出：设 $f$ 在内点 $x_0$ 连续，并在点 $x_0$ 取极大值，则当 $f^{[2]}(x_0)$ 存在时，有 $f^{[2]}(x_0)\le0$.

若 $f$ 在某内点 $x_1$ 连续，且在点 $x_1$ 取极小值，则当 $f^{[2]}(x_1)$ 存在时，有 $f^{[2]}(x_1)\ge0$.

\subsection*{单元练习 3.2}

\noindent\textbf{关于函数零值点（方程根）的存在唯一性}

\noindent\textbf{$\star$3.2.1}\quad 若 $f(x)$ 在 $[a,b]$ 上连续，且 $f(a)=f(b)=0,\ f'(a)\cdot f'(b)>0$，证明：存在 $\xi\in(a,b)$，使 $f(\xi)=0$.（哈尔滨工业大学，华中科技大学，华中师范大学）

\noindent\textbf{提示}\quad $f'(x)$ 在 $a,b$ 同号 $\Rightarrow$ 在 $a$ 右侧 $\exists x_1$，$b$ 左侧 $\exists x_2$，使得 $f(x_1),f(x_2)$ 异号.

\noindent\textbf{再提示}\quad 例如 $f'(a)>0$，即 $\lim\limits_{x\to a+}\dfrac{f(x)-f(a)}{x-a}>0$，则 $x_1>a$ 且与 $a$ 充分接近时有 $f(x_1)>f(a)=0$，这时 $f'(b)>0$. 类似地，$\exists x_2<b$（与 $b$ 充分接近）使得 $f(x_2)<f(b)=0$，从而由连续介值性可得欲证之结论.

\noindent\textbf{$\star$3.2.2}\quad 设 $a,b,c$ 为三个实数，证明：方程 $e^x=ax^2+bx+c$ 的根不超过三个.（浙江大学，武汉理工大学）

\noindent\textbf{提示}\quad 可用反证法.

\noindent\textbf{再提示}\quad 否则 $F(x)=e^x-ax^2-bx-c$ 的零点个数 $\ge4$，反复使用 Rolle 定理，可知 $F''(x)=e^x-2a$ 应有一个零点，矛盾.

\noindent\textbf{3.2.3}\quad 设 $f(x)$ 与 $g(x)$ 在 $(a,b)$ 内可微，$f(x)g'(x)\ne f'(x)g(x)$. 证明：$f(x)=0$ 的两个根之间至少夹 $g(x)=0$ 的一根.（上海交通大学）

\noindent\textbf{提示}\quad 考虑辅助函数 $\dfrac{f(x)}{g(x)}$.

\noindent\textbf{再提示}\quad 若 $f(x)=0$ 的某两个根 $x_1,x_2$ 之间无 $g$ 之零点，对函数 $\dfrac{f(x)}{g(x)}$ 在 $[x_1,x_2]$ 上应用 Rolle 定理，可得出 $\exists \xi\in(x_1,x_2)$，使 $f'(\xi)g(\xi)=f(\xi)g'(\xi)$（矛盾）.

\noindent\textbf{3.2.4}\quad 设 $a^2-3b<0$，试证：$x^3+ax^2+bx+c=0$ 仅有唯一实根.

\noindent\textbf{提示}\quad 一方面：虚根，必成对出现，故三次方程至少有 1 个实根.

\noindent\textbf{再提示}\quad 另一方面：题设 $a^2-3b<0$，若
\[
  f'(x)\overset{\text{记}}{=}(x^3+ax^2+bx+c)'=3x^2+2ax+b=0,
\]
其判别式 $\Delta=\dfrac49(a^2-3b)\overset{\text{按题设}}{<}0$，说明 $f'(x)\ne0$（无零点），故 $f'(x)$ 不可能变号（否则由导数的介值性 $f'(x)$ 有零点）. 因此 $f(x)$ 单调，从而原方程 $f(x)=0$ 最多有一（实）根. 而上面说：至少有 1 个实根，故有且仅有唯一实根.

\noindent\textbf{*3.2.5}\quad 设 $f(x)=\dfrac{\mathrm d^n}{\mathrm dx^n}(e^{-x}x^n),\ x\in(0,+\infty)$，证明：函数 $f(x)$ 在 $(0,+\infty)$ 中恰有 $n$ 个零点.（清华大学）

\noindent\textbf{提示}\quad 反复使用 Rolle 定理及例 3.2.1 之结果.

\noindent\textbf{证}\quad $n=1,2$ 的情况明显，只需证明高阶的情况. $F_1(x)=(e^{-x}x^n)'=e^{-x}(-x^n+nx^{n-1})$，$0,+\infty$ 为其零点；因而 $F_2(x)=(e^{-x}x^n)''=e^{-x}[x^n-2nx^{n-1}+n(n-1)x^{n-2}]$ 包括 0 和 $+\infty$，共有三个零点，$0<c<+\infty$；如此递推下去，设 $F_{k-1}(x)\equiv(e^{-x}x^n)^{(k-1)}$ 有 $k$ 个零点，$0<c_1<\cdots<c_{k-2}<+\infty$. 下面来证：当 $k\le n-1$ 时，$F_k(x)$ 必有 $k+1$ 个不同的零点，$0<d_1<d_2<\cdots<d_{k-1}<+\infty$. 事实上，利用 Leibniz 求导公式，
\[
\begin{aligned}
F_k(x)
&=(e^{-x}x^n)^{(k)}
=\sum_{i=0}^k C_k^i(e^{-x})^{(k-i)}(x^n)^{(i)}
=\sum_{i=0}^k(-1)^{k-i}C_k^i e^{-x}\frac{n!}{(n-i)!}x^{n-i}\\
&=e^{-x}\bigl((-1)^k x^n+(-1)^{k-1}knx^{n-1}+\cdots+n(n-1)\cdots(n-k+1)x^{n-k}\bigr).
\end{aligned}
\]
（注意 $k\le n-1$ 时，$n-k\ge1$）. 可见 $F_k(+\infty)=0$，当 $n-k\ge1$（即 $k\le n-1$）时 $F_k(0)=0$. 根据 Rolle 定理及其推广形式（例 3.2.1），每两个 $F_{k-1}(x)$ 的零点之间，有 $F_k(x)$ 的一个零点，记作 $d_i$，故 $F_k(x)$ 有 $0<d_1<\cdots<d_{k-1}<+\infty$ 共 $k+1$ 个零点. 因此当 $k=n-1$ 时，其有 $n$ 个不同零点（包括 0 和 $+\infty$）. 再用一次 Rolle 定理及例 3.2.1，可知
\[
  F_n(x)=e^{-x}\bigl[(-1)^n x^n+(-1)^{n-1}nC_n^1x^{n-1}+\cdots+n!\bigr]
\]
在 $(0,+\infty)$ 中恰有 $n$ 个不同的零点.

\noindent\textbf{*3.2.6}\quad 证明 Chebyshev-Laguerre 多项式
\[
  L_n(x)=e^x\frac{\mathrm d^n}{\mathrm dx^n}(e^{-x}x^n)
\]
所有的根都为正.

\noindent\textbf{提示}\quad 利用上题.

\noindent\textbf{*3.2.7}\quad 试证：Chebyshev-Hermite 多项式
\[
  H_n(x)=(-1)^n e^{x^2}\frac{\mathrm d^n}{\mathrm dx^n}(e^{-x^2})
\]
的所有根都是实数.

\noindent\textbf{提示}\quad 类似上题.

\noindent\textbf{3.2.8}\quad 证明：当 $\dfrac{a_0}{n+1}+\dfrac{a_1}{n}+\cdots+\dfrac{a_{n-1}}2+a_n=0$ 时，方程 $a_0x^n+a_1x^{n-1}+\cdots+a_n=0$ 在 $(0,1)$ 内至少有一实根.（南京邮电大学）

\noindent\textbf{提示}\quad 对函数 $F(x)=\int_0^x(a_0t^n+a_1t^{n-1}+\cdots+a_n)\,\mathrm dt$ 在 $[0,1]$ 上应用 Rolle 定理.

\noindent\textbf{$\star$3.2.9}\quad 设函数 $f(x)$ 在 $[a,+\infty)$ 上连续，且当 $x>a$ 时，$f'(x)>k>0$（$k$ 为常数），证明：当 $f(a)<0$ 时，方程 $f(x)=0$ 在区间 $\left(a,a-\dfrac{f(a)}k\right)$ 内有且只有一个根.（湘潭大学，西安交通大学，西安电子科技大学等）

\noindent\textbf{提示}\quad 在 $\left[a,a-\dfrac{f(a)}k\right]$ 上应用 Lagrange 公式，知 $f\left(a-\dfrac{f(a)}k\right)$ 与 $f(a)$ 异号.

\noindent\textbf{$\star$3.2.10}\quad 设 $f(x)$ 在 $(-\infty,+\infty)$ 内具有二阶导数，且 $f''(x)>0,\ \lim\limits_{x\to+\infty}f'(x)=\alpha>0,\ \lim\limits_{x\to-\infty}f'(x)=\beta<0$，又存在一点 $x_0$，使 $f(x_0)<0$. 试证明：方程 $f(x)=0$ 在 $(-\infty,+\infty)$ 内有且只有两个实根.（上海交通大学，浙江大学）

\noindent\textbf{提示}\quad 可证 $f(-\infty)=f(+\infty)=+\infty$，从而由 $f(x_0)<0$ 知有两个根. 又由 $f''(x)>0$，可证只有两个实根.

\noindent\textbf{再提示}\quad 由 $f'(+\infty)>0$ 可知 $\exists b>0$，使得 $f'(b)>0$. 又 $x>b$ 时 $\exists \xi:x>\xi>b$，使得
\[
  f(x)=f(b)+f'(\xi)(x-b),
  \quad\text{但 }f''(x)>0,\ f'(x)\nearrow,
\]
\[
  f(x)\ge f(b)+f'(b)(x-b)\to+\infty\quad(\text{当 }x\to+\infty\text{ 时}).
\]
因此，$f(+\infty)=+\infty$；同样由 $f'(-\infty)<0$ 亦可推出 $f(-\infty)=+\infty$.

此外，由 $f''>0$ 可得 $f'$ 严格增，又因 $f'(-\infty)<0,\ f'(+\infty)>0$，故存在唯一一点 $c$，使得
\[
  f'(x)
  \begin{cases}
  <0,& x<c,\\
  =0,& x=c,\\
  >0,& x>c.
  \end{cases}
\]
因此 $c$ 是 $f$ 的最小值点，$f(c)\le f(x_0)<0$，$f$ 在 $(-\infty,c]$ 和 $[c,+\infty)$ 分别由 $+\infty\searrow$ 负，又由负 $\nearrow+\infty$，故 $f(x)=0$ 有且仅有两个实根.

\noindent\textbf{推导新的中值形式}

\noindent\textbf{$\star$3.2.11}\quad 设 $f(x)$ 在 $[0,+\infty)$ 内可微，且满足不等式
\[
  0\le f(x)\le \ln\frac{2x+1}{x+\sqrt{1+x^2}},\quad \forall x\in(0,+\infty),
\]
试证明存在一点 $\xi\in(0,+\infty)$，使得 $f'(\xi)=\dfrac2{2\xi+1}-\dfrac1{\sqrt{1+\xi^2}}$.（北京大学）

\noindent\textbf{提示}\quad 参考例 3.2.8.

\noindent\textbf{$\star$3.2.12}\quad 设函数 $f(x)$ 在 $[a,b]$ 上连续，在 $(a,b)$ 内可微，且 $f(a)<0,\ f(b)<0$，又有一点 $c\in(a,b)$，$f(c)>0$. 证明：存在一点 $\xi\in(a,b)$，使得 $f(\xi)+f'(\xi)=0$.（西北大学）

\noindent\textbf{提示}\quad （函数 $e^x$ 的妙用）考虑辅助函数 $F(x)=e^xf(x)$，分别在 $[a,c]$ 和 $[c,b]$ 上两端点异号；再用连续函数介值性及 Rolle 定理.

\noindent\textbf{3.2.13}\quad 设 $f(x)$ 在 $[0,1]$ 上连续，在 $(0,1)$ 内可微，证明：$\exists \xi\in(0,1)$，使得 $f'(\xi)f(1-\xi)=f(\xi)f'(1-\xi)$.（华中科技大学）

\noindent\textbf{提示}\quad 考虑辅助函数 $F(x)=f(x)f(1-x)$.

\noindent\textbf{3.2.14}\quad 假设函数 $f(x)$ 和 $g(x)$ 在 $[a,b]$ 上存在二阶导数，并且 $g''(x)\ne0,\ f(a)=f(b)=g(a)=g(b)=0$，试证

1）在开区间 $(a,b)$ 内 $g(x)\ne0$；

2）在 $(a,b)$ 内至少存在一点 $\xi$，使得 $\dfrac{f(\xi)}{g(\xi)}=\dfrac{f''(\xi)}{g''(\xi)}$.（数学一）

\noindent\textbf{提示}\quad 1）可用反证法，用 Rolle 定理推出 $g''$ 有零点，得出矛盾；2）借助辅助函数 $F(x)=f(x)g' (x)-f'(x)g(x)$.

\noindent\textbf{$\star$3.2.15}\quad 设 $f(x)$ 在 $[a,b]$ 上非负且三阶可导，方程 $f(x)=0$ 在 $(a,b)$ 内有两个不同实根，证明存在 $\xi\in(a,b)$，使得 $f^{(3)}(\xi)=0$.（华中师范大学）

\noindent\textbf{提示}\quad $f(x)=0$ 的两根之间有 $f'$ 的零点. 因 $f\ge0$，故 $f$ 的零点也是 $f$ 的极小点. 根据 Fermat 定理，其上导数也应为零. 因而 $f'$ 有三个零点.

\noindent\textbf{new 3.2.16}\quad 设 $f(x)$ 在 $[0,1]$ 上可导，$f(1)=2f(0)$，求证：$\exists \xi\in(0,1)$，使得 $(1+\xi)f'(\xi)=f(\xi)$.（武汉大学）

\noindent\textbf{提示}\quad $(1+\xi)f'(\xi)=f(\xi)\Leftrightarrow \left.\left(\dfrac{f(x)}{1+x}\right)'\right|_{x=\xi}=0$.

\noindent\textbf{3.2.17}\quad 设 $f(x)$ 在 $[a,b]$ 上二阶可导，过点 $A(a,f(a))$ 与 $B(b,f(b))$ 的直线与曲线 $y=f(x)$ 相交于 $C(c,f(c))$，其中 $a<c<b$. 证明：在 $(a,b)$ 中至少存在一点 $\xi$，使得 $f''(\xi)=0$.（华中师范大学）

\noindent\textbf{提示}\quad $A,B,C$ 三点共线，在 $[a,c],[c,b]$ 上分别应用 Lagrange 公式，可知 $\exists x_1,x_2:a<x_1<c<x_2<b$ 使得 $f'(x_1)=f'(x_2)$（弦 $AC,CB$ 的斜率）.

\noindent\textbf{$\star$3.2.18}\quad 设函数 $f(x)$ 在闭区间 $[a,b]$ 上连续，在开区间 $(a,b)$ 内二阶可微，并且 $f(a)=f(b)$，证明：若存在一点 $c\in(a,b)$，使得 $f(c)>f(a)$，则必存在三点 $\xi,\eta,\zeta\in(a,b)$，使得 $f'(\xi)>0,\ f'(\eta)<0,\ f''(\zeta)<0$.（吉林大学，北京师范大学，国防科技大学）

\noindent\textbf{提示}\quad 因 $f(c)>f(a)=f(b)$，曲线 $y=f(x)$ 在 $[a,c],[c,b]$ 上弦的斜率异号，可在两段上分别应用 Lagrange 公式.

\noindent\textbf{3.2.19}\quad 函数 $f(x)$ 在 $[0,x]$ 上的拉格朗日中值公式为
\[
  f(x)-f(0)=f'(\theta x)x,\quad \text{其中 }0<\theta<1,
\]
且 $\theta$ 是与 $f(x)$ 及 $x$ 有关的量. 对 $f(x)=\arctan x$，求 $x\to0^+$ 时 $\theta$ 的极限值.（武汉大学）

\noindent\textbf{提示}\quad $\theta^2=\dfrac{x-\arctan x}{x^2\arctan x}$.

\noindent\textbf{$\star$3.2.20}\quad 设 $f(x)$ 在 $[1,2]$ 上连续，在 $(1,2)$ 内可微，证明：存在 $\xi\in(1,2)$ 使得 $f(2)-f(1)=\dfrac12\xi^2 f'(\xi)$.（北京科技大学）

\noindent\textbf{提示}\quad $\dfrac{f(2)-f(1)}{\frac12-\frac11}=\dfrac{f'(\xi)}{-\frac1{\xi^2}}$（参看例 3.2.26）.

\noindent\textbf{3.2.21}\quad 设 $f(x)$ 在 $[a,b]$ 上有三阶导数. 试证：必存在点 $\xi\in(a,b)$，使得
\[
  f(b)=f(a)+\frac12(b-a)[f'(a)+f'(b)]-\frac1{12}(b-a)^3 f'''(\xi).
\tag{1}
\]
（郑州大学）

\noindent\textbf{提示}\quad 参看例 3.2.27，改写式 (1)，用 Cauchy 中值定理.

\noindent\textbf{再提示}\quad 第一次对
\[
  F(x)=f(x)-f(a)-\frac12(x-a)[f'(x)+f'(a)]
  \quad\text{和}\quad
  G(x)=-\frac1{12}(x-a)^3
\]
在 $[a,x]$ 上应用 Cauchy 中值公式.
\[
\begin{aligned}
\frac{F(x)}{G(x)}
&=\frac{F(x)-F(a)}{G(x)-G(a)}
=\frac{F'(\eta)}{G'(\eta)}
=\frac{f'(\eta)-\frac12[f'(\eta)+f'(a)]-\frac12(\eta-a)f''(\eta)}
{-\frac14(\eta-a)^2}\\
&=\frac{F'(\eta)-F'(a)}{G'(\eta)-G'(a)}
=\frac{F''(\xi)}{G''(\xi)}
=f'''(\xi).
\end{aligned}
\]
（利用 $F(a)=F'(a)=G(a)=G'(a)=0$.）

\noindent\textbf{3.2.22}\quad $f(x)$ 在 $[a,b]$ 上连续，$f''(x)$ 在 $(a,b)$ 内存在，试证：$\forall c:a<c<b,\ \exists \xi\in(a,b)$，使得
\[
  \frac12f''(\xi)=
  \frac{f(a)}{(a-b)(a-c)}
  +\frac{f(b)}{(b-c)(b-a)}
  +\frac{f(c)}{(c-a)(c-b)}.
\tag{1}
\]

\noindent\textbf{提示}\quad 将式 (1) 改写成
\[
  \frac12f''(\xi)=
  \frac{f(a)(c-b)+f(b)(a-c)+f(c)(b-a)}
  {(a-b)(b-c)(c-a)},
\]
然后类似上题，反复利用 Cauchy 中值定理.

\noindent\textbf{$\star$3.2.23}\quad 设 $f(x)$ 在 $[a,b]$ 上连续，在 $(a,b)$ 内可微，$b>a>0$，证明：在 $(a,b)$ 内存在 $x_1,x_2,x_3$，使得
\[
  \frac{f'(x_1)}{2x_1}
  =(b^2+a^2)\frac{f'(x_2)}{4x_2^3}
  =\frac{\ln(b/a)}{b^2-a^2}x_3\cdot f'(x_3).
\tag{1}
\]
（四川大学）

\noindent\textbf{提示}\quad 参看例 3.2.28, 3.2.29.

\noindent\textbf{再提示}\quad 将式 (1) 改写成
\[
  f(b)-f(a)
  =
  \frac{f'(x_1)}{2x_1}(b^2-a^2)
  =
  \frac{f'(x_2)}{4x_2^3}(b^4-a^4)
  =
  \frac{f'(x_3)}{1/x_3}(\ln b-\ln a),
\]
在 $[a,b]$ 上，分别对 $f(x)$ 与 $x^2$；$f(x)$ 与 $x^4$；$f(x)$ 与 $\ln x$ 应用 Cauchy 中值定理，以证明上式后三项分别等于第一项 $f(b)-f(a)$.

\noindent\textbf{3.2.24}\quad 设 $f(x)$ 在区间 $[a,b]$ 上连续，在 $(a,b)$ 内可导，且 $a\ge0$（或 $b\le0$）. 试证：

1）$\exists x_1,x_2,x_3\in(a,b)$，使得
\[
  f'(x_1)=(b+a)\frac{f'(x_2)}{2x_2}
  =(b^2+ba+a^2)\frac{f'(x_3)}{3x_3^2};
\]
（南京航空航天大学）

2）$\forall n\in\{1,2,\cdots\}$，$\exists x_1,x_2,\cdots,x_n\in(a,b)$，使得
\[
  f'(x_1)=(b+a)\frac{f'(x_2)}{2x_2}
  =\cdots
  =(b^{n-1}+b^{n-2}a+\cdots+ba^{n-2}+a^{n-1})\frac{f'(x_n)}{nx_n^{n-1}}.
\tag{1}
\]

\noindent\textbf{提示}\quad 将式 (1) 诸量同时乘以 $(b-a)$，改写成
\[
  \frac{f'(x_1)}1(b-a)
  =\frac{f'(x_2)}{2x_2}(b^2-a^2)
  =\cdots
  =\frac{f'(x_n)}{nx_n^{n-1}}(b^n-a^n).
\]

\noindent\textbf{小结}\quad 以上 5 题主要练习利用 Cauchy 中值定理推导新的中值公式.

\noindent\textbf{微分中值定理的灵活应用}

\noindent\textbf{3.2.25}\quad 设 $f(x)$ 在有限区间 $(a,b)$ 内可微，试证：

1）若 $f'(x)$ 在 $(a,b)$ 内有界，则 $f(x)$ 在 $(a,b)$ 内亦有界；（北京师范大学）

2）若 $f(x)$ 在 $(a,b)$ 内无界，则 $f'(x)$ 在 $(a,b)$ 内亦无界.（华东师范大学）

\noindent\textbf{提示}\quad 1）$\exists M>0,\ |f'(x)|\le M\Rightarrow |f(x)|\le |f(a+0)|+M(b-a)$.

\noindent\textbf{*3.2.26}\quad 设 $f$ 在 $[a,b]$ 中任意两点都具有介值性：$c_1,c_2\in[a,b]$，$\forall r:f(c_1)<r<f(c_2)$，$\exists c$（在 $c_1,c_2$ 之间），使得 $f(c)=r$；而且 $f$ 在 $(a,b)$ 内可导，$|f'(x)|\le k$（正常数），$\forall x\in(a,b)$. 试证：$f$ 在 $a$ 处右连续（同理在 $b$ 处左连续）.（华东师范大学）

\noindent\textbf{注}\quad “用 $|f(x)-f(a)|=|f'(\xi)||x-a|$”是不对的！

\noindent\textbf{证}\quad $\forall \varepsilon>0$，取 $\delta=\dfrac{\varepsilon}{2k}$，则 $\forall x\in(a,a+\delta)$ 有 $|f(x)-f(a)|<\varepsilon$. 事实上，若 $|f(x)-f(a)|<\dfrac{\varepsilon}{2}$，问题自明；否则，由介值性条件，必 $\exists x'\in(a,x)$ 使得 $|f(x')-f(a)|<\dfrac{\varepsilon}{2}$.
\[
  |f(x)-f(a)|
  \le |f(x)-f(x')|+|f(x')-f(a)|
  \le |f'(\xi)(x-x')|+\frac{\varepsilon}{2}
  < k\cdot\frac{\varepsilon}{2k}+\frac{\varepsilon}{2}
  =\varepsilon.
\]

\noindent\textbf{*3.2.27}\quad 设 $f(x)$ 是 $(-\infty,+\infty)$ 内的可微函数.

1）若 $\lim\limits_{x\to+\infty} f(x)$ 存在且有限，问 $\lim\limits_{x\to+\infty} f'(x)$ 是否必定存在？（云南大学）

2）如果 $\lim\limits_{x\to+\infty} f(x)$ 与 $\lim\limits_{x\to+\infty} f'(x)$ 都存在且有限，那么必有 $\lim\limits_{x\to+\infty} f'(x)=0$，试证明之.（云南大学，哈尔滨工业大学等）

\noindent\textbf{提示}\quad 1）否. 例如
\[
  f(x)=
  \begin{cases}
    \dfrac{\sin x^2}{x},& x\ne0,\\
    0,& x=0.
  \end{cases}
\]
2）利用 Cauchy 准则及中值定理.

\noindent\textbf{证}\quad $\forall\varepsilon>0$，$f(+\infty)$ 存在 $\Rightarrow\exists\Delta_1>0$，当 $x',x''>\Delta_1$ 时，$|f(x')-f(x'')|<\dfrac{\varepsilon}{2}$；

$f'(+\infty)$ 存在 $\Rightarrow\exists\Delta_2>0$，当 $x',x''>\Delta_2$ 时，$|f'(x')-f'(x'')|<\dfrac{\varepsilon}{2}$.

令 $\Delta=\max\{\Delta_1,\Delta_2\}$，则 $x>\Delta$ 时，$\exists \xi\in(x,x+1)$ 使得
\[
  |f'(\xi)|=|f(x+1)-f(x)|<\frac{\varepsilon}{2},\quad \xi>x>\Delta\ge\Delta_2,
\]
故
\[
  |f'(x)|\le |f'(x)-f'(\xi)|+|f'(\xi)|<\frac{\varepsilon}{2}+\frac{\varepsilon}{2}=\varepsilon.
\]

\noindent\textbf{3.2.28}\quad 设 $f(x)$ 于 $(0,1)$ 内可微，且满足 $|f'(x)|\le1$，求证：$\lim\limits_{n\to\infty}f\left(\dfrac1n\right)$ 存在.（哈尔滨工业大学）

\noindent\textbf{提示}\quad
\[
  \left|f\left(\frac1{n+p}\right)-f\left(\frac1n\right)\right|
  =|f'(\xi_n)|\frac{p}{n(n+p)}
  \le\frac1n\to0\quad(\text{关于 }p\text{ 一致}).
\]

\noindent\textbf{$\star$3.2.29}\quad 设 $f(x)$ 在 $[0,+\infty)$ 上连续，在 $(0,+\infty)$ 内可微，$f(0)=0$，试证

1）若 $f'(x)$ 单调增加，则 $g(x)=\dfrac{f(x)}x$ 在 $(0,+\infty)$ 内单调增加；（同济大学，武汉大学，四川大学等）

2）若 $f'(x)$ 单调递减，则 $\dfrac{f(x)}x$ 在 $(0,+\infty)$ 内单调递减.（中国科学院）

\noindent\textbf{提示}\quad 1）$g'(x)=\dfrac1x\left[f'(x)-\dfrac{f(x)-f(0)}x\right]=\dfrac1x[f'(x)-f'(\xi)]>0$.

\noindent\textbf{3.2.30}\quad 设 $f(x)$ 在 $[a,b]$ 上连续，且 $(a,b)$ 内可微. 证明：若存在极限 $\lim\limits_{x\to a+}f'(x)=l$，则右导数 $f'_+(a)$ 存在且等于 $l$.（北京大学，湖北大学）

\noindent\textbf{提示}\quad $f'_+(a)=\lim\limits_{x\to a+}\dfrac{f(x)-f(a)}{x-a}=\lim\limits_{x\to a+}f'(\xi)=f'(a+0)=l$.

由此可见 $f$ 在 $[a,b]$ 上连续，在 $(a,b)$ 内可微，只要 $f'(a+0),f'(b-0)$ 存在，则 $f$ 在 $[a,b]$ 上可微.

\noindent\textbf{$\star$3.2.31}\quad 设
\[
  f(x)=
  \begin{cases}
    |x|,& x\ne0,\\
    1,& x=0,
  \end{cases}
\]
证明：不存在一个函数以 $f$ 为其导函数.（中国科学院）

\noindent\textbf{提示}\quad 若 $\exists g$ 使 $g'=f$，则 $1=f(0)=g'(0)=g'_+(0)=g'(0+0)=0$，矛盾.

\noindent\textbf{3.2.32}\quad 证明：若 $f''(0)$ 存在（有限），则 $\lim\limits_{h\to0}\dfrac{f(2h)-2f(0)+f(-2h)}{4h^2}=f''(0)$.（北京师范大学）

\noindent\textbf{提示}\quad 可对 $F(x)=f(2x)-2f(0)+f(-2x),\ G(x)=4x^2$ 在 $[0,x]$ 上应用 Cauchy 微分中值定理（一次）. 注意，$f''(0)$ 存在意味着在 $x=0$ 的邻域里 $f'(x)$ 存在. 但题目中未假定在 $x=0$ 的邻域里有二阶导数，因此用了一次 Cauchy 微分中值定理之后，不能用第二次. 剩下的问题可用导数定义解决.

\noindent\textbf{※3.2.33}\quad 将上题结果推广到一般情况，即若 $f^{(n)}(0)$ 存在（有限），则（$n$ 是自然数）
\[
  \lim_{h\to0+}\sum_{k=0}^n C_n^k\frac{(-1)^k f[(n-2k)h]}{(2h)^n}=f^{(n)}(0).
\]
（北京师范大学）

\noindent\textbf{提示}\quad （用数学归纳法.）注意：$C_n^0=1,\ C_n^k+C_n^{k+1}=C_{n+1}^{k+1}\ (k=0,1,\cdots,n-1),\ C_{n+1}^{n+1}=1$.

\noindent\textbf{证}\quad $1^\circ$ 当 $n=1,2$ 时，原式成立. 因为当 $n=1$ 时，
\[
\begin{aligned}
\left.\lim_{h\to0+}\frac1{(2h)^n}\sum_{k=0}^n(-1)^kC_n^k[f((n-2k)h)]\right|_{n=1}
&=\lim_{h\to0+}\frac{f(h)-f(-h)}{2h}\\
&=\lim_{h\to0+}\frac{f(h)-f(0)+[f(0)-f(-h)]}{2h}
=f'(0).
\end{aligned}
\]
当 $n=2$ 时，
\[
\begin{aligned}
\left.\lim_{h\to0+}\frac1{(2h)^n}\sum_{k=0}^n(-1)^kC_n^k[f((n-2k)h)]\right|_{n=2}
&=\lim_{h\to0+}\frac{f(2h)-2f(0)+f(-2h)}{(2h)^2}\\
&=\lim_{h\to0+}\frac1{2h}\frac{\Delta_h}{2h}[f(h)-f(-h)]
=f^{[2]}(0)=f''(0).
\end{aligned}
\]

$2^\circ$ 设原式对 $n$ 成立：
\[
  f^{[n]}(x)=\lim_{h\to0+}\frac1{(2h)^n}
  \sum_{k=0}^n(-1)^kC_n^k[f((n-2k)h)].
\tag{1}
\]
下面来证对 $n+1$ 也必成立. 事实上，这时
\[
  f^{[n+1]}(x)=
  \lim_{h\to0+}\frac{\Delta_h}{2h}
  \left\{\frac1{(2h)^n}\sum_{k=0}^n(-1)^kC_n^k[f((n-2k)h)]\right\}.
\]
下面证明：此式
\[
  =
  \lim_{h\to0+}\frac1{(2h)^{n+1}}
  \sum_{k=0}^{n+1}(-1)^kC_{n+1}^k[f((n+1-2k)h)].
\]
在式 (1) 里，
\[
\begin{aligned}
&\Delta_h\sum_{k=0}^n(-1)^kC_n^k[f((n-2k)h)]\\
={}&
\sum_{k=0}^n(-1)^kC_n^k[f((n-2k)h+h)]
-
\sum_{k=0}^n(-1)^kC_n^k[f((n-2k)h-h)]\\
={}&[C_n^0f((n+1)h)-C_n^1f((n+1)h-2h)+\cdots+(-1)^nC_n^nf((n+1)h-2nh)]\\
&-[C_n^0f((n-1)h)+\cdots+(-1)^{n-1}C_n^{n-1}f((n-1)h-2(n-1)h)+(-1)^nf(-(n+1)h)]\\
={}&
C_{n+1}^0f((n+1)h)-C_{n+1}^1f((n+1)h-2h)+\cdots\\
&+(-1)^nC_{n+1}^nf((n+1)h-2nh)+(-1)^{n+1}C_{n+1}^{n+1}f(-(n+1)h)\\
={}&\sum_{k=0}^{n+1}(-1)^kC_{n+1}^kf((n+1-2k)h).
\end{aligned}
\]
所以
\[
  f^{[n+1]}(x)=
  \lim_{h\to0+}\frac1{(2h)^{n+1}}
  \sum_{k=0}^{n+1}(-1)^kC_{n+1}^kf((n+1-2k)h).
\]

\noindent\textbf{※3.2.34}\quad （Schwarz 定理）若 $f(x)$ 在 $[a,b]$ 上连续，$f(x)$ 的广义二阶导数
\[
  f^{[2]}(x)=\lim_{h\to0+}\frac{f(x+2h)-2f(x)+f(x-2h)}{4h^2}
\]
存在，且恒为零（对 $x\in(a,b)$）. 试证：$f(x)$ 是一次线性函数，可写为
\[
  f(x)=\alpha x+\beta,
\]
其中 $\alpha,\beta$ 为常数.（浙江大学）

\noindent\textbf{new 证}\quad （经典传统证法）作两辅助函数
\begin{equation}
  F_\pm(x)=\pm[f(x)-g(x)]+\varepsilon(x-a)(x-b)\quad(\varepsilon>0),
\tag{1}
\end{equation}
其中 $g(x)=\dfrac{f(b)-f(a)}{b-a}(x-a)+f(a)$（过两端点的直线）.

$1^\circ$ 已知在 $[a,b]$ 上 $(f(x))^{[2]}\equiv0$，因此
\[
  (f(x)-g(x))^{[2]}=f^{[2]}(x)-g^{[2]}(x)=0-\left[\frac{f(b)-f(a)}{b-a}(x-a)+f(a)\right]''\equiv0.
\]
故
\[
  (F_\pm(x))^{[2]}\equiv\pm(f(x)-g(x))^{[2]}+[\varepsilon(x-a)(x-b)]''\equiv2\varepsilon>0.
\tag{2}
\]

$2^\circ$ $F_\pm(x)=\pm[f(x)-g(x)]+\varepsilon(x-a)(x-b)$ 在端点 $x=a,x=b$ 处为零，为连续函数.

（i）若 $F_+(x)$ 在 $(a,b)$ 有正值，必有正的最大值点 $x_0\in(a,b)$，$F_+(x_0)\ge F_+(x_0\pm2h)$. 故
\[
  F_+^{[2]}(x_0)
  =
  \lim_{h\to0+}
  \frac{F_+(x_0+2h)-2F_+(x_0)+F_+(x_0-2h)}{4h^2}
  \le0,
\]
跟式 (2) 矛盾！故 $F_+(x)$ 没有正值.

（ii）若 $F_+(x)$ 在 $(a,b)$ 内有负值，则 $F_+(x)$ 在 $(a,b)$ 内必有负的最小值点 $x_0\in(a,b)$，
\[
  F_+(x_0)\le F_+(x_0\pm2h).
\]
注意此时有 $F_-(x_0)\ge F_-(x_0\pm2h)$，于是
\[
  F_-^{[2]}(x_0)
  =
  \lim_{h\to0+}
  \frac{F_-(x_0+2h)-2F_-(x_0)+F_-(x_0-2h)}{4h^2}
  \le0,
\]
也跟式 (2) 矛盾！故 $F_+(x)$ 没有负值.

总之，只可能是：$f(x)\equiv g(x)$（为一次线性函数）（亦即：过端点的直线），因此
\[
  f(x)=\alpha x+\beta\quad(\text{其中 }\alpha,\beta\text{ 是常数}),
\]
且
\[
  \alpha=\frac{f(b)-f(a)}{b-a},\quad
  \beta=f(a)-a\frac{f(b)-f(a)}{b-a}.
\]

\noindent\textbf{*3.2.35}\quad 设 $f(x)$ 在 $[a,b]$ 上连续，$f(a)<f(b)$；又设对一切 $x\in(a,b)$，$\lim\limits_{t\to0}\dfrac{f(x+t)-f(x-t)}t$ 存在，用 $g(x)$ 表示这一极限值. 试证：存在 $c\in(a,b)$，使得 $g(c)\ge0$.（南开大学）

（注意题目未假定导数存在.）

\noindent\textbf{提示}\quad （利用确界定理.）取 $\mu:f(a)<\mu<f(b)$，在 $[a,b]$ 上，让 $x$ 从 $a$ 点往右移动，保持 $f(x)\le\mu$，直到不能再前进为止，$x$ 到达它的上确界（记作 $c$），则 $c$ 能满足要求.

\noindent\textbf{证 I}\quad 在常值子区间 $(\alpha,\beta)$ 内，明显有 $g(x)=0$，问题自明. 下面的讨论将假设在 $(a,b)$ 内无常值子区间. 证明如下：

任取 $\mu:f(a)<\mu<f(b)$，作集合 $E=\{x\mid x\in[a,b]:\forall t\in[a,x],f(t)\le\mu\}$（意指 $[a,b]$ 内的点要成为 $E$ 的成员，除要求在此点有 $f(x)\le\mu$ 外，还要求在左侧区间 $[a,x]$ 有 $f(x)\le\mu$）. 因 $f(a)<\mu$，又 $f(x)$ 连续有保号性，故在 $a$ 的右侧充分小的邻域内，有 $x$ 使得 $f(x)<\mu$，说明 $E$ 非空；又 $E\subset[a,b]$，知 $E$ 有界. 根据确界定理，存在 $c=\sup E$，在 $[a,c]$ 上必有 $f(x)\le\mu$；在 $c$ 的右邻域至少存在串趋向 $c$ 的数列 $\{x_n=c+h_n\}\ (h_n>0)$ 使得 $f(x_n)>\mu$（不然 $E$ 的上确界会更大些），如上，$f(c-h_n)\le\mu$ 和 $f(c+h_n)>\mu\ (h_n>0,h_n\to0)$. 于是
\[
  g(c)=\lim_{n\to\infty}\frac{f(c+h_n)-f(c-h_n)}{2h_n}\ge0.
\]

\noindent\textbf{证 II}\quad （对偶地）将 $E$ 换成 $F=\{x\mid x\in[a,b]:\forall t\in[x,b],f(t)\ge\mu\}$，可类似找到欲求的点 $c$.

\noindent\textbf{注}\quad 若将 $E$ 改写为 $E'=\{x\mid x\in[a,b],f(x)\le\mu\}$ 或 $F'=\{x\mid x\in[a,b],f(x)\ge\mu\}$，将会导致错误. 如图 3.2.6，用 $E$ 可得出正确答案：$c=\sup E,\ \mu=f(c),\ g(c)>0$. 若采用 $E'$ 替代 $E$，得到的却是 $c'=\sup E'$，因 $f(x)$ 多次震荡，此时 $c'$ 恰是 $f$（最后一个）严格极小值点，故 $c'$ 处不能保证 $f$ 的对称差 $>0$；$F'$ 情况类似.

\begin{center}
  \includegraphics[width=0.42\textwidth]{3-2-6.png}
\end{center}

\noindent\textbf{new *3.2.36}\quad 设 $f(x)$ 在 $(a,b)$ 上连续，且 $\forall x\in(a,b)$，都有下极限：
\[
  \lim_{h\to0+}\frac{f(x+h)-f(x-h)}{2h}\ge0.
\tag{1}
\]
求证：$f(x)$ 在 $(a,b)$ 上单调不减.（北京大学）

\noindent\textbf{预备}\quad （对称导数是利用普通极限来定义的，但是本题是下极限，如果想用上题的思想方法，需用下极限定义（新的）对称导数，如下：
\[
  f^{[1]}_{\text{下}}(x)=\liminf_{h\to0+}\frac{f(x+h)-f(x-h)}{2h}.
\]
易证：若 $f(x)$ 在某点 $x_0$ 可导，则
\[
  f^{[1]}_{\text{下}}(x_0)=f^{[1]}(x_0)=f'(x_0).
\]
并且具有“线性可加性”，即 $[af(x)\pm bg(x)+c]^{[1]}_{\text{下}}=af^{[1]}_{\text{下}}(x)\pm bf^{[1]}_{\text{下}}(x)$.

\noindent\textbf{证}\quad $1^\circ$（反证法）若命题不成立（即：$f(x)$ 在 $(a,b)$ 上非“单调不减”），亦即至少存在两点 $\alpha,\beta$，使得 $a<\alpha<\beta<b$，但有 $f(\alpha)>f(\beta)$. 取实数 $k$：
\begin{equation}
  \frac{f(\alpha)-f(\beta)}{\beta-\alpha}<k<0.
\tag{2}
\end{equation}
作辅助函数
\begin{equation}
  F(x)=f(x)-f(\alpha)-k(x-\alpha).
\tag{3}
\end{equation}
则 $F(\alpha)=0,\ F(\beta)\overset{\text{式}(2)}{<}0$. 取 $\mu:F(\alpha)>\mu>F(\beta)$，作集合
\[
  E=\{x\mid x\in[\alpha,\beta]:\forall t\in[\alpha,x],F(t)\ge\mu\},
\]
明显 $\alpha\in E\subset[\alpha,b]$，故 $E$ 非空、有界，记 $c=\sup E$.（如上题所叙）可取充分小的数列 $\{h_n\}:h_n>0,\ h_n\to0$，使得
\[
  F(c+h_n)-F(c-h_n)<0.
\tag{4}
\]
$2^\circ$ 因题设，$f^{[1]}_{\text{下}}(x)$ 存在，根据式 (3) $F^{[1]}_{\text{下}}(x)$ 也存在，且 $F^{[1]}_{\text{下}}(c)=f^{[1]}_{\text{下}}(c)-k$. 既然
\[
  \lim_{t\to0+}\frac{F(c+t)-F(c-t)}{2t}
\]
存在，则 $\lim\limits_{n\to\infty}\dfrac{F(c+h_n)-F(c-h_n)}{2h_n}$ 也存在，且两者相等；由式 (4) 知
\[
  F^{[1]}_{\text{下}}(c)
  =
  \lim_{t\to0+}\frac{F(c+t)-F(c-t)}{2t}
  =
  \lim_{n\to\infty}\frac{F(c+h_n)-F(c-h_n)}{2h_n}
  \le0,
\]
其中分子 $F(c+h_n)-F(c-h_n)=f(c+h_n)-f(c-h_n)-2kh_n$. 故
\[
  \lim_{h\to0+}\frac{f(c+h)-f(c-h)}{2h}
  =
  \lim_{n\to\infty}\frac{F(c+h_n)-F(c-h_n)}{2h_n}
  \le k<0.
\]
与式 (1) 矛盾. 证毕.

\noindent\fbox{\parbox{0.86\textwidth}{
\noindent\textbf{小结}\quad 可放松为：若 $f'(x)\ge0$，则函数 $f(x)$ 单调不减；若 $f^{[1]}(x)\ge0$，则函数 $f(x)$ 单调不减；进一步放松为若 $f_{\text{下}}^{[1]}(x)\ge0$，则函数 $f(x)$ 单调不减.
}}

\section*{§3.3\quad Taylor 公式}

本节主要讨论带 Lagrange 余项与带 Peano 余项的 Taylor 公式在解题中的若干应用.

\noindent\textbf{要点}\quad $1^\circ$（Taylor 公式）若 $f^{(n)}(x)$ 在 $[a,b]$ 上连续，$f^{(n+1)}(x)$ 在 $(a,b)$ 内存在，则 $\forall x,x_0\in[a,b]$，$\exists \xi$ 位于 $x$ 与 $x_0$ 之间，使得下式成立：
\begin{equation}
  f(x)=f(x_0)+f'(x_0)(x-x_0)+\cdots+\frac1{n!}f^{(n)}(x_0)(x-x_0)^n+R_n(x),
\tag{1}
\end{equation}
其中 $R_n(x)=\dfrac1{(n+1)!}f^{(n+1)}(\xi)(x-x_0)^{n+1}$ 为 Lagrange 余项.

若 $f(x)$ 在 $x_0$ 处有 $n$ 阶导数 $f^{(n)}(x_0)$，则在 $x_0$ 邻域内 Taylor 公式 (1) 成立，其中
\[
  R_n(x)=o((x-x_0)^n)\quad(\text{当 }x\to x_0\text{ 时})
\]
为 Peano 余项.

$2^\circ$ 若把 $x_0$ 看成定点，$x$ 看成动点，则式 (1) 通过定点 $x_0$ 处的函数值 $f(x_0)$ 及导数值 $f'(x_0),\cdots,f^{(n)}(x_0)$ 表示动点 $x$ 处的函数值 $f(x)$. 当问题涉及二阶以上的导数时，通常可考虑用 Taylor 公式求解. 这里关键在于选取函数 $f$，点 $x_0$，展开的阶次 $n$ 以及余项形式. 根据需要，$x_0$ 一般应选在有特点的地方，例如使某 $f^{(i)}(x_0)=0$ 的地方等.

\subsection*{一、证明中值公式}

\noindent\textbf{例 3.3.1}\quad 设 $f(x)$ 在 $[a,b]$ 上三次可导，试证：$\exists c\in(a,b)$，使得
\begin{equation}
  f(b)=f(a)+f'\left(\frac{a+b}{2}\right)(b-a)+\frac1{24}f'''(c)(b-a)^3.
\tag{1}
\end{equation}

\noindent\textbf{证}\quad （待定常数法）设 $k$ 为使下式成立的实数：
\begin{equation}
  f(b)-f(a)-f'\left(\frac{a+b}{2}\right)(b-a)-\frac1{24}k(b-a)^3=0.
\tag{2}
\end{equation}
这时，我们的问题归为证明：$\exists c\in(a,b)$，使得
\begin{equation}
  k=f'''(c).
\tag{3}
\end{equation}
令
\begin{equation}
  g(x)=f(x)-f(a)-f'\left(\frac{a+x}{2}\right)(x-a)-\frac{k}{24}(x-a)^3.
\tag{4}
\end{equation}
则
\[
  g(a)=g(b)=0.
\]
根据 Rolle 定理，$\exists \xi\in(a,b)$，使得 $g'(\xi)=0$，由式 (4)，即
\begin{equation}
  f'(\xi)-f'\left(\frac{a+\xi}{2}\right)-f''\left(\frac{a+\xi}{2}\right)\frac{\xi-a}{2}
  -\frac{k}{8}(\xi-a)^2=0.
\tag{5}
\end{equation}
这是关于 $k$ 的方程，注意到 $f'(\xi)$ 在点 $\dfrac{a+\xi}{2}$ 处的 Taylor 公式：
\[
  f'(\xi)=f'\left(\frac{a+\xi}{2}\right)
  +f''\left(\frac{a+\xi}{2}\right)\frac{\xi-a}{2}
  +\frac12f'''(c)\left(\frac{\xi-a}{2}\right)^2,
\tag{6}
\]
其中 $c\in(a,b)$. 比较 (5), (6) 可得式 (3). 证毕.

\noindent\textbf{$\star$例 3.3.2}\quad 设 $f(x)$ 在 $[a,b]$ 上有三阶导数. 试证：$\exists c\in(a,b)$，使得
\begin{equation}
  \int_a^b f(x)\,\mathrm dx
  =(b-a)f\left(\frac{a+b}{2}\right)+\frac1{24}f''(c)(b-a)^3.
\tag{1}
\end{equation}
（西安电子科技大学，西安理工大学，东北大学）

\noindent\textbf{证 I}\quad 对函数 $F(x)=\int_a^x f(t)\,\mathrm dt$ 利用上例结果，或重复上例的证明即得.

\noindent\textbf{证 II}\quad 将函数 $F(x)=\int_a^x f(t)\,\mathrm dt$ 在点 $x_0=\dfrac{a+b}{2}$ 处按 Taylor 公式展开，记 $h=\dfrac{b-a}{2}$，则
\[
  F(x_0+h)=F(x_0)+f(x_0)h+\frac12f'(x_0)h^2+\frac16f''(\xi)h^3,
\]
\[
  F(x_0-h)=F(x_0)-f(x_0)h+\frac12f'(x_0)h^2-\frac16f''(\eta)h^3,
\]
其中 $\xi,\eta\in(a,b)$. 于是
\[
\begin{aligned}
\int_a^b f(x)\,\mathrm dx
&=F(x_0+h)-F(x_0-h)\\
&=(b-a)f(x_0)+\frac{(b-a)^3}{48}\bigl(f''(\xi)+f''(\eta)\bigr).
\end{aligned}
\tag{2}
\]
注意到导函数的介值性，$\exists c\in(a,b)$，使得
\[
  f''(c)=\frac{f''(\xi)+f''(\eta)}2.
\]
代入式 (2) 即得欲证的式 (1).

\noindent\textbf{注}\quad 例 3.3.1 亦可用此法证明.

\noindent\textbf{证 III}\quad 记 $x_0=\dfrac{a+b}{2}$，在 Taylor 展开式
\[
  f(x)=f(x_0)+f'(x_0)(x-x_0)+\frac12f''(\xi)(x-x_0)^2
\]
两端同时取 $[a,b]$ 上的积分. 注意右端第二项积分为 0. 第三项的积分，由于导数有介值性，第一积分中值定理成立：$\exists c\in(a,b)$，使得
\[
  \int_a^b f''(\xi)(x-x_0)^2\,\mathrm dx
  =
  f''(c)\int_a^b (x-x_0)^2\,\mathrm dx
  =
  \frac1{12}f''(c)(b-a)^3.
\]
因此式 (1) 成立.

\subsection*{二、用 Taylor 公式证明不等式}

\noindent\textbf{$\star$例 3.3.3}\quad 设 $f(x)$ 在 $[a,b]$ 上二次可微，$f''(x)<0$. 试证：$\forall a\le x_1<x_2<\cdots<x_n\le b,\ k_i\ge0,\ \sum\limits_{i=1}^n k_i=1$，有
\[
  f\left(\sum_{i=1}^n k_ix_i\right)>\sum_{i=1}^n k_if(x_i).
\]
（哈尔滨工业大学，北京科技大学）

\noindent\textbf{证}\quad 取 $x_0=\sum\limits_{i=1}^n k_ix_i$，将 $f(x_i)$ 在 $x=x_0$ 处展开：
\[
  f(x_i)=f(x_0)+f'(x_0)(x_i-x_0)+\frac12f''(\xi_i)(x_i-x_0)^2
  <f(x_0)+f'(x_0)(x_i-x_0)\quad(i=1,2,\cdots,n).
\]
以 $k_i$ 乘此式两端，然后 $n$ 个不等式相加，注意 $\sum\limits_{i=1}^n k_i=1$，
\[
  \sum_{i=1}^n k_i(x_i-x_0)=\sum_{i=1}^n k_ix_i-x_0=0,
\]
得
\[
  \sum_{i=1}^n k_if(x_i)<f(x_0)=f\left(\sum_{i=1}^n k_ix_i\right).
\]

\noindent\textbf{例 3.3.4}\quad 设 $f(x)$ 有二阶导数，
\begin{equation}
  f(x)\le\frac12[f(x-h)+f(x+h)].
\tag{1}
\end{equation}
试证 $f''(x)\ge0$.（北京师范大学）

\noindent\textbf{证}\quad
\[
  f(x\pm h)=f(x)\pm f'(x)h+\frac12f''(x)h^2+o(h^2).
\]
两式相加，并除以 $h^2$，注意 (1)，有 $f''(x)+o(1)\ge0$. 令 $h\to0$，取极限得 $f''(x)\ge0$.

\subsection*{三、用 Taylor 公式作导数的中值估计}

\noindent\textbf{例 3.3.5}\quad 设 $f(x)$ 二次可微，且 $f(0)=f(1)=0,\ \max\limits_{0\le x\le1}f(x)=2$，试证 $\min\limits_{0\le x\le1}f''(x)\le-16$.

\noindent\textbf{证}\quad 因 $f(x)$ 在 $[0,1]$ 上连续，有最大、最小值. 又因 $\max\limits_{0\le x\le1}f(x)=2,\ f(0)=f(1)=0$，故最大值在 $(0,1)$ 内部达到. 所以，$\exists x_0\in(0,1)$ 使得 $f(x_0)=\max\limits_{0\le x\le1}f(x)$，于是 $f(x_0)$ 为极大值. 由 Fermat 定理，有 $f'(x_0)=0$.

在 $x=x_0$ 处按 Taylor 公式展开，$\exists \xi,\eta\in(0,1)$ 使得
\[
  0=f(0)=f(x_0)+\frac12f''(\xi)(0-x_0)^2=2+\frac12f''(\xi)x_0^2,
\]
\[
  0=f(1)=f(x_0)+\frac12f''(\eta)(1-x_0)^2=2+\frac12f''(\eta)(1-x_0)^2.
\]
因此，
\[
  \min_{0\le x\le1}f''(x)
  \le
  \min\{f''(\xi),f''(\eta)\}
  =
  \min\left\{-\frac4{x_0^2},-\frac4{(1-x_0)^2}\right\}.
\]
而 $x_0\in\left[\dfrac12,1\right]$ 时，
\[
  \min\left\{-\frac4{x_0^2},-\frac4{(1-x_0)^2}\right\}
  =
  -\frac4{(1-x_0)^2}\le-16;
\]
$x_0\in\left[0,\dfrac12\right]$ 时，
\[
  \min\left\{-\frac4{x_0^2},-\frac4{(1-x_0)^2}\right\}
  =
  -\frac4{x_0^2}\le-16,
\]
所以
\[
  \min_{0\le x\le1}f''(x)\le-16.
\]

\noindent\textbf{读者试证}\quad 若 $f(x)$ 在 $[0,1]$ 上有二阶导数，$f(0)=f(1)=0,\ \min\limits_{0\le x\le1}f(x)=-1$，则 $\max\limits_{0\le x\le1}f''(x)\ge8$.

\noindent\textbf{new 练习}\quad 已知 $f(x)$ 在 $[0,1]$ 二阶可导，且 $f(0)=0,\ f(1)=3,\ \min\limits_{x\in[0,1]}f(x)=-1$. 求证：存在一点 $c\in(0,1)$，满足 $f''(c)\ge18$.（南开大学）

\noindent\textbf{提示}\quad 用例 3.3.5 的方法.

\noindent\textbf{再提示}\quad 必存在 $x_0\in(0,1)$，使得 $f(x_0)=\min\limits_{x\in[0,1]}f(x)=-1$，且 $x_0$ 是极值点，有 $f'(x_0)=0$.

若 $x_0\in\left(0,\dfrac13\right)$，用 Taylor 公式：$\exists \xi\in(0,x_0)$，使得
\[
  0=f(0)=f(x_0)+f'(x_0)(0-x_0)+\frac12f''(\xi)(0-x_0)^2
  =-1+\frac12f''(\xi)x_0^2
  \le -1+\frac1{18}f''(\xi)
\]
（因 $x_0<\dfrac13$），即有 $c=\xi$，使得 $f''(c)\ge18$.

类似地，若 $x_0\in\left[\dfrac13,1\right)$，由 $f(1)=3$，可得 $f''(\eta)\dfrac{8}{(1-x_0)^2}\ge18$. 令 $c=\eta$，则 $f''(c)\ge18$.

\noindent\textbf{$\star$例 3.3.6}\quad 若 $f(x)$ 在 $[a,b]$ 上有二阶导数，$f'(a)=f'(b)=0$，试证：$\exists \xi\in(a,b)$，使得
\begin{equation}
  |f''(\xi)|\ge\frac4{(b-a)^2}|f(b)-f(a)|.
\tag{1}
\end{equation}

\noindent\textbf{证 I}\quad 应用 Taylor 公式，将 $f\left(\dfrac{a+b}{2}\right)$ 分别在点 $a,b$ 展开，注意 $f'(a)=f'(b)=0$，
\[
  \exists \zeta,\eta:a<\zeta<\frac{a+b}{2}<\eta<b,
\]
使得
\[
  f\left(\frac{a+b}{2}\right)=f(a)+\frac12f''(\zeta)\left(\frac{b-a}{2}\right)^2,
\tag{2}
\]
\[
  f\left(\frac{a+b}{2}\right)=f(b)+\frac12f''(\eta)\left(\frac{b-a}{2}\right)^2.
\tag{3}
\]
(3)-(2) 得 $f(b)-f(a)+\dfrac18[f''(\eta)-f''(\zeta)](b-a)^2=0$. 故
\[
  \frac{4|f(b)-f(a)|}{(b-a)^2}
  \le\frac12(|f''(\zeta)|+|f''(\eta)|)\le |f''(\xi)|,
\]
其中
\[
  \xi=
  \begin{cases}
    \zeta,& |f''(\zeta)|\ge |f''(\eta)|,\\
    \eta,& |f''(\zeta)|<|f''(\eta)|.
  \end{cases}
\]

\noindent\textbf{证 II}\quad 若 $f(a)=f(b)$，问题自明. 设 $f(a)<f(b)$（$f(a)>f(b)$ 类似可证），记 $c=\dfrac{a+b}{2}$.

$1^\circ$ 若 $f(c)\ge\dfrac{f(a)+f(b)}2$（此时 $2[f(c)-f(a)]\ge f(b)-f(a)$），由式 (2) 可知
\[
  |f(b)-f(a)|=f(b)-f(a)\le2[f(c)-f(a)]
  =f''(\zeta)\left(\frac{b-a}{2}\right)^2.
\]
所以
\[
  |f''(\zeta)|\ge\frac4{(b-a)^2}|f(b)-f(a)|.
\]

$2^\circ$ 若 $f(c)<\dfrac{f(a)+f(b)}2$，类似可由式 (3) 推出
\[
  |f''(\eta)|\ge\frac4{(b-a)^2}|f(b)-f(a)|.
\]

\noindent\textbf{证 III}\quad （采用辅助函数.）设 $f(a)<f(b),\ c=\dfrac{a+b}{2}$.

$1^\circ$ 若 $f(c)\ge\dfrac{f(a)+f(b)}2$，作辅助函数
\[
  F(x)=f(x)-\frac{k}{2}(x-a)^2
  \quad\left(k=4\frac{f(b)-f(a)}{(b-a)^2}\right).
\]
（只要证明 $\exists \xi\in(a,b)$ 使得 $F''(\xi)\ge0$ 即可.）因 $F'(a)=0$，
\[
  F(c)=f(c)-\frac{k}{2}\frac{(b-a)^2}{4}
  \ge\frac{f(a)+f(b)}2-\frac{f(b)-f(a)}2
  =f(a)=F(a),
\]
所以
\[
  0\le F(c)-F(a)=\frac12F''(\xi)(c-a)^2\quad(\xi\in(a,c)).
\]
故 $F''(\xi)\ge0$. 即 $|f''(\xi)|\ge f''(\xi)\ge k=\dfrac4{(b-a)^2}|f(b)-f(a)|$.

$2^\circ$ 若 $f(c)<\dfrac{f(a)+f(b)}2$，可作 $F(x)=f(x)+\dfrac{k}{2}(x-b)^2$，类似可证.

（本题还可写出一些别的证法，也可不用 Taylor 公式.）读者试用本例结果证明：若火车从起点到终点共走了 $t\ \mathrm s$，行驶了 $s\ \mathrm m$，则途中必有一个时刻，其加速度的绝对值不低于 $\dfrac{4s}{t^2}\ \mathrm{m/s^2}$.

\subsection*{四、关于界的估计}

\noindent\textbf{$\star$例 3.3.7}\quad 设 $f(x)$ 在 $[0,1]$ 上有二阶导数，$0\le x\le1$ 时，$|f(x)|\le1,\ |f''(x)|<2$. 试证：当 $0\le x\le1$ 时，$|f'(x)|\le3$.（南京航空航天大学）

\noindent\textbf{证}\quad
\[
  f(1)=f(x)+f'(x)(1-x)+\frac12f''(\xi)(1-x)^2,
\]
\[
  f(0)=f(x)+f'(x)(-x)+\frac12f''(\eta)(-x)^2,
\]
所以
\[
  f(1)-f(0)=f'(x)+\frac12f''(\xi)(1-x)^2-\frac12f''(\eta)x^2,
\]
\[
  |f'(x)|\le |f(1)|+|f(0)|+\frac12|f''(\xi)|(1-x)^2+\frac12|f''(\eta)|x^2
  \le2+(1-x)^2+x^2\le2+1=3.
\]

\noindent\textbf{例 3.3.8}\quad 设 $f(x)$ 在 $[0,1]$ 上具有二阶连续导数，且满足 $f(1)=f(0)$ 及 $|f''(x)|\le M\ (x\in[0,1])$. 试证：对一切 $x\in[0,1]$，有 $|f'(x)|\le\dfrac M2$.（上海师范大学）

\noindent\textbf{提示}\quad 与上例类似展开.

\noindent\textbf{*例 3.3.9}\quad 设 $f(x)\ (-\infty<x<+\infty)$ 为二次可微函数，
\[
  M_k=\sup_{-\infty<x<+\infty}|f^{(k)}(x)|<+\infty\quad(k=0,2),
\]
试证：
\[
  M_1=\sup_{-\infty<x<+\infty}|f'(x)|<+\infty\quad\text{且 }M_1^2\le2M_0M_2,
\]
其中 $f^{(0)}(x)$ 表示 $f(x)$.（北京大学）

\noindent\textbf{证 I}\quad
\[
  f(x+h)=f(x)+f'(x)h+\frac12f''(\xi)h^2\quad(\xi\text{ 在 }x\text{ 与 }x+h\text{ 之间}),
\]
\[
  f(x-h)=f(x)-f'(x)h+\frac12f''(\eta)h^2\quad(\eta\text{ 在 }x-h\text{ 与 }x\text{ 之间}),
\]
两式相减，得
\[
  f(x+h)-f(x-h)=2f'(x)h+\frac{h^2}{2}\bigl[f''(\xi)-f''(\eta)\bigr],
\]
从而
\[
  2|f'(x)|h\le |f(x+h)|+|f(x-h)|+\frac{h^2}{2}\bigl(|f''(\xi)|+|f''(\eta)|\bigr)
  \le 2M_0+h^2M_2. \tag{1}
\]
因此，对一切 $h>0$，二次三项式
\[
  M_2h^2-2|f'(x)|h+2M_0\ge0,
\]
故其判别式小于等于零，即
\[
  |f'(x)|^2-2M_0M_2\le0.
\]
故对一切 $x\in(-\infty,+\infty)$，有
\[
  |f'(x)|^2\le2M_0M_2.
\]
所以
\[
  M_1=\sup_{-\infty<x<+\infty}|f'(x)|<+\infty,\qquad M_1^2\le2M_0M_2.
\]

\noindent\textbf{证 II}\quad 由式 (1) 可得
\[
  |f'(x)|\le \frac{M_0}{h}+\frac{hM_2}{2},\quad \forall h>0. \tag{2}
\]
由于
\[
  \frac{M_0}{h}\cdot\frac{hM_2}{2}=\frac12M_0M_2
\]
为定值，故当 $\dfrac{M_0}{h}=\dfrac{hM_2}{2}$ 时，式 (2) 右端取得最小值. 令
\[
  h=\sqrt{\frac{2M_0}{M_2}},
\]
则
\[
  |f'(x)|\le\sqrt{2M_0M_2},\quad \forall x\in(-\infty,+\infty),
\]
故 $M_1^2\le2M_0M_2$.

\noindent\textbf{例 3.3.10}\quad 设函数 $f(x)$ 在 $(-\infty,+\infty)$ 内有三阶导数，并且 $f(x)$ 和 $f'''(x)$ 在 $(-\infty,+\infty)$ 内有界. 证明：$f'(x)$ 和 $f''(x)$ 也在 $(-\infty,+\infty)$ 内有界.

\noindent\textbf{证}\quad 由 Taylor 公式，
\[
  f(x+h)=f(x)+f'(x)h+\frac12f''(x)h^2+\frac{1}{3!}f'''(\xi)h^3.
\]
令 $h=\pm1$，得
\[
  f(x+1)=f(x)+f'(x)+\frac12f''(x)+\frac16f'''(\xi), \tag{1}
\]
\[
  f(x-1)=f(x)-f'(x)+\frac12f''(x)-\frac16f'''(\eta). \tag{2}
\]
式 (1) $-$ 式 (2)，得
\[
  f(x+1)-f(x-1)=2f'(x)+\frac16\bigl[f'''(\xi)+f'''(\eta)\bigr].
\]
所以
\[
  2|f'(x)|\le2M_0+\frac13M_3
  \quad(\forall x\in(-\infty,+\infty)),
\]
其中
\[
  M_k=\sup_{-\infty<x<+\infty}|f^{(k)}(x)|\quad(k=0,3).
\]
同理，由式 (1) $+$ 式 (2)，得
\[
  |f''(x)|\le4M_0+\frac13M_3
  \quad(\forall x\in(-\infty,+\infty)).
\]
故 $f'(x),f''(x)$ 有界.

\noindent\textbf{*练习}\quad 设 $f(x)$ 在 $(-1,1)$ 内有二阶导数，且
\[
  f(0)=f'(0)=0,\qquad |f''(x)|^2\le |f(x)\cdot f'(x)|.
\]
试证：在 $(-1,1)$ 内 $f(x)\equiv0$.（南开大学）

\noindent\textbf{分析}\quad 因为 $f(0)=f'(0)=0$，由 Taylor 公式知，对任意 $x\in(-1,1)$，存在 $\xi$ 在 $0$ 与 $x$ 之间，使得
\[
  f(x)=f(0)+f'(0)x+\frac12f''(\xi)x^2=\frac12f''(\xi)x^2.
\]
从而，在 $(-1,1)$ 内，只要证明 $f''(x)\equiv0$，则必有 $f(x)\equiv0$.

\noindent\textbf{证}\quad 对任一闭区间 $[-a,a]\subset(-1,1)$，记
\[
  M=\max_{-a\le x\le a}|f(x)|,\quad
  M_1=\max_{-a\le x\le a}|f'(x)|,\quad
  M_2=\sup_{-a\le x\le a}|f''(x)|.
\]
则
\[
  |f''(x)|^2\le |f(x)\cdot f'(x)|\le MM_1,\quad \forall x\in[-a,a],
\]
取上确界得
\[
  M_2\le\sqrt{MM_1}. \tag{1}
\]
又
\[
  |f(x)|=|f(x)-f(0)|=|f'(\xi)x|\le aM_1,
\]
取最大值得
\[
  M\le aM_1. \tag{2}
\]
同理
\[
  |f'(x)|=|f'(x)-f'(0)|=|f''(\eta)x|\le aM_2,
\]
故
\[
  M_1\le aM_2. \tag{3}
\]
由式 (1)、(2)、(3)，得
\[
  0\le M_2\le\sqrt{MM_1}\le a^{3/2}M_2
\]
其中 $0<a^{3/2}<1$，故 $M_2=\sup_{-a\le x\le a}|f''(x)|=0$. 因此 $f''(x)\equiv0\ (-a\le x\le a)$. 由于 $[-a,a]\subset(-1,1)$ 任意，故 $f''(x)\equiv0\ (-1<x<1)$，由分析可知 $f(x)\equiv0$.

\subsection*{五、求无穷远处的极限}

Taylor 公式在求极限里有广泛应用，在第一章讲极限时已作了介绍，请见例 1.3.13. 这里只作适当补充.

\noindent\textbf{$\star$例 3.3.11}\quad 设函数 $\varphi(x)$ 在 $[0,+\infty)$ 上二次连续可微，如果 $\lim\limits_{x\to+\infty}\varphi(x)$ 存在，且 $\varphi''(x)$ 在 $[0,+\infty)$ 上有界. 试证：
\[
  \lim_{x\to+\infty}\varphi'(x)=0.
\]
（吉林大学，北京理工大学）

\noindent\textbf{证 I}\quad 即证
\[
  \forall \varepsilon>0,\ \exists \Delta>0,\text{ 当 }x>\Delta\text{ 时， }|\varphi'(x)|<\varepsilon.
\]
应用 Taylor 公式，
\[
  \varphi(x+h)=\varphi(x)+\varphi'(x)h+\frac12\varphi''(\xi)h^2,
\]
即
\[
  \varphi'(x)=\frac1h[\varphi(x+h)-\varphi(x)]-\frac12\varphi''(\xi)h. \tag{1}
\]
记 $A=\lim\limits_{x\to+\infty}\varphi(x)$. 因 $\varphi''(x)$ 有界，故存在 $M>0$，使得
\[
  |\varphi''(x)|\le M\quad(\forall x\ge0).
\]
由式 (1)，得
\[
  |\varphi'(x)|\le\frac1h\bigl(|\varphi(x+h)-A|+|A-\varphi(x)|\bigr)+\frac12Mh. \tag{2}
\]
对任意 $\varepsilon>0$，取 $h>0$ 充分小，使得 $\dfrac12Mh<\dfrac{\varepsilon}{2}$. 固定此 $h$，由 $\lim\limits_{x\to+\infty}\varphi(x)=A$ 知，存在 $\Delta>0$，当 $x>\Delta$ 时，
\[
  \frac1h\bigl(|\varphi(x+h)-A|+|A-\varphi(x)|\bigr)<\frac{\varepsilon}{2}.
\]
于是由式 (2) 得 $|\varphi'(x)|<\varepsilon$.

\noindent\textbf{证 II}\quad $1^\circ$ 由于 $\varphi''(x)$ 有界，故存在 $M>0$，使得 $|\varphi''(x)|\le M\ (\forall x\ge0)$. 因此，对任意 $\varepsilon>0$，令
\[
  \delta=\frac{\varepsilon}{2(M+1)}>0.
\]
当 $x_1>x\ge0$ 且 $x_1-x\le\delta$ 时，有
\[
  |\varphi'(x_1)-\varphi'(x)|=|\varphi''(\xi)(x_1-x)|
  \le \frac{M\varepsilon}{2(M+1)}<\frac{\varepsilon}{2}. \tag{3}
\]

$2^\circ$ 由已知 $\lim\limits_{x\to+\infty}\varphi(x)$ 存在，根据 Cauchy 准则，对于上述 $\varepsilon>0$，存在 $\Delta>0$，当 $x_1>x>\Delta$ 时，有
\[
  |\varphi(x_1)-\varphi(x)|<\frac{\varepsilon^2}{4(M+1)}.
\]
特别地，令 $x_1=x+\dfrac{\varepsilon}{2(M+1)}$，有
\[
  \left|\varphi'(\xi_x)\frac{\varepsilon}{2(M+1)}\right|
  =\left|\varphi(x)-\varphi\left(x+\frac{\varepsilon}{2(M+1)}\right)\right|
  <\frac{\varepsilon^2}{4(M+1)},
\]
其中 $x<\xi_x<x+\dfrac{\varepsilon}{2(M+1)}$，因此
\[
  |\varphi'(\xi_x)|<\frac{\varepsilon}{2}. \tag{4}
\]
又因 $|x-\xi_x|<\dfrac{\varepsilon}{2(M+1)}=\delta$，由式 (3) 得
\[
  |\varphi'(x)-\varphi'(\xi_x)|<\frac{\varepsilon}{2}. \tag{5}
\]
由式 (4)、(5)，得
\[
  |\varphi'(x)|\le |\varphi'(x)-\varphi'(\xi_x)|+|\varphi'(\xi_x)|<\varepsilon.
\]
故 $\lim\limits_{x\to+\infty}\varphi'(x)=0$.

\noindent\textbf{注}\quad 1）本题曾在博士数学论坛讨论过，证 II 由张祖锦先生提供.

2）读者熟知反常积分后，可证明下述练习.

\noindent\textbf{*练习}\quad 若 $g(x)$ 在 $[0,+\infty)$ 上一致连续，且反常积分
\[
  \int_0^{+\infty}g(t)\,\mathrm dt
\]
收敛，则
\[
  \lim_{x\to+\infty}g(x)=0.
\]
（南开大学）

\noindent\textbf{证}\quad $1^\circ$ 令
\[
  f(x)=\int_0^xg(t)\,\mathrm dt,
\]
按题设，$f'(x)=g(x)$ 在 $[0,+\infty)$ 上一致连续：$\forall\varepsilon>0$，存在 $\delta:0<\delta<\varepsilon$，$\forall x_1>x\ge0$，当 $x_1-x\le\delta$ 时，有
\[
  |f'(x)-f'(x_1)|=|g(x)-g(x_1)|<\frac{\varepsilon}{2}. \tag{1}
\]

$2^\circ$ $\lim\limits_{x\to+\infty}f(x)=\lim\limits_{x\to+\infty}\int_0^xg(t)\,\mathrm dt$ 存在，根据 Cauchy 准则，对上述 $\varepsilon>0$ 和 $\delta(0<\delta<\varepsilon)$，存在 $\Delta>0$，当 $x_1>x>\Delta$ 时，有
\[
  |f(x_1)-f(x)|<\frac{\delta^2}{4}.
\]
特别地，令 $x_1=x+\dfrac{\delta}{2}$，有
\[
  \left|f'(\xi_x)\frac{\delta}{2}\right|
  =\left|f(x)-f\left(x+\frac{\delta}{2}\right)\right|
  <\frac{\delta^2}{4}
  \quad\left(x<\xi_x<x+\frac{\delta}{2}\right).
\]
因此
\[
  |f'(\xi_x)|<\frac{\delta}{2}<\frac{\varepsilon}{2}. \tag{2}
\]
由 $|x-\xi_x|<\dfrac{\delta}{2}$ 和式 (1)，得
\[
  |f'(x)-f'(\xi_x)|<\frac{\varepsilon}{2}. \tag{3}
\]
利用式 (2) 和 (3)，得
\[
  |f'(x)|\le |f'(x)-f'(\xi_x)|+|f'(\xi_x)|=\frac{\varepsilon}{2}+\frac{\varepsilon}{2}=\varepsilon.
\]
故
\[
  \lim_{x\to+\infty}g(x)=\lim_{x\to+\infty}f'(x)=0.
\]

\noindent\textbf{*例 3.3.12}\quad 设 $f(x)$ 至少有 $k$ 阶导数，且对某个实数 $\alpha$ 有
\[
  \lim_{x\to+\infty}x^\alpha f(x)=0,\qquad
  \lim_{x\to+\infty}x^\alpha f^{(k)}(x)=0, \tag{1}
\]
试证：
\[
  \lim_{x\to+\infty}x^\alpha f^{(i)}(x)=0,\quad i=0,1,2,\cdots,k,
\]
$f^{(0)}(x)$ 表示 $f(x)$.

\noindent\textbf{证}\quad 根据已知条件 (1)，要证明 $\lim\limits_{x\to+\infty}x^\alpha f^{(i)}(x)=0$，只要把 $f^{(i)}(x)$ 写成 $f(x)$ 与 $f^{(k)}(x)$ 的线性组合即可. 应用 Taylor 公式，
\[
  f(x+m)=f(x)+mf'(x)+\frac{m^2}{2!}f''(x)+\cdots
  +\frac{m^{k-1}}{(k-1)!}f^{(k-1)}(x)+\frac{m^k}{k!}f^{(k)}(\xi_m), \tag{2}
\]
其中 $x<\xi_m<x+m,\ m=1,2,\cdots,k$. 这是关于 $f'(x),f''(x),\cdots,f^{(k-1)}(x)$ 的线性方程组，其系数行列式为
\[
  \begin{vmatrix}
  1&1&\dfrac1{2!}&\cdots&\dfrac1{(k-1)!}\\
  1&2&\dfrac{2^2}{2!}&\cdots&\dfrac{2^{k-1}}{(k-1)!}\\
  \vdots&\vdots&\vdots& &\vdots\\
  1&k&\dfrac{k^2}{2!}&\cdots&\dfrac{k^{k-1}}{(k-1)!}
  \end{vmatrix}
\]
\[
  =\frac{1}{1!\,2!\cdots(k-1)!}
  \begin{vmatrix}
  1&1&1&\cdots&1\\
  1&2&2^2&\cdots&2^{k-1}\\
  1&3&3^2&\cdots&3^{k-1}\\
  \vdots&\vdots&\vdots& &\vdots\\
  1&k&k^2&\cdots&k^{k-1}
  \end{vmatrix}=1.
\]
（其中后一行列式是著名的 Vandermonde 行列式，它的值为 $1!2!\cdots(k-1)!$.）故从方程组 (2)，可把 $f'(x),f''(x),\cdots,f^{(k-1)}(x)$ 写成 $f(x+m)\ (m=1,2,\cdots,k)$ 与 $f^{(k)}(\xi_m)\ (m=1,2,\cdots,k)$ 的线性组合. 至此，我们只要证明
\[
  \lim_{x\to+\infty}x^\alpha f(x+m)=\lim_{x\to+\infty}x^\alpha f^{(k)}(\xi_m)=0\quad(m=1,2,\cdots,k)
\]
即可. 事实上，设 $x\le t\le x+k$，则
\[
  \lim_{x\to+\infty}x^\alpha f^{(i)}(t)
  =\lim_{x\to+\infty}\left(\frac{x}{t}\right)^\alpha t^\alpha f^{(i)}(t)
  =\lim_{x\to+\infty}\left(\frac{x}{t}\right)^\alpha\cdot\lim_{t\to+\infty}t^\alpha f^{(i)}(t)
  =1\cdot0=0\quad(i=0,k).
\]
在此式中令 $t=x+m,\ i=0$ 和 $t=\xi_m,\ i=k$，则得
\[
  \lim_{x\to+\infty}x^\alpha f(x+m)
  =\lim_{x\to+\infty}x^\alpha f^{(k)}(\xi_m)=0\quad(m=1,2,\cdots,k).
\]
证毕.

\subsection*{六、中值点的极限}

\noindent\textbf{*例 3.3.13}\quad 设

1）$f(x)$ 在 $(x_0-\delta,x_0+\delta)$ 内是 $n$ 阶连续可微函数，此处 $\delta>0$；

2）当 $k=2,3,\cdots,(n-1)$ 时，有 $f^{(k)}(x_0)=0$，但是 $f^{(n)}(x_0)\ne0$；

3）当 $0\ne |h|<\delta$ 时，有
\[
  \frac{f(x_0+h)-f(x_0)}{h}=f'(x_0+h\cdot\theta(h))\quad(0<\theta(h)<1). \tag{1}
\]
证明：
\[
  \lim_{h\to0}\theta(h)=\sqrt[n-1]{\frac1n}.
\]
（西安电子科技大学）

\noindent\textbf{证}\quad 我们要设法从式 (1) 中解出 $\theta(h)$. 为此，我们将式 (1) 左端的 $f(x_0+h)$ 及右端的 $f'(x_0+h\cdot\theta(h))$ 在 $x_0$ 处展开. 注意条件 2），知存在 $\theta_1,\theta_2\in(0,1)$，使得
\[
  f(x_0+h)=f(x_0)+hf'(x_0)+\frac{h^n}{n!}f^{(n)}(x_0+\theta_1h),
\]
\[
  f'(x_0+h\theta(h))=f'(x_0)+\frac{h^{n-1}\cdot(\theta(h))^{n-1}}{(n-1)!}f^{(n)}(x_0+\theta_2h\theta(h)).
\]
于是，式 (1) 变成
\[
  f'(x_0)+\frac{h^{n-1}}{n!}f^{(n)}(x_0+\theta_1h)
  =f'(x_0)+\frac{h^{n-1}\cdot(\theta(h))^{n-1}}{(n-1)!}f^{(n)}(x_0+\theta_2h\theta(h)),
\]
从而
\[
  \theta(h)=
  \sqrt[n-1]{\frac{f^{(n)}(x_0+\theta_1h)}{n\cdot f^{(n)}(x_0+\theta_2h\theta(h))}}.
\]
因 $\theta_1,\theta_2,\theta(h)\in(0,1)$，利用 $f^{(n)}(x)$ 的连续性，由此可得
\[
  \lim_{h\to0}\theta(h)=\sqrt[n-1]{\frac1n}.
\]

\noindent\textbf{例 3.3.14}\quad 设
\[
  f(x+h)=f(x)+hf'(x)+\cdots+\frac{h^n}{n!}f^{(n)}(x+\theta\cdot h)\quad(0<\theta<1),
\]
且 $f^{(n+1)}(x)\ne0$，证明 $\lim\limits_{h\to0}\theta=\dfrac1{n+1}$.

\noindent\textbf{提示}\quad
\[
  f(x+h)=f(x)+hf'(x)+\cdots+\frac{h^n}{n!}f^{(n)}(x)
  +\frac{h^{n+1}}{(n+1)!}f^{(n+1)}(x)+o(h^{n+1}).
\]
从而有
\[
  \theta h\frac{f^{(n)}(x+\theta h)-f^{(n)}(x)}{\theta h}
  =\frac{h}{n+1}f^{(n+1)}(x)+o(h).
\]

\subsection*{七、函数方程中的应用}

\noindent\textbf{例 3.3.15}\quad 设 $f(x)$ 在 $(-\infty,+\infty)$ 内有连续三阶导数，且满足方程：
\[
  f(x+h)=f(x)+hf'(x+\theta h),\quad 0<\theta<1\ (\theta\text{ 与 }h\text{ 无关}). \tag{1}
\]
试证：$f(x)$ 是一次或二次函数.

\noindent\textbf{证}\quad 问题在于证明：$f''(x)\equiv0$ 或 $f'''(x)\equiv0$. 为此将式 (1) 对 $h$ 求导，注意 $\theta$ 与 $h$ 无关，有
\[
  f'(x+h)=f'(x+\theta h)+\theta hf''(x+\theta h). \tag{2}
\]
从而
\[
  \frac{f'(x+h)-f'(x)+f'(x)-f'(x+\theta h)}{h}
  =\theta f''(x+\theta h).
\]
令 $h\to0$ 取极限，得
\[
  f''(x)-\theta f''(x)=\theta f''(x),\quad f''(x)=2\theta f''(x).
\]
若 $\theta\ne\dfrac12$，由此知 $f''(x)\equiv0$，$f(x)$ 为一次函数；若 $\theta=\dfrac12$，式 (2) 给出
\[
  f'(x+h)=f'\left(x+\frac12h\right)+\frac12hf''\left(x+\frac12h\right).
\]
此式两端同时对 $h$ 求导，减去 $f''(x)$，除以 $h$，然后令 $h\to0$，取极限即得 $f'''(x)\equiv0$，$f(x)$ 为二次函数.

在一定条件下证明某函数 $f(x)\equiv0$ 的问题，我们称之为归零问题. 因此上例实际上是 $f',f''$ 的归零问题. 例 3.2.18 也是归零问题，下面让我们再看一例.

\noindent\textbf{例 3.3.16}\quad 已知函数 $f(x)$ 在区间 $(-1,1)$ 内有二阶导数，且 $f(0)=f'(0)=0$，
\[
  |f''(x)|\le |f(x)|+|f'(x)|. \tag{1}
\]
试证：$\exists\delta>0$，使得 $(-\delta,\delta)$ 内 $f(x)\equiv0$.（武汉大学赛题）

\noindent\textbf{证}\quad 为了证明 $f(x)$ 在 $x=0$ 的邻域内恒为零，我们将式 (1) 右端的 $f(x),f'(x)$ 在 $x=0$ 处按 Taylor 公式展开. 注意到 $f(0)=f'(0)=0$，有
\[
  f(x)=f(0)+f'(0)x+\frac{f''(\xi)}2x^2=\frac12f''(\xi)x^2,
\]
\[
  f'(x)=f'(0)+f''(\eta)x=f''(\eta)x.
\]
从而
\[
  |f(x)|+|f'(x)|=\left|\frac12f''(\xi)x^2\right|+|f''(\eta)x|. \tag{2}
\]
令限制 $x\in\left[-\dfrac14,\dfrac14\right]$，则 $|f(x)|+|f'(x)|$ 在 $\left[-\dfrac14,\dfrac14\right]$ 上连续有界，$\exists x_0\in\left[-\dfrac14,\dfrac14\right]$，使得
\[
  |f(x_0)|+|f'(x_0)|=\max_{-\frac14\le x\le\frac14}\{|f(x)|+|f'(x)|\}=M.
\]
我们只要证明 $M=0$ 即可. 事实上，
\[
  M=|f(x_0)|+|f'(x_0)|
  \overset{\text{式}(2)}{\le}
  \left|\frac12f''(\xi_0)x_0^2\right|+|f''(\eta_0)x_0|
\]
\[
  \le \frac14\bigl(|f''(\xi_0)|+|f''(\eta_0)|\bigr)
  \le \frac14\bigl(|f'(\xi_0)|+|f(\xi_0)|+|f'(\eta_0)|+|f(\eta_0)|\bigr)
\]
\[
  \le \frac14\cdot2M=\frac12M,
\]
即 $0\le M\le\dfrac12M$. 所以 $M=0$，在 $\left[-\dfrac14,\dfrac14\right]$ 上 $f(x)\equiv0$，证毕.

\subsection*{八、Taylor 展开的唯一性问题}

\noindent\textbf{$\star$例 3.3.17}\quad 设 $f(x)$ 有连续的 $n$ 阶导数，$f(x)$ 在 $x=x_0$ 处有展开式：
\[
  f(x)=a_0+a_1(x-x_0)+a_2(x-x_0)^2+\cdots+a_n(x-x_0)^n+R_n(x), \tag{1}
\]
且余项 $R_n(x)$ 满足
\[
  \lim_{x\to x_0}\frac{R_n(x)}{(x-x_0)^n}=0, \tag{2}
\]
则必有
\[
  a_k=\frac{f^{(k)}(x_0)}{k!}\quad(k=0,1,2,\cdots,n), \tag{3}
\]
其中 $f^{(0)}(x)\equiv f(x)$.

\noindent\textbf{证}\quad 根据 Taylor 公式，$f(x)$ 在 $x=x_0$ 处可展开成
\[
  f(x)=\sum_{i=0}^n\frac{f^{(i)}(x_0)}{i!}(x-x_0)^i+o((x-x_0)^n). \tag{4}
\]
式 (1) 与式 (4) 联立即得
\[
  \sum_{i=0}^n a_i(x-x_0)^i+R_n(x)
  =\sum_{i=0}^n\frac{f^{(i)}(x_0)}{i!}(x-x_0)^i+o((x-x_0)^n).
\]
此式令 $x\to x_0$，取极限得 $a_0=f(x_0)$. 两边消去首项，再同时除以 $(x-x_0)$，然后令 $x\to x_0$，取极限又得 $a_1=f'(x_0)$. 继续这样下去，顺次可得式 (3).

\noindent\textbf{注}\quad $1^\circ$ 该例具有重要理论意义. 它表明：不论用何种途径、何种方式得到的形如式 (1) 的展开式，只要余项满足条件 (2)，则此展开式的系数必是唯一确定的，它们是式 (3) 给出的 Taylor 系数.

$2^\circ$ $x_0=0$ 时结论自然也成立. 由此可知，对于任何多项式 $P(x)=a_0+a_1x+\cdots+a_nx^n$ 而言，必有 $a_k=\dfrac{P^{(k)}(0)}{k!}\ (k=0,1,2,\cdots,n)$，$P^{(0)}(x)\equiv P(x)$.

\noindent\textbf{$\star$例 3.3.18}\quad 设 $P(x)$ 是 $n$ 次多项式，试证：
\[
  \sum_{k=0}^n\frac{P^{(k)}(0)}{(k+1)!}x^{k+1}
  =\sum_{k=0}^n(-1)^k\frac{P^{(k)}(x)}{(k+1)!}x^{k+1}, \tag{1}
\]
$P^{(0)}(x)=P(x)$.（中山大学）

\noindent\textbf{证 I}\quad $P(x)=a_0+a_1x+\cdots+a_nx^n$，则
\[
  a_k=\frac{P^{(k)}(0)}{k!}\quad(k=0,1,2,\cdots,n).
\]
令
\[
  f(x)=a_0x+a_1\frac{x^2}{2}+\cdots+a_n\frac{x^{n+1}}{n+1},
\]
则
\[
  f^{(k+1)}(x)=P^{(k)}(x)\quad(k=0,1,2,\cdots,n),
\]
\[
  f(x)=\sum_{k=0}^n a_k\frac{x^{k+1}}{k+1}
  =\sum_{k=0}^n\frac{P^{(k)}(0)}{(k+1)!}x^{k+1}. \tag{2}
\]
$\forall x_0$，将 $f$ 在 $x_0$ 处展开，有
\[
  f(x)=f(x_0)+\sum_{i=1}^{n+1}\frac{f^{(i)}(x_0)}{i!}(x-x_0)^i.
\]
令 $x=0$，得
\[
  0=f(0)=f(x_0)+\sum_{i=1}^{n+1}\frac{(-1)^if^{(i)}(x_0)}{i!}x_0^i.
\]
注意 $x_0$ 的任意性，不妨把 $x_0$ 改记为 $x$，移项则得
\[
  f(x)=-\sum_{i=1}^{n+1}\frac{(-1)^if^{(i)}(x)}{i!}x^i
  =\sum_{k=0}^n\frac{(-1)^k f^{(k+1)}(x)}{(k+1)!}x^{k+1}
\]
\[
  =\sum_{k=0}^n\frac{(-1)^kP^{(k)}(x)}{(k+1)!}x^{k+1}.
\]
与 (2) 比较，即得 (1).

\noindent\textbf{证 II}\quad (1) 式两端均为 $n+1$ 次多项式，只需证明 $(k+1)$ 次幂的系数相等即可 $(k=0,1,2,\cdots,n)$.

左端 $k+1$ 次幂的系数为 $\dfrac1{k+1}a_k$，右端 $k+1$ 次幂的系数为
\[
  a_k\left[1-\frac{k}{2!}+\frac{k(k-1)}{3!}-\cdots+(-1)^k\frac{k!}{(k+1)!}\right]
\]
\[
  =\frac{a_k}{k+1}\left[1-\left(1-C_{k+1}^1+C_{k+1}^2-\cdots+(-1)^{k+1}C_{k+1}^{k+1}\right)\right]
  =\frac{a_k}{k+1}\left[1-(1-1)^{k+1}\right]=\frac{a_k}{k+1}.
\]
故左端 = 右端，证毕.

\subsection*{九、符号 ``$O$'' 与 ``$o$'' 的含义和应用}

\noindent\textbf{要点}\quad 1）符号 ``$f(x)=O(1)$'' 表示在所讨论过程中，``$f(x)$ 是有界量''，即 $\exists M>0$，使得 $|f(x)|\le M$（在此过程中保持成立）；

符号 ``$o(1)$'' 代表在所讨论过程中，它是``无穷小量''. 例如：$\alpha=o(1)$，意指：（在所讨论过程中）$\alpha$ 是无穷小量.

2）$f(x)=O(g(x))$ 代表 $\dfrac{f(x)}{g(x)}=O(1)$，$f(x)=o(g(x))$ 代表 $\dfrac{f(x)}{g(x)}=o(1)$.

\noindent\textbf{new $\star$例 3.3.19}\quad 设函数 $f$ 在 $(0,+\infty)$ 上可微，且 $f'(x)=O(x)$（当 $x\to+\infty$ 时），试证：$f(x)=O(x^2)$（当 $x\to+\infty$ 时）.（中国科学院技术大学）

\noindent\textbf{证}\quad 已知 $\dfrac{f'(x)}x=O(1)\ (x\to+\infty)$，取 $x_0>0$，则 $\forall x>x_0$，有
\[
  f(x)=f(x_0)+f'(\xi)(x-x_0)\quad(x_0<\xi<x),\quad \left|\frac{\xi}{x}\right|\le1,\quad \left|\frac{x-x_0}{x}\right|\le1.
\]
故
\[
  \frac{f(x)}{x^2}=\frac{f(x_0)}{x^2}
  +\left(\frac{f'(\xi)}{\xi}\right)\left(\frac{\xi}{x}\right)\left(\frac{x-x_0}{x}\right)
  =o(1)+O(1)=O(1),
\]
即 $f(x)=O(x^2)$（当 $x\to+\infty$ 时）.

\noindent\textbf{new $\star$例 3.3.20}\quad 设 $x_n=f\left(\dfrac1{n^2}\right)+f\left(\dfrac2{n^2}\right)+\cdots+f\left(\dfrac n{n^2}\right)$，其中 $f(x)$ 在 $x=0$ 处有连续导数，且 $f(0)=0,\ f'(0)=1$. 试证：$\lim x_n$ 存在，并求极限值.（武汉大学）

\noindent\textbf{注}\quad $f(x)$ 在 $x=0$ 处有连续导数，意指：$f(x)$ 在 $x=0$ 处及其附近有导数，并且导函数 $f'(x)$ 至少在 $x=0$ 处连续.

\noindent\textbf{证 I}\quad 因 $f(0)=0,f'(0)=1$，应用 Taylor 公式，得
\[
  x_n=\sum_{i=1}^n f\left(\frac{i}{n^2}\right)
  =\sum_{i=1}^n\left[f(0)+f'(0)\frac{i}{n^2}+o\left(\frac{i}{n^2}\right)\right]
  =\sum_{i=1}^n\left[\frac{i}{n^2}+o\left(\frac{i}{n^2}\right)\right]
\]
\[
  =\sum_{i=1}^n\frac{i}{n^2}+\sum_{i=1}^n o\left(\frac{i}{n^2}\right)
  \xrightarrow[\text{见注}]{\frac{n(n+1)}{2n^2}+\frac{n(n+1)}{2n^2}o(1)}
  \frac12.
\]

\noindent\textbf{注}\quad 当 $n\to\infty$ 时，$o(1)$ 代表一个无穷小量，且 $o\left(\dfrac{i}{n^2}\right)=\dfrac{i}{n^2}o(1)$. 因此
\[
  \sum_{i=1}^n o\left(\frac{i}{n^2}\right)
  =\sum_{i=1}^n\frac{i}{n^2}o(1)
  =\frac1{n^2}\sum_{i=1}^n i\cdot o(1)
\]
\[
  =\frac1{n^2}[o(1)+2o(1)+3o(1)+\cdots+no(1)]
  =\frac1{n^2}\frac{n(n+1)}2o(1)\to0\quad(n\to\infty).
\]

\noindent\textbf{证 II}\quad
\[
  x_n\overset{f(0)=0}{=}\sum_{i=1}^n\left[f\left(\frac{i}{n^2}\right)-f(0)\right]
  \overset{\text{Lagrange 中值定理}}{=}
  \sum_{i=1}^n\frac{i}{n^2}f'(\xi_i)\quad\left(0<\xi_i<\frac{i}{n^2}\right)
\]
\[
  =\sum_{i=1}^n\frac{i}{n^2}\bigl(f'(\xi_i)-f'(0)\bigr)+\sum_{i=1}^n\frac{i}{n^2}f'(0)
  \xrightarrow[n\to\infty]{}0+\frac12=\frac12.
\]
（因为 $f'(x)$ 在原点连续，$\forall\varepsilon>0$，当 $n$ 充分大时，对 $|x|<\dfrac1n$，有 $|f'(x)-f'(0)|<\varepsilon$. 而 $0<\xi_i<\dfrac{i}{n^2}\le\dfrac1n$，故有 $|f'(\xi_i)-f'(0)|<\varepsilon$. 于是
\[
  0\le \left|\sum_{i=1}^n\frac{i}{n^2}\bigl(f'(\xi_i)-f'(0)\bigr)\right|
  \le\sum_{i=1}^n\frac{i}{n^2}|f'(\xi_i)-f'(0)|
  \le\sum_{i=1}^n\frac{i}{n^2}\varepsilon=\frac{n(n+1)}{2n^2}\varepsilon<\varepsilon.
\]
因此 $\lim\limits_{n\to\infty}\sum\limits_{i=1}^n\dfrac{i}{n^2}\bigl(f'(\xi_i)-f'(0)\bigr)=0$.）

\subsection*{单元练习 3.3}

\noindent Taylor 公式及其应用

\noindent\textbf{3.3.1}\quad 求 $e^{2x-x^2}$ 包含 $x^5$ 项的 Taylor 展开式.（北京大学）

\[
  \left(e^{2x-x^2}=1+2x+x^2-\frac23x^3-\frac56x^4-\frac1{15}x^5+o(x^5)\right)
\]

\noindent\textbf{提示}\quad $e^x=\sum\limits_{k=0}^5\dfrac1{k!}x^k+o(x^5)$ 中以 $2x-x^2$ 代入 $x$，展开 $(2x-x^2)^k\ (k=1,2,\cdots,5)$，合并同类项.

\noindent\textbf{$\star$3.3.2}\quad 设 $f(x)$ 在无穷区间 $(x_0,+\infty)$ 内可微分两次，$\lim\limits_{x\to x_0+0}f(x)=\lim\limits_{x\to+\infty}f(x)$ 存在且有限，试证：在区间 $(x_0,+\infty)$ 内至少有一点 $\xi$，满足 $f''(\xi)=0$.（山东大学）

\noindent\textbf{提示}\quad 若 $f''(x)$ 在 $(x_0,+\infty)$ 内变号，由导数的介值性（Darboux 定理，例 3.2.24），知 $\exists\xi$ 使得 $f''(\xi)=0$. 若 $f''(x)$ 不变号（恒大于 0，或恒小于 0），必导致矛盾.

再提示，$f(+\infty)=f(x_0+0)$. 根据例 3.2.10 可知 $\exists\eta\in(0,+\infty)$，使得 $f'(\eta)=0$. 若 $f''(x)$ 恒大于 0，则 $f'$ 严增，当 $x>\eta$ 时，有 $f'(x)>f'(\eta)=0$，从而
\[
  f(x)=f(x_1)+f'(\xi)(x-x_1)+\frac{f''(\zeta)}2(x-x_1)^2>f(x_1)+f'(x_1)(x-x_1)\to+\infty\quad(x\to+\infty).
\]
与 $\lim\limits_{x\to+\infty}f(x)$ 存在且有限相矛盾. 同理，$f''(x)$ 恒小于 0 也不可能.

\noindent\textbf{3.3.3}\quad 设 $f(x)$ 在 $[0,+\infty)$ 上具有连续二阶导数，又设 $f(0)>0,\ f'(0)<0,\ f''(x)<0\ (x\in[0,+\infty))$. 试证：在区间 $\left(0,-\dfrac{f(0)}{f'(0)}\right)$ 内至少有一个点 $\xi$，使 $f(\xi)=0$.（厦门大学）

\noindent\textbf{提示}\quad 利用 Taylor 公式，$\forall x>0$，$\exists\xi\in(0,x)$，使
\[
  f(x)=f(0)+f'(0)x+\frac{f''(\xi)}{2!}x^2.
\]
由此可知 $f$ 在区间 $\left[0,-\dfrac{f(0)}{f'(0)}\right]$ 两端点异号.

\noindent\textbf{3.3.4}\quad 设 $f(x)$ 在 $x_0$ 的邻域里存在四阶导数，且 $|f^{(4)}(x)|\le M$，试证：对于此邻域内异于 $x_0$ 的任何 $x$，均有
\[
  \left|f''(x_0)-\frac{f(x)-2f(x_0)+f(x')}{(x-x_0)^2}\right|
  \le \frac{M}{12}(x-x_0)^2,
\]
其中 $x'$ 与 $x$ 关于 $x_0$ 对称.

\noindent\textbf{提示}\quad 利用 Taylor 公式，将 $f(x)$ 与 $f(x')$ 在 $x_0$ 处展开到三次项，余项利用 Lagrange 形式，然后代入待证不等式左端. 注意：由于 $x'$ 与 $x$ 关于 $x_0$ 对称，含 $(x-x_0)$ 的项及 $(x'-x_0)$ 的项相互抵消，同理三次项也被抵消.

\noindent\textbf{3.3.5}\quad 设（1）$f(x),f'(x)$ 在 $[a,b]$ 上连续；（2）$f''(x)$ 在 $(a,b)$ 内存在；（3）$f(a)=f(b)=0$；（4）在 $(a,b)$ 内存在点 $c$，使 $f(c)>0$. 求证：在 $(a,b)$ 内存在 $\xi$，使 $f''(\xi)<0$.（四川大学）

\noindent\textbf{提示}\quad 由条件 (3)、(4) 知最大值必在 $(a,b)$ 内，设为 $x_0$，则 $f'(x_0)=0$（Fermat 定理），$f(x_0)>0$.

\noindent\textbf{再提示}\quad 由
\[
  f(b)=f(x_0)+\frac{f''(\xi)}2(b-x_0)^2,\quad \xi\in(x_0,b),
\]
即得 $f''(\xi)<0$.

\noindent\textbf{$\star$3.3.6}\quad 设 $f(x)$ 在 $[0,1]$ 上二阶可导，$f(0)=f(1)=0$，$\min\limits_{0\le x\le1}f(x)=-1$，求证：$\max\limits_{0\le x\le1}f''(x)\ge8$.（华中师范大学，湖南大学，北京师范大学）

\noindent\textbf{提示}\quad 参考例 3.3.5.

\noindent\textbf{3.3.7}\quad 设函数 $f(x)$ 在区间 $[0,1]$ 上有二阶导数，且当 $0\le x\le1$ 时，恒有 $|f(x)|\le a$，$|f''(x)|\le b$. 证明：当 $0<x<1$ 时，$|f'(x)|\le2a+\dfrac b2$.（数学一）

\noindent\textbf{提示}\quad 参考例 3.3.7.

\noindent\textbf{$\star$3.3.8}\quad 设 $f(x)$ 在 $[0,1]$ 上二次可微，$|f''(x)|\le M\ (0\le x\le1)$，$M>0$，$f(0)=f(1)=f\left(\dfrac12\right)=0$. 证明：$|f'(x)|<\dfrac M2\ (0\le x\le1)$.（华中科技大学）

\noindent\textbf{注}\quad 与例 3.3.8 比较，由于增加了 $f\left(\dfrac12\right)=0$ 的条件，结论里的 $|f'(x)|\le\dfrac M2$，改进为 $|f'(x)|<\dfrac M2$.

\noindent\textbf{提示}\quad 可分别在 $\left[0,\dfrac12\right]$，$\left[\dfrac12,1\right]$ 上应用例 3.3.7 的证法.

\noindent\textbf{再提示}\quad 设 $x\in\left[0,\dfrac12\right]$，则
\[
  0=f(0)=f(x)+f'(x)(0-x)+\frac{f''(\xi)}2(0-x)^2, \tag{1}
\]
\[
  0=f\left(\frac12\right)=f(x)+f'(x)\left(\frac12-x\right)
  +\frac{f''(\eta)}2\left(\frac12-x\right)^2. \tag{2}
\]
(2) $-$ (1) 得
\[
  |f'(x)|=\left|f''(\eta)\left(\frac12-x\right)^2-f''(\xi)x^2\right|
  \le M\left[\left(\frac12-x\right)^2+x^2\right]\le\frac M4<\frac M2.
\]
类似可证 $x\in\left[\dfrac12,1\right]$ 的情况.

\noindent\textbf{*3.3.9}\quad 设函数 $f(x),g(x),p(x)$ 有连续二阶导数，试求
\[
  \lim_{h\to0}\frac1{h^3}
  \begin{vmatrix}
  f(x)&g(x)&p(x)\\
  f(x+h)&g(x+h)&p(x+h)\\
  f(x+2h)&g(x+2h)&p(x+2h)
  \end{vmatrix}.
\]
（华中师范大学）

\noindent\textbf{提示}\quad
\[
  f(x+h)=f(x)+f'(x)h+\frac{f''(\xi_1)}2h^2\overset{\text{记}}{=}f_0+f_1h+f_2h^2,
\]
\[
  f(x+2h)=f(x)+f'(x)2h+\frac{f''(\xi_2)}2\,4h^2
  \overset{\text{记}}{=}f_0+f_1\,2h+f_3h^2.
\]
对 $g,p$ 有类似展开，代入原行列式，化简. 提取公因子，利用乘法分配律进行分解，将极限符号引入行列式内.

\noindent\textbf{再提示}\quad 将原式里的行列式记为 $H$，则
\[
  H=
  \begin{vmatrix}
  f_0&g_0&p_0\\
  f_0+f_1h+f_2h^2&g_0+g_1h+g_2h^2&p_0+p_1h+p_2h^2\\
  f_0+f_1\,2h+f_3h^2&g_0+g_1\,2h+g_3h^2&p_0+p_1\,2h+p_3h^2
  \end{vmatrix}.
\]
在第二行里分别减去第一行的对应元素，第三行也如此，然后提取公因数，得
\[
  H=h^2
  \begin{vmatrix}
  f_0&g_0&p_0\\
  f_1+f_2h&g_1+g_2h&p_1+p_2h\\
  2f_1+f_3h&2g_1+g_3h&2p_1+p_3h
  \end{vmatrix}\quad(\text{分解})
\]
\[
  =2h^3
  \begin{vmatrix}
  f_0&g_0&p_0\\
  f_2&g_2&p_2\\
  f_1&g_1&p_1
  \end{vmatrix}
  +h^3
  \begin{vmatrix}
  f_0&g_0&p_0\\
  f_1&g_1&p_1\\
  f_3&g_3&p_3
  \end{vmatrix}
  +h^4
  \begin{vmatrix}
  f_0&g_0&p_0\\
  f_2&g_2&p_2\\
  f_3&g_3&p_3
  \end{vmatrix}
\]
（为零的一项未写出）. 注意这里
\[
  f_2=\frac12f''(\xi_1)\to\frac12f''(x)\quad(h\to0\text{ 时})(\xi_1\text{ 介于 }x\text{ 与 }x+h\text{ 之间}),
\]
\[
  f_3=\frac42f''(\xi_2)\to2f''(x)\quad(h\to0\text{ 时})(\xi_2\text{ 在 }x\text{ 与 }x+2h\text{ 之间}).
\]
同理，$g_2\to\dfrac12g''(x),g_3\to2g''(x),p_2\to\dfrac12p''(x),p_3\to2p''(x)$（$g_i,p_i$ 的中间点应记作 $\zeta_i,\eta_i\ (i=1,2)$，以区别于 $\xi_i(i=1,2)$）.

最后可知
\[
  \text{原式}=\lim_{h\to0}\frac1{h^3}H=
  \begin{vmatrix}
  f(x)&g(x)&p(x)\\
  f'(x)&g'(x)&p'(x)\\
  f''(x)&g''(x)&p''(x)
  \end{vmatrix}.
\]

\noindent\textbf{$\star$3.3.10}\quad 若要使 $x\to0$ 时 $e^x-\dfrac{1+ax}{1+bx}$ 为尽可能高阶的无穷小，问数 $a,b$ 应取什么值？用 $x$ 的幂函数写出此时的等价无穷小.

\noindent\textbf{提示}\quad 用 Taylor 公式展开到 $x^3$ 次项.

\noindent\textbf{解}\quad
\[
  e^x=\sum_{k=0}^3\frac1{k!}x^k+o(x^3),
\]
\[
  \frac{1+ax}{1+bx}=\frac{1+bx+(a-b)x}{1+bx}
  =1+(a-b)x\left[\sum_{k=0}^2(-bx)^k+o(x^2)\right]
\]
\[
  =1+(a-b)\sum_{k=0}^2(-1)^kb^kx^{k+1}+o(x^3),
\]
\[
  e^x-\frac{1+ax}{1+bx}
  =[1-(a-b)]x+\left[\frac12+(a-b)b\right]x^2
  +\left(\frac1{3!}-ab^2+b^3\right)x^3+o(x^3).
\]
令 $x$ 的一、二次项系数为零，解得 $a=\dfrac12,\ b=-\dfrac12$. 此时
\[
  e^x-\frac{1+ax}{1+bx}=-\frac1{12}x^3+o(x^3)\sim-\frac1{12}x^3\quad(x\to0\text{ 时}).
\]

\noindent\textbf{new *3.3.11}\quad 设函数 $f(x)$ 在 $[a,b]$ 上有连续导数，在开区间 $(a,b)$ 内二阶可导，且
\[
  f'\left(\frac{a+b}{2}\right)=0,
\]
证明存在 $\xi\in(a,b)$，使得
\[
  |f''(\xi)|\ge\frac4{(b-a)^2}|f(b)-f(a)|. \tag{1}
\]
（南开大学）

\noindent\textbf{提示}\quad 可参考例 3.3.6 的证明方法.

\noindent\textbf{注}\quad 下一章习题 4.3.33 对该题有进一步讨论. 例 3.3.6 和习题 4.3.7 是类似问题.

\section*{§3.4\quad 不等式与凸函数}

不等式是数学分析中经常遇到而又比较困难的问题之一. 本节我们将用微分方法讨论不等式，以及与不等式密切相关的凸函数问题；在积分学里，我们将重新回到这些问题上来. 请参看 §4.3，§4.4.

\subsection*{$\star$一、不等式}

\subsubsection*{a. 利用单调性证明不等式}

\noindent\textbf{要点}\quad 若 $f'(x)\ge0$（或 $f'(x)>0$），则当 $x_1<x_2$ 时，有 $f(x_1)\le f(x_2)$（或 $f(x_1)<f(x_2)$）. 由此可获得不等式.

\noindent\textbf{例 3.4.1}\quad 证明：
\[
  \frac{|a+b|}{1+|a+b|}\le \frac{|a|}{1+|a|}+\frac{|b|}{1+|b|}.
\]

\noindent\textbf{证}\quad 记 $f(x)=\dfrac{x}{1+x}$，则 $f'(x)=\dfrac1{(1+x)^2}>0$，$f(x)=\dfrac{x}{1+x}\nearrow$. 于是，由 $|a+b|\le |a|+|b|$ 知
\[
  \frac{|a+b|}{1+|a+b|}\le\frac{|a|+|b|}{1+|a|+|b|}
  =\frac{|a|}{1+|a|+|b|}+\frac{|b|}{1+|a|+|b|}
\]
\[
  \le \frac{|a|}{1+|a|}+\frac{|b|}{1+|b|}.
\]

\subsubsection*{b. 利用微分中值定理证明不等式}

\noindent\textbf{要点}\quad $1^\circ$ 若 $f(x)$ 在 $[a,b]$ 上连续，在 $(a,b)$ 内可导，则
\[
  f(x)=f(a)+f'(\xi)(x-a)\quad(\xi\in(a,b)).
\]
故当 $f(a)=0$，$(a,b)$ 内 $f'(x)>0$ 时，有 $f(x)>0\ (\forall x\in(a,b])$.

$2^\circ$ 在上述条件下，有
\[
  \frac{f(b)-f(a)}{b-a}=f'(\xi),\quad \text{其中 }a<\xi<b.
\]
若 $f'(x)$ 严增，则
\[
  f'(a)<\frac{f(b)-f(a)}{b-a}<f'(b).
\]
以上原理，在证明不等式时经常采用.

\noindent\textbf{$\star$例 3.4.2}\quad 1）证明：当 $0<x<1$ 时，有 $x-\dfrac1x<2\ln x$；

2）设 $f(x)$ 在 $(0,+\infty)$ 内单调下降、可微，如果当 $x\in(0,+\infty)$ 时，$0<f(x)<|f'(x)|$ 成立，则当 $0<x<1$ 时，必有 $xf(x)>\dfrac1x f\left(\dfrac1x\right)$；（北京大学）

3）证明：当 $s>0$ 时，
\[
  \frac{n^{s+1}}{s+1}<1^s+2^s+\cdots+n^s<\frac{(n+1)^{s+1}}{s+1};
\]
（武汉理工大学）

\noindent\textbf{new 4）}\quad 若 $\lambda=\sum\limits_{k=1}^n\dfrac1k$，证明 $e^\lambda>n+1$.（中国科学院）

\noindent\textbf{证}\quad 1）只需证明 $f(x)\equiv x-\dfrac1x-2\ln x<0$. 事实上，$f(1)=0$. 而当 $0<x<1$ 时，
\[
  f'(x)=1+\frac1{x^2}-\frac2x=\frac{(x-1)^2}{x^2}>0.
\]
因此，$f(x)<0$（当 $0<x<1$ 时）.

\noindent\textbf{注}\quad 从证明看到，该不等式对于 $x>1$，有 $f(x)>0$ 成立. 若令 $x=\sqrt t$，以上两不等式还可统一写成
\[
  \frac{\ln t}{t-1}\le \frac1{\sqrt t}\quad(t>0,t\ne1).
\]
这曾是一道国外赛题.

2）目标在于证明 $(0,1)$ 内
\[
  \frac{f(1/x)}{f(x)}<x^2\quad\text{或}\quad
  \ln\frac{f(1/x)}{f(x)}<\ln x^2=2\ln x.
\]
事实上，因 $f\searrow$，$f'(x)<0$，有 $|f'(x)|=-f'(x)$，
\[
  \ln\frac{f(1/x)}{f(x)}
  =\ln f(1/x)-\ln f(x)
  \overset{\text{Lagrange 定理}}{=}
  \frac{f'(\xi)}{f(\xi)}\left(\frac1x-x\right).
\]
注意到 $0<f(x)<|f'(x)|=-f'(x)$，$\dfrac{f'(x)}{f(x)}<-1$，$\dfrac1x-x>0\ (0<x<1)$，再利用题 1）的结果，知
\[
  \ln\frac{f(1/x)}{f(x)}<x-\frac1x<2\ln x.
\]
证毕.

\noindent\textbf{另证}\quad 由 $0<f<|f'|=-f'$，$\dfrac{f'}f<-1$，$-(\ln f)'>1$，知
\[
  -\int_x^{1/x}(\ln f(t))'\,dt>\int_x^{1/x}dt,
\]
即
\[
  \ln\frac{f(x)}{f(1/x)}>\frac1x-x>-2\ln x=\ln x^{-2}.
\]
从而 $\dfrac{f(x)}{f(1/x)}>x^{-2}$，即 $xf(x)>\dfrac1x f(1/x)$.

3）在 $[k,k+1]$ 上对函数 $f(x)\equiv x^{s+1}$ 应用 Lagrange 公式，
\[
  (k+1)^{s+1}-k^{s+1}=(s+1)\xi^s,
\]
其中 $k<\xi<k+1$，从而 $k^s<\xi^s<(k+1)^s\ (k=0,1,2,\cdots)$. 因此
\[
  k^s<\frac{(k+1)^{s+1}-k^{s+1}}{s+1}<(k+1)^s\quad(k=0,1,2,\cdots). \tag{1}
\]
在左边的不等式里令 $k=0,1,\cdots,n$，$n+1$ 个不等式相加即得
\[
  \sum_{k=1}^n k^s<\frac{(n+1)^{s+1}}{s+1}; \tag{2}
\]
在右边不等式里令 $k=0,1,\cdots,n-1$，$n$ 个不等式相加即得
\[
  \frac{n^{s+1}}{s+1}<\sum_{k=0}^{n-1}(k+1)^s=\sum_{k=1}^n k^s. \tag{3}
\]
联结 (2)、(3) 两式，即为所求.

4）注意：$e^\lambda>n+1\Leftrightarrow \ln(n+1)<\lambda$. 而
\[
  \ln(n+1)=\sum_{k=1}^n[\ln(k+1)-\ln k]
  \overset{\text{Lagrange 定理}}{=}\sum_{k=1}^n\frac1{\xi_k}
  <\sum_{k=1}^n\frac1k\overset{\text{记}}{=}\lambda.
\]
即得 $e^\lambda>n+1$. 证毕.

\subsubsection*{c. 利用 Taylor 公式证明不等式}

\noindent\textbf{要点}\quad 若 $f(x)$ 在 $[a,b]$ 上有连续 $n$ 阶导数，且 $f(a)=f'(a)=\cdots=f^{(n-1)}(a)=0$，$f^{(n)}(x)>0$（当 $x\in(a,b]$ 时），则
\[
  f(x)=\frac{f^{(n)}(\xi)}{n!}(x-a)^n>0\quad(\text{当 }x\in(a,b]\text{ 时}).
\]
利用此原理，可以证明一些不等式.

\noindent\textbf{$\star$例 3.4.3}\quad 求证 $\dfrac{\tan x}{x}>\dfrac{x}{\sin x}$，$\forall x\in\left(0,\dfrac\pi2\right)$.（上海师范大学）

\noindent\textbf{证}\quad 原式等价于 $f(x)\equiv \sin x\cdot\tan x-x^2>0$，因 $f(0)=f'(0)=f''(0)=0$，
\[
  f'''(x)=\sin x(5\sec^2x-1)+6\sin^3x\sec^4x>0,
\]
故 $f(x)>0\left(x\in\left(0,\dfrac\pi2\right)\right)$. 原式获证.

\noindent\textbf{练习}\quad 试用此证法证明下章例 4.3.19 中的两个不等式.

\subsubsection*{d. 用求极值的方法证明不等式}

\noindent\textbf{要点}\quad 要证明 $f(x)\ge g(x)$，只要求函数 $F(x)\equiv f(x)-g(x)$ 的极值，证明 $\min F(x)\ge0$. 这是证明不等式的基本方法.

\noindent\textbf{例 3.4.4}\quad 设 $a>\ln2-1$ 为任一常数，试证：$x^2-2ax+1<e^x\ (x>0)$.

\noindent\textbf{证}\quad 问题是证明
\[
  f(x)\equiv e^x-x^2+2ax-1>0\quad(\text{当 }x>0\text{ 时}).
\]
因 $f(0)=0$，所以只要证明
\[
  f'(x)=e^x-2x+2a>0\quad(\text{当 }x>0\text{ 时})
  \quad\text{或}\quad \min_{x>0}f'(x)>0.
\]
令 $f''(x)=e^x-2=0$，得唯一稳定点 $x=\ln2$.

当 $x<\ln2$ 时，$f''(x)<0$；当 $x>\ln2$ 时，$f''(x)>0$. 所以
\[
  \min_{x>0}f'(x)=f'(\ln2)=2-2\ln2+2a=2(1-\ln2)+2a>0.
\]
证毕.

\noindent\textbf{$\star$例 3.4.5}\quad 设 $n$ 为自然数，试证：
\[
  e^{-t}-\left(1-\frac tn\right)^n\le \frac{t^2}{n}e^{-t}\quad(\text{当 }t\le n\text{ 时}).
\]
（吉林大学）

\noindent\textbf{证}\quad 原式等价于 $1-\left(1-\dfrac tn\right)^n e^t\le \dfrac{t^2}{n}$. 故只要证明
\[
  f(t)\equiv \frac{t^2}{n}-\left[1-\left(1-\frac tn\right)^n e^t\right]\ge0\quad(t\le n).
\]
\[
  f'(t)=\frac{2t}{n}+e^t\left[\left(1-\frac tn\right)^{n-1}(-1)+\left(1-\frac tn\right)^n\right]
  =\frac tn\left[2-e^t\left(1-\frac tn\right)^{n-1}\right].
\]
故用 $\xi$ 表示方程
\[
  2-e^t\left(1-\frac tn\right)^{n-1}=0 \tag{1}
\]
的根，则极值的可疑点为 $t=0,t=\xi$ 及 $t=n$. 但 $f(0)=0$，
\[
  f(\xi)=\frac{\xi^2}{n}-\left[1-\left(1-\frac{\xi}{n}\right)^n e^\xi\right]
  =\frac{\xi^2}{n}-\left[1-2\left(1-\frac{\xi}{n}\right)\right]\quad(\text{因式 }(1))
\]
\[
  =\left(1-\frac{\xi}{n}\right)^2+\frac{\xi^2}{n^2}(n-1)\ge0,
\]
\[
  f(n)=n-1\ge0,\quad f(-\infty)=+\infty.
\]
由此 $f(t)\ge \min_{t\le n}f(t)=f(0)=0\ (t\le n)$. 问题证毕.

\subsubsection*{e. 利用单调极限证明不等式}

\noindent\textbf{要点}\quad 若当 $x<b$ 时，$f(x)\nearrow$（或严 $\nearrow$），且当 $x\to b-0$ 时 $f(x)\to A$（以上条件今后简记作 $f(x)\nearrow A$（或 $f(x)$ 严 $\nearrow A$，当 $x\to b-0$ 时）），则
\[
  f(x)\le A\quad(\text{当 }x\le b\text{ 时})\quad
  (\text{或 }f(x)<A\quad(\text{当 }x<b\text{ 时})).
\]
对于递减或严格递减，也有类似结论. 利用这一原理可以证明一些不等式.

\noindent\textbf{例 3.4.6}\quad 证明：当 $x>0,t\le x$ 时，$e^{-t}-\left(1-\dfrac tx\right)^x\ge0$.（吉林大学）

\noindent\textbf{证}\quad 当 $t=0$ 或 $t=x$ 时，不等式自明. 只需证明 $x>0,t<x,t\ne0$ 的情况. 为此，只需证明当 $x\nearrow+\infty$ 时，$f(x)\equiv\left(1-\dfrac tx\right)^x\nearrow e^{-t}$ 即可. 事实上：

$1^\circ$ 当 $x>0,t\ne0,t<x$ 时，
\[
  [\ln f(x)]'=\left[\ln\left(1-\frac tx\right)^x\right]'_x
  =\left[x\ln\left(1-\frac tx\right)\right]'_x
\]
\[
  =\ln(x-t)-\ln x+\frac{t}{x-t}\quad(\text{应用 Lagrange 公式})
\]
\[
  =-\frac{t}{\xi}+\frac{t}{x-t}
  \quad(0<t<x\text{ 时 }0<x-t<\xi<x,\ t<0\text{ 时 }0<x<\xi<x-t)
\]
\[
  \ge -\frac{t}{x-t}+\frac{t}{x-t}=0.
\]

$2^\circ$
\[
  \lim_{x\to+\infty}\left(1-\frac tx\right)^x
  =\lim_{x\to+\infty}\left[\left(1-\frac tx\right)^{-\frac xt}\right]^{-t}
  =e^{-t}.
\]
故 $x\nearrow+\infty$ 时，$\left(1-\dfrac tx\right)^x\nearrow e^{-t}$. 证毕.

\noindent\textbf{$\star$例 3.4.7}\quad 证明：集合
\[
  A\equiv\left\{\alpha\mid \forall x>0,\ \left(1+\frac1x\right)^{x+\alpha}>e\right\}
\]
有最小值，并求最小值.（北京师范大学）

\noindent\textbf{证}\quad $1^\circ$ 不等式 $\left(1+\dfrac1x\right)^{x+\alpha}>e$ 等价于 $(x+\alpha)\ln\left(1+\dfrac1x\right)>1$，亦即
\[
  \alpha>\frac1{\ln\left(1+\frac1x\right)}-x\quad(\forall x>0).
\]
所以 $\alpha\in A$，等价于 $\alpha$ 为
\[
  f(x)\equiv\frac1{\ln\left(1+\frac1x\right)}-x\quad(x>0)
\]
的上界. 按确界的定义，即
\[
  \min A=\sup_{x>0}f(x).
\]

$2^\circ$ 由例 3.4.2 中 1）可知
\[
  f'(x)=\frac1{\ln^2\left(1+\frac1x\right)\cdot x(1+x)}-1>0,
\]
所以 $f(x)\nearrow$，
\[
  \sup_{x>0}f(x)=\lim_{x\to+\infty}f(x).
\]

$3^\circ$
\[
  f(x)=\frac1{\ln\left(1+\frac1x\right)}-x
  =x\left[\frac1{x\ln\left(1+\frac1x\right)}-1\right]
\]
\[
  =x\left[\frac1{x\left(\frac1x-\frac1{2x^2}+o\left(\frac1{x^2}\right)\right)}-1\right]
  =x\left[\frac1{1-\frac1{2x}+o\left(\frac1x\right)}-1\right]
\]
\[
  =x\left[1+\frac1{2x}+o\left(\frac1x\right)-1\right]
  =\frac12+o(1)\to\frac12\quad(\text{当 }x\to+\infty\text{ 时}).
\]
总之，$A$ 有最小值，$\min A=\dfrac12$.

\subsection*{二、凸函数}

凸函数是一重要的概念. 它在许多学科里有着重要的应用. 在研究生入学考试题中，也时有涉及. 考虑目前多数教材的情况，本段拟对凸函数最基本的内容作一概述，主要包括：凸函数几种不同的定义及它们的关系；凸函数各种等价描述；凸函数的性质及应用等.

\subsubsection*{a. 凸函数的几种定义以及它们的关系}

凸函数有几种不同的定义.

\noindent\textbf{$\star$定义 1}\quad 设 $f(x)$ 在区间 $I$ 上有定义，$f(x)$ 在 $I$ 上称为凸函数，当且仅当 $\forall x_1,x_2\in I,\ \forall\lambda\in(0,1)$，有
\[
  f(\lambda x_1+(1-\lambda)x_2)\le \lambda f(x_1)+(1-\lambda)f(x_2). \tag{A}
\]
若式 (A) 中的 ``$\le$'' 改成 ``$<$''，则是严格凸函数的定义；若 ``$\le$'' 改成 ``$\ge$'' 或 ``$>$''，则分别是凹函数与严格凹函数的定义. 由于凸与凹是对偶的概念，对一个有什么结论，对另一个亦有相应结论. 今后，只对凸函数进行论述.

\noindent\textbf{几何意义}\quad 设 $x_1<x_2$，因为 $\lambda\in(0,1)$，所以
\[
  x=\lambda x_1+(1-\lambda)x_2<\lambda x_2+(1-\lambda)x_2=x_2.
\]
同理可证 $x>x_1$. 因此 $x\in(x_1,x_2)$，且当 $\lambda$ 从 $0$ 连续变化到 $1$ 时 $x$ 也从 $x_2$ 连续变化到 $x_1$.

我们联结曲线 $y=f(x)\ (x\in I)$ 上两点 $A(x_1,f(x_1)),B(x_2,f(x_2))$，作弦 $AB$（如图 3.4.1），则 $AB$ 的方程为
\[
  \frac{y-f(x_2)}{f(x_1)-f(x_2)}
  =\frac{x-x_2}{x_1-x_2}.
\]
将此式的比值记为 $\lambda$，则可得 $AB$ 的参数方程：
\[
  \begin{cases}
  y=\lambda f(x_1)+(1-\lambda)f(x_2),\\
  x=\lambda x_1+(1-\lambda)x_2.
  \end{cases}
\]
这表明：在点 $x=\lambda x_1+(1-\lambda)x_2$ 处，弦 $AB$ 的高度为 $y=\lambda f(x_1)+(1-\lambda)f(x_2)$. 可见，不等式 (A) 说明在 $(x_1,x_2)$ 内每点 $x$ 处，曲线 $y=f(x)$ 的高度不超过弦 $AB$ 的高度. 换句话说，曲线在弦 $AB$ 以下. 对曲线上任意两点 $A,B$ 都是如此. 因此，凸函数意味着函数图形向下凸.

\begin{center}
  \includegraphics[width=0.42\textwidth]{3-4-1.png}

  图 3.4.1
\end{center}

现代数学多数采用这种定义. 本书也一律采用此定义. 除此定义之外，还有其他形式的定义.

\noindent\textbf{定义 2}\quad 设 $f(x)$ 在区间 $I$ 上有定义，$f(x)$ 称为 $I$ 上的凸函数，当且仅当 $\forall x_1,x_2\in I$，有
\[
  f\left(\frac{x_1+x_2}{2}\right)\le \frac{f(x_1)+f(x_2)}2. \tag{B}
\]
式 (B) 中的 ``$\le$'' 改为 ``$<$'' 便是严格凸的定义.

\noindent\textbf{定义 3}\quad 设 $f(x)$ 在区间 $I$ 上有定义，$f(x)$ 称为凸函数，当且仅当 $\forall x_1,x_2,\cdots,x_n\in I$，有
\[
  f\left(\frac{x_1+x_2+\cdots+x_n}{n}\right)
  \le \frac{f(x_1)+f(x_2)+\cdots+f(x_n)}{n}. \tag{C}
\]
式 (C) 中的 ``$\le$'' 改为 ``$<$'' 便是严格凸的定义.

\noindent\textbf{定义 4}\quad 设 $f(x)$ 在区间 $I$ 上有定义，当且仅当曲线 $y=f(x)$ 的切线恒保持在曲线以下，称 $f(x)$ 为凸函数. 若除切点之外，切线严格保持在曲线的下方，则称 $f(x)$ 为严格凸的.

下面我们将要证明定义 2、3 是等价的. 当 $f(x)$ 连续时定义 1、2、3 等价，当 $f(x)$ 处处可导时，定义 1、2、3、4 都等价（见定理 4 的推论 1）.

\noindent\textbf{定理 1}\quad 定义 2 与定义 3 等价.

\noindent\textbf{注}\quad 定义 3 $\Rightarrow$ 定义 2 明显. 只要证明：定义 2 $\Rightarrow$ 定义 3. 应用通常的数学归纳法，有一定的困难. 这里采用反向归纳法，其要点是：（1）证明命题对于自然数的某个子列成立（本定理证明式 (C) 对于 $n=2^k\ (k=1,2,\cdots)$ 皆成立）；（2）证明命题当 $n=k+1$ 成立时，必然对 $n=k$ 成立.

\noindent\textbf{证}\quad $1^\circ$ 由式 (B) 知式 (C) 当 $n=2$ 时成立. 现证 $n=4$ 时式 (C) 成立. 事实上，$\forall x_1,x_2,x_3,x_4\in I$，由式 (B)，有
\[
  f\left(\frac{x_1+x_2+x_3+x_4}{4}\right)
  =f\left(\frac{\frac{x_1+x_2}{2}+\frac{x_3+x_4}{2}}2\right)
  \le\frac{f\left(\frac{x_1+x_2}{2}\right)+f\left(\frac{x_3+x_4}{2}\right)}2
\]
\[
  \le\frac{f(x_1)+f(x_2)+f(x_3)+f(x_4)}4.
\]
此即式 (C) 对 $n=4$ 成立. 一般来说，对任一自然数 $k$，重复上面的方法，应用式 (B) $k$ 次，可知
\[
  f\left(\frac{x_1+x_2+\cdots+x_{2^k}}{2^k}\right)
  \le\frac{f(x_1)+f(x_2)+\cdots+f(x_{2^k})}{2^k}.
\]
这说明式 (C) 对一切 $n=2^k$ 皆成立.

$2^\circ$ （证明式 (C) 对 $n=k+1$ 成立时，必对 $n=k$ 也成立.）记
\[
  A=\frac{x_1+x_2+\cdots+x_k}{k},
\]
则 $x_1+x_2+\cdots+x_k=kA$，所以
\[
  A=\frac{x_1+x_2+\cdots+x_k+A}{k+1}.
\]
因式 (C) 对 $n=k+1$ 成立，故
\[
  f(A)=f\left(\frac{x_1+x_2+\cdots+x_k+A}{k+1}\right)
  \le\frac{f(x_1)+f(x_2)+\cdots+f(x_k)+f(A)}{k+1}.
\]
不等式两边同乘 $k+1$，减去 $f(A)$，最后除以 $k$. 注意 $A=\dfrac{x_1+x_2+\cdots+x_k}{k}$，得
\[
  f\left(\frac{x_1+x_2+\cdots+x_k}{k}\right)
  \le\frac{f(x_1)+f(x_2)+\cdots+f(x_k)}{k}.
\]
此式表示式 (C) 对 $n=k$ 成立. 证毕.

\noindent\textbf{定理 2}\quad 若 $f(x)$ 连续，则定义 1、2、3 等价.

\noindent\textbf{证}\quad $1^\circ$ （定义 1 $\Rightarrow$ 定义 2、3）在定义 1 中令 $\lambda=\dfrac12$，则由式 (A) 得
\[
  f\left(\frac{x_1+x_2}{2}\right)
  =f[\lambda x_1+(1-\lambda)x_2]
  \le\lambda f(x_1)+(1-\lambda)f(x_2)
  =\frac{f(x_1)+f(x_2)}2\quad(\forall x_1,x_2\in I).
\]
此式表明式 (B) 成立. 所以定义 1 蕴涵定义 2. 而定义 2、3 等价，故定义 1 也蕴涵定义 3.

$2^\circ$ （定义 2、3 $\Rightarrow$ 定义 1）设 $x_1,x_2\in I$ 为任意两点，为了证明式 (A) 对于任意实数 $\lambda\in(0,1)$ 成立. 我们先来证明：式 (A) 当 $\lambda$ 为有理数即 $\lambda=\dfrac mn\in(0,1)$（$m<n$ 为自然数）时成立. 事实上，
\[
  f(\lambda x_1+(1-\lambda)x_2)
  =f\left(\frac mn x_1+\left(1-\frac mn\right)x_2\right)
  =f\left(\frac{mx_1+(n-m)x_2}{n}\right)
\]
\[
  =f\left(\frac{\overbrace{x_1+\cdots+x_1}^{m\text{个}}+\overbrace{x_2+\cdots+x_2}^{n-m\text{个}}}{n}\right)
  \le\frac{\overbrace{f(x_1)+\cdots+f(x_1)}^{m\text{个}}+\overbrace{f(x_2)+\cdots+f(x_2)}^{n-m\text{个}}}{n}
\]
\[
  =\frac{mf(x_1)+(n-m)f(x_2)}{n}
  =\lambda f(x_1)+(1-\lambda)f(x_2),
\]
$\lambda$ 为有理数的情况获证.

若 $\lambda\in(0,1)$ 为无理数，则 $\exists$ 有理数 $\lambda_n\in(0,1)\ (n=1,2,\cdots)$ 使得 $\lambda_n\to\lambda$（当 $n\to\infty$ 时），从而由 $f(x)$ 的连续性：
\[
  f(\lambda x_1+(1-\lambda)x_2)
  =f\left(\lim_{n\to\infty}(\lambda_nx_1+(1-\lambda_n)x_2)\right)
  =\lim_{n\to\infty}f(\lambda_nx_1+(1-\lambda_n)x_2).
\]
对于有理数 $\lambda_n\in(0,1)$，上面已证明有
\[
  f(\lambda_nx_1+(1-\lambda_n)x_2)
  \le \lambda_n f(x_1)+(1-\lambda_n)f(x_2).
\]
此式中令 $n\to\infty$，取极限，联系上式，有
\[
  f(\lambda x_1+(1-\lambda)x_2)
  \le \lambda f(x_1)+(1-\lambda)f(x_2),
\]
即式 (A) 对任意无理数 $\lambda\in(0,1)$ 也成立. 这就证明了定义 2、3 蕴涵定义 1.

\noindent\textbf{注}\quad 上述证明里看到从定义 1 $\Rightarrow$ 定义 2、3 无须连续性，定义 2、3 $\Rightarrow$ 定义 1 才需要连续性. 可见定义 1 强于定义 2、3.

\subsubsection*{b. 凸函数的等价描述}

\noindent\textbf{定理 3}\quad 如图 3.4.2，设 $f(x)$ 在区间 $I$ 上有定义，则以下条件等价（其中各不等式要求 $\forall x_1,x_2,x_3\in I,\ x_1<x_2<x_3$ 保持成立）：

i）$f(x)$ 在 $I$ 上为凸函数；

ii）
\[
  \frac{f(x_2)-f(x_1)}{x_2-x_1}
  \le \frac{f(x_3)-f(x_1)}{x_3-x_1};
\]

iii）
\[
  \frac{f(x_3)-f(x_1)}{x_3-x_1}
  \le \frac{f(x_3)-f(x_2)}{x_3-x_2};
\]

iv）
\[
  \frac{f(x_2)-f(x_1)}{x_2-x_1}
  \le \frac{f(x_3)-f(x_2)}{x_3-x_2};
\]

v）曲线 $y=f(x)$ 上三点 $A(x_1,f(x_1)),B(x_2,f(x_2)),C(x_3,f(x_3))$ 所围的有向面积
\[
  \frac12
  \begin{vmatrix}
  1&x_1&f(x_1)\\
  1&x_2&f(x_2)\\
  1&x_3&f(x_3)
  \end{vmatrix}\ge0.
\]

\begin{center}
  \includegraphics[width=0.42\textwidth]{3-4-2.png}

  图 3.4.2
\end{center}

（对于严格凸函数，有类似结论，只要将 ``$\le$'' 改为 ``$<$'' 即可.）

\noindent\textbf{证}\quad $1^\circ$（证明 i）与 ii）等价.）对 $I$ 中任意 $x_1<x_3$，根据凸函数定义，条件 i）等价于
\[
  f(\lambda x_1+(1-\lambda)x_3)\le \lambda f(x_1)+(1-\lambda)f(x_3). \tag{1}
\]
另一方面，将条件 ii）中的不等式乘 $(x_2-x_1)$，移项变形，可知它等价于
\[
  f(x_2)\le \frac{x_2-x_1}{x_3-x_1}f(x_3)
  +\frac{x_3-x_2}{x_3-x_1}f(x_1). \tag{2}
\]
可见，$\forall x_2\in(x_1,x_3)$，令
\[
  \lambda=\frac{x_2-x_1}{x_3-x_1},\quad
  1-\lambda=\frac{x_3-x_2}{x_3-x_1},
\]
\[
  \lambda x_3+(1-\lambda)x_1
  =\frac{x_2-x_1}{x_3-x_1}x_3
  +\frac{x_3-x_2}{x_3-x_1}x_1=x_2.
\]
从而由 (1) 可推到 (2). 反之，$\forall\lambda\in(0,1)$，若令 $x_2=\lambda x_3+(1-\lambda)x_1$，则 $\lambda=\dfrac{x_2-x_1}{x_3-x_1}$，从而可由 (2) 推得 (1). 故 i）与 ii）等价.

$2^\circ$ 类似可证 iii）、iv）与 i）等价.

$3^\circ$（证明 ii）与 v）等价.）将 ii）中的不等式乘 $(x_2-x_1)(x_3-x_1)$ 并移项，可知 ii）中不等式等价于
\[
  (x_3-x_2)f(x_1)+(x_1-x_3)f(x_2)+(x_2-x_1)f(x_3)\ge0,
\]
此即
\[
  \begin{vmatrix}
  1&x_1&f(x_1)\\
  1&x_2&f(x_2)\\
  1&x_3&f(x_3)
  \end{vmatrix}\ge0.
\]

\noindent\textbf{推论 1}\quad 若 $f(x)$ 在区间 $I$ 上为凸函数，则 $I$ 上任意三点 $x_1<x_2<x_3$，有
\[
  \frac{f(x_2)-f(x_1)}{x_2-x_1}
  \le \frac{f(x_3)-f(x_1)}{x_3-x_1}
  \le \frac{f(x_3)-f(x_2)}{x_3-x_2}.
\]

\noindent\textbf{注}\quad 对曲线 $y=f(x)$ 上任意一条弦 $AB$，若用 $k_{AB}$ 表示弦 $AB$ 的斜率，那么此不等式的几何意义即为（如图 3.4.2）$k_{AB}\le k_{AC}\le k_{BC}$.

\noindent\textbf{推论 2}\quad 若 $f(x)$ 为区间 $I$ 上的凸函数，则 $\forall x_0\in I$，过 $x_0$ 的弦的斜率
\[
  k(x)=\frac{f(x)-f(x_0)}{x-x_0}
\]
是 $x$ 的增函数（若 $f$ 为严格凸函数，则 $k(x)$ 严 $\nearrow$）.

\noindent\textbf{推论 3}\quad 若 $f(x)$ 是区间 $I$ 上的凸函数，则 $I$ 上任意四点 $s<t<u<v$，有
\[
  \frac{f(t)-f(s)}{t-s}\le \frac{f(v)-f(u)}{v-u}.
\]
（若 $f$ 为严格凸函数，则 ``$\le$'' 可改为 ``$<$''.）

\noindent\textbf{推论 4}\quad 若 $f(x)$ 是区间 $I$ 上的凸函数，则对 $I$ 的任一内点 $x$，单侧导数 $f'_+(x),f'_-(x)$ 皆存在，皆为增函数，且
\[
  f'_-(x)\le f'_+(x)\quad(\forall x\in I^\circ).
\]
这里 $I^\circ$ 表示 $I$ 的全体内点组成之集合（若 $f$ 为严格凸的，则 $f'_+(x)$ 与 $f'_-(x)$ 为严格递增的）.

\noindent\textbf{证}\quad 因 $x$ 为内点，故 $\exists x_1,x_2\in I$，使得 $x_1<x<x_2$. 从而（利用推论 2）
\[
  \frac{f(x_1)-f(x)}{x_1-x}
  \le\frac{f(x_2)-f(x)}{x_2-x}.
\]
再由推论 2 所述，当 $x_1\nearrow x$ 时，$\dfrac{f(x_1)-f(x)}{x_1-x}\nearrow$. 故由单调有界原理知如下极限存在，且
\[
  f'_-(x)=\lim_{x_1\to x-}\frac{f(x_1)-f(x)}{x_1-x}
  \le\frac{f(x_2)-f(x)}{x_2-x}.
\]
同理，在此式中，令 $x_2\searrow x$，可知 $f'_+(x)$ 存在，且 $f'_-(x)\le f'_+(x)$. 最后在推论 3 中的不等式里取相应的极限，可知 $f'_+(x)$ 与 $f'_-(x)$ 皆为增函数.

\noindent\textbf{推论 5}\quad 若 $f(x)$ 在区间 $I$ 上为凸的，则 $f$ 在任一内点 $x\in I^\circ$ 上连续.

事实上，由推论 4 知 $f'_+(x),f'_-(x)$ 存在，所以 $f$ 在 $x$ 处左、右都连续.

\noindent\textbf{注}\quad 将凸函数在区间端点的值改成较大的值时，仍能保持其凸性；因此推论 4、5 的结论对于区间的端点，一般不成立.

\noindent\textbf{定理 4}\quad 设函数 $f(x)$ 在区间 $I$ 上有定义，则 $f(x)$ 为凸函数的充要条件是：$\forall x_0\in I^\circ$，$\exists$ 实数 $\alpha$，使得 $\forall x\in I$ 有
\[
  f(x)\ge \alpha(x-x_0)+f(x_0).
\]

\noindent\textbf{证}\quad $1^\circ$（必要性）因 $f(x)$ 为凸函数，由上面刚得的推论 4，知 $\forall x_0\in I^\circ$，$f'_-(x_0)$ 存在，且
\[
  \frac{f(x)-f(x_0)}{x-x_0}\nearrow f'_-(x_0).
\]
由此，任取一 $\alpha\ge f'_-(x_0)$，则当 $x<x_0$ 时有
\[
  f(x)\ge \alpha(x-x_0)+f(x_0).
\]
同理，取 $\alpha\le f'_+(x_0)$，则当 $x>x_0$ 时有
\[
  f(x)\ge \alpha(x-x_0)+f(x_0).
\]
因 $f'_-(x_0)\le f'_+(x_0)$，所以对任一 $\alpha:f'_-(x_0)\le\alpha\le f'_+(x_0)$，$\forall x\in I$，恒有
\[
  f(x)\ge \alpha(x-x_0)+f(x_0).
\]

$2^\circ$（充分性）设 $x_1<x_2<x_3$ 是区间 $I$ 上任意三点，由已知条件，对 $x_2$，存在 $\alpha$，使得
\[
  f(x)\ge \alpha(x-x_2)+f(x_2)\quad(\forall x\in I).
\]
由此，令 $x=x_1$ 和 $x=x_3$，可得
\[
  \frac{f(x_3)-f(x_2)}{x_3-x_2}\ge \alpha\ge \frac{f(x_1)-f(x_2)}{x_1-x_2}.
\]
根据定理 3，可知 $f(x)$ 为凸的.

\noindent\textbf{推论 1}\quad 设 $f(x)$ 在区间 $I$ 内部可导，则 $f(x)$ 在 $I$ 上为凸的充要条件是：$\forall x_0\in I^\circ$，有
\[
  f(x)\ge f'(x_0)(x-x_0)+f(x_0)\quad(\forall x\in I).
\]
由此可见，若 $f(x)$ 可导，则凸函数的定义 1、2、3、4 等价.

\noindent\textbf{推论 2}\quad 若 $f(x)$ 在区间 $I$ 上为凸的，则 $\forall x_0\in I^\circ$，在曲线 $y=f(x)$ 上一点 $(x_0,f(x_0))$ 可作一条直线
\[
  L:y=\alpha(x-x_0)+f(x_0),
\]
使曲线 $y=f(x)$ 位于直线 $L$ 上方.

若 $f$ 为严格凸函数，则除点 $(x_0,f(x_0))$ 之外曲线严格地在直线 $L$ 的上方. 这是著名分离性定理. 直线 $L$ 称为 $y=f(x)$ 的支撑.

\noindent\textbf{$\star$定理 5}\quad 设 $f(x)$ 在区间 $I$ 上有导数，则 $f(x)$ 在 $I$ 上为凸函数的充要条件是 $f'(x)\nearrow\ (x\in I)$.

\noindent\textbf{证}\quad $1^\circ$（充分性）$\forall x_1,x_2\in I$（不妨设 $x_1<x_2$）及 $\lambda\in(0,1)$，记 $x=\lambda x_1+(1-\lambda)x_2$，来证
\[
  f(x)=f[\lambda x_1+(1-\lambda)x_2]\le \lambda f(x_1)+(1-\lambda)f(x_2),
\]
即
\[
  f(x)-\lambda f(x_1)-(1-\lambda)f(x_2)\le0. \tag{1}
\]
注意 $f(x)=\lambda f(x)+(1-\lambda)f(x)$，式 (1) 等价于
\[
  \lambda[f(x)-f(x_1)]+(1-\lambda)[f(x)-f(x_2)]\le0. \tag{2}
\]
应用 Lagrange 定理，$\exists\xi,\eta:x_1<\xi<x<\eta<x_2$，使得
\[
  \lambda[f(x)-f(x_1)]+(1-\lambda)[f(x)-f(x_2)]
  =\lambda f'(\xi)(x-x_1)+(1-\lambda)f'(\eta)(x-x_2).
\]
但
\[
  x-x_1=[\lambda x_1+(1-\lambda)x_2]-x_1=(1-\lambda)(x_2-x_1),
\]
\[
  x-x_2=[\lambda x_1+(1-\lambda)x_2]-x_2=\lambda(x_1-x_2),
\]
故式 (2) 左端
\[
  \lambda[f(x)-f(x_1)]+(1-\lambda)[f(x)-f(x_2)]
\]
\[
  =\lambda f'(\xi)(1-\lambda)(x_2-x_1)
  +(1-\lambda)f'(\eta)\lambda(x_1-x_2)
\]
\[
  =\lambda(1-\lambda)(x_2-x_1)[f'(\xi)-f'(\eta)]. \tag{3}
\]
按已知条件 $f'(x)\nearrow$，得 $f'(\xi)\le f'(\eta)$. 从而上式 $\le0$，式 (2) 获证.

$2^\circ$（必要性）根据定理 3 的推论 4，$f'_+(x)$ 在 $I^\circ$ 内为递增的. 因 $f'(x)$ 存在，故 $f'(x)=f'_+(x)$ 亦在 $I^\circ$ 内为递增的. 若 $I$ 有右端点 $b$，按已知条件，$f$ 在 $b$ 点有左导数，$\forall x\in I^\circ$，易知
\[
  f'(x)=f'_+(x)\le \frac{f(x)-f(b)}{x-b}\le f'_-(b)=f'(b).
\]
同理，若 $I$ 有左端点 $a$，则 $f'(a)\le f'(x)$. 如此我们证明了：$f'(x)$ 在 $I$ 为递增的（不论 $I$ 为有限无限，开或闭，或半开半闭）.

\noindent\textbf{$\star$推论}\quad 若 $f(x)$ 在区间 $I$ 上有二阶导数，则 $f(x)$ 在 $I$ 上为凸函数的充要条件是 $f''(x)\ge0$.

从该定理上述证明过程中可以看出，若 $f(x)$ 在区间 $I$ 上可导，则 $f(x)$ 在 $I$ 上为严格凸函数的充要条件是 $f'(x)$ 在 $I$ 上严格递增. 从而若 $f(x)$ 在 $I$ 上有二阶导数，则 $f$ 为严格凸的充分条件是 $f''(x)>0$. 等价的充要条件是：$f''(x)\ge0$，且在任何子区间上 $f''(x)\ne0$（即曲线 $y=f(x)$ 的斜率 $\nearrow$，且不含直线段）.

\noindent\textbf{定理 6}\quad 若 $f(x)$ 在 $I$ 上有定义，则以下三条件等价：

i）$f(x)$ 在 $I$ 上为凸函数；

ii）$\forall q_i\ge0:q_1+q_2+\cdots+q_n=1,\ \forall x_1,x_2,\cdots,x_n\in I$，有
\[
  f(q_1x_1+q_2x_2+\cdots+q_nx_n)
  \le q_1f(x_1)+q_2f(x_2)+\cdots+q_nf(x_n);
\]

iii）$\forall p_i\ge0\ (i=1,2,\cdots,n)$ 不全为零，$\forall x_1,x_2,\cdots,x_n\in I$，有
\[
  f\left(\frac{p_1x_1+p_2x_2+\cdots+p_nx_n}{p_1+p_2+\cdots+p_n}\right)
  \le \frac{p_1f(x_1)+p_2f(x_2)+\cdots+p_nf(x_n)}
  {p_1+p_2+\cdots+p_n}.
\]

\noindent\textbf{证}\quad ii）$\Rightarrow$i），只要令 $n=2$ 即得.

i）$\Rightarrow$ii），用数学归纳法.

ii）$\Leftrightarrow$iii）明显.

以上全部结论，对于凹函数都有对偶结论，只要将函数值的不等式反向即得.

\subsubsection*{c. 凸函数的性质及应用}

利用凸性，很容易获得一些不等式.

\noindent\textbf{$\star$例 3.4.8}\quad 设 $x_i>0\ (i=1,2,\cdots,n)$，证明：
\[
  \frac{n}{\dfrac1{x_1}+\dfrac1{x_2}+\cdots+\dfrac1{x_n}}
  \le \sqrt[n]{x_1x_2\cdots x_n}
  \le \frac{x_1+x_2+\cdots+x_n}{n},
\]
其中等号当且仅当 $x_i$ 全部相等时成立.

\noindent\textbf{提示}\quad 原式第二个不等式是 $\{a_k\}_{k=1}^n$ 的均值不等式. 第一个不等号的两端同时取对数，可变为
\[
  -\ln\frac{\frac1{x_1}+\frac1{x_2}+\cdots+\frac1{x_n}}n
  \le \frac1n\left(-\ln\frac1{x_1}-\ln\frac1{x_2}-\cdots-\ln\frac1{x_n}\right).
\]
从而由函数 $f(x)=-\ln x$ 在 $(0,+\infty)$ 上的（严格）凸性可得；右边的不等式可直接由 $g(x)=\ln x$ 在 $(0,+\infty)$ 上的（严格）凹性得到.

\noindent\textbf{例 3.4.9}\quad 设函数 $f(x)$ 在区间 $I$ 为凸函数，试证：$f(x)$ 在 $I$ 的任一闭子区间上有界.（华中师范大学）

\noindent\textbf{证}\quad 设 $[a,b]\subset I$ 为任一闭子区间.

$1^\circ$（证明 $f(x)$ 在 $[a,b]$ 上有上界.）$\forall x\in[a,b]$，取 $\lambda=\dfrac{x-a}{b-a}\in[0,1]$，则
\[
  x=\lambda b+(1-\lambda)a.
\]
因 $f$ 为凸函数，所以
\[
  f(x)=f[\lambda b+(1-\lambda)a]\le \lambda f(b)+(1-\lambda)f(a)\le \lambda M+(1-\lambda)M=M,
\]
其中 $M=\max\{f(a),f(b)\}$. 即在 $[a,b]$ 上有上界 $M$.

$2^\circ$（证明 $f(x)$ 在 $[a,b]$ 上有下界.）记 $c=\dfrac{a+b}{2}$ 为 $a,b$ 的中点，则 $\forall x\in[a,b]$，有关于 $c$ 的对称点 $x'$. 因 $f$ 为凸函数，所以
\[
  f(c)\le \frac{f(x)+f(x')}{2}\le \frac12f(x)+\frac12M,
\]
从而 $f(x)\ge2f(c)-M\overset{\text{记}}{=}m$. 即 $m$ 为 $f(x)$ 在 $[a,b]$ 上的下界.

\noindent\textbf{例 3.4.10}\quad 设 $f(x)$ 为区间 $(a,b)$ 内的凸函数，试证：$f(x)$ 在 $I$ 的任一内闭区间 $[\alpha,\beta]\subset(a,b)$ 上满足 Lipschitz 条件.

\noindent\textbf{证}\quad 要证明 $f(x)$ 在 $[\alpha,\beta]$ 上满足 Lipschitz 条件，即要证明：$\exists L>0$，使得 $\forall x_1,x_2\in[\alpha,\beta]$，有
\[
  |f(x_1)-f(x_2)|\le L|x_1-x_2|. \tag{1}
\]
因为 $[\alpha,\beta]\subset(a,b)$，故可取 $h>0$ 充分小，使得 $[\alpha-h,\beta+h]\subset(a,b)$. 于是 $\forall x_1,x_2\in[\alpha,\beta]$，若 $x_1<x_2$，取 $x_3=x_2+h$，根据 $f$ 的凸性，
\[
  \frac{f(x_2)-f(x_1)}{x_2-x_1}
  \le \frac{f(x_3)-f(x_2)}{x_3-x_2}\le \frac{M-m}{h}
\]
（其中 $M,m$ 分别表示 $f(x)$ 在 $[\alpha-h,\beta+h]$ 上的上、下界），从而
\[
  f(x_2)-f(x_1)\le \frac{M-m}{h}|x_2-x_1|. \tag{2}
\]

若 $x_2<x_1$，可取 $x_3=x_2-h$，由 $f$ 的凸性，有
\[
  \frac{f(x_2)-f(x_3)}{x_2-x_3}
  \le \frac{f(x_1)-f(x_2)}{x_1-x_2}.
\]
从而
\[
  \frac{f(x_2)-f(x_1)}{x_1-x_2}
  \le \frac{f(x_3)-f(x_2)}{x_2-x_3}\le \frac{M-m}{h}.
\]
由此亦可推得式 (2) 成立.

若 $x_1=x_2$，则式 (2) 明显成立. 这就证明了式 (2) 对一切 $x_1,x_2\in[\alpha,\beta]$ 皆成立. 因此，式 (2) 当 $x_1$ 与 $x_2$ 交换位置时也应成立，故有
\[
  |f(x_2)-f(x_1)|\le \frac{M-m}{h}|x_2-x_1|.
\]
令 $L=\dfrac{M-m}{h}$，则式 (1) 获证.

\noindent\textbf{注}\quad 由本例可知，若 $f(x)$ 在 $(a,b)$ 内为凸的，则 $f(x)$ 在 $(a,b)$ 内连续. 但注意端点的情况不一样：即令 $f$ 在 $[a,b]$ 上为凸的，不能保证 $a,b$ 处连续，因为端点处 $f(a),f(b)$ 改为更大的数不会改变凸性.

\noindent\textbf{例 3.4.11}\quad 设 $f(0)=0,\ f(x)$ 在 $[0,+\infty)$ 上为非负的严格凸函数，$F(x)=\dfrac{f(x)}x$（当 $x>0$ 时）. 试证：$f(x),F(x)$ 为严格递增的.

\noindent\textbf{证}\quad 因 $f(x)$ 严格凸，$f(0)=0$，所以 $F(x)=\dfrac{f(x)}x=\dfrac{f(x)-f(0)}{x-0}$ 为严格递增的.

因为 $f(x)$ 非负，所以 $\forall x>0$，有 $f(x)\ge0=f(0)$. 若存在某点 $x_1>0$ 使得 $f(x_1)=0$，则在 $[0,x_1]$ 上有 $f(x)\equiv0$，与 $f(x)$ 为严格凸函数矛盾. 所以，$\forall x>0$，有 $f(x)>0$. 最后设 $x_2>x_1>0$，则
\[
  \frac{f(x_2)-f(x_1)}{x_2-x_1}
  >\frac{f(x_1)-f(0)}{x_1-0}=\frac{f(x_1)}{x_1}>0,
\]
知 $f(x)$ 为严格递增的（当 $x\in[0,+\infty)$ 时）.

\noindent\textbf{例 3.4.12}\quad 设 $f(x)$ 在 $[a,b]$ 上二次可微，对 $[a,b]$ 中每个 $x$，$f(x)$ 与 $f''(x)$ 同号或同时为零，又 $f(x)$ 在 $[a,b]$ 的任何子区间内不恒为零. 试证：$f(x)=0$ 在 $(a,b)$ 内如果有根，则必唯一.（广西师范大学）

\noindent\textbf{证}\quad （反证法）设 $f(x)=0$ 在 $(a,b)$ 内有两相异实根 $x_1,x_2\in(a,b)$（不妨设 $x_1<x_2$）. 因为 $f(x)$ 在 $[x_1,x_2]$ 上连续，在 $[x_1,x_2]$ 上有最大、最小值，而 $f(x_1)=f(x_2)=0$，所以最大、最小值至少有一个在内部达到（否则 $f(x)\equiv0$，与已知条件矛盾）. 例如在 $\xi\in(x_1,x_2)$ 处有最大值 $f(\xi)>0$. 根据连续函数局部保号性，必存在 $\xi$ 的某个邻域 $U=(\xi-\varepsilon,\xi+\varepsilon)$，使得在 $U$ 上恒有 $f(x)>0$（从而按已知条件，在 $U$ 上 $f''(x)>0$，$f$ 为凸函数），又因 $f(x_1)=f(x_2)=0$，邻域 $U$ 可取得足够大，以致在 $U$ 上 $f(x)\ne f(\xi)$，于是 $\exists\xi_1\in U$ 使得 $0<f(\xi_1)<f(\xi)$. 记 $\xi_1$ 关于 $\xi$ 的对称点为 $\xi_2$，则 $\xi_2\in(\xi-\varepsilon,\xi+\varepsilon)$，有 $0<f(\xi_2)\le f(\xi)$. 从而
\[
  \frac{f(\xi_1)+f(\xi_2)}2<f(\xi)
  =f\left(\frac{\xi_1+\xi_2}2\right),
\]
与凸性矛盾.

对于 $(x_1,x_2)$ 内部达到负的最小值，可以类似证明.

\noindent\textbf{例 3.4.13}\quad 设 $f(x)$ 在区间 $(a,b)$ 内为凸函数，并且有界，试证极限 $\lim\limits_{x\to a+}f(x)$ 与 $\lim\limits_{x\to b-}f(x)$ 存在.

\noindent\textbf{证}\quad 设当 $x\in(a,b)$ 时 $f(x)\le M$，$x>x_1>x_0$ 为 $(a,b)$ 内任意三点. 根据 $f(x)$ 的凸性，当 $x\nearrow b$ 时，$\dfrac{f(x)-f(x_0)}{x-x_0}\nearrow$.

又因为
\[
  \frac{f(x)-f(x_0)}{x-x_0}
  \le \frac{M-f(x_0)}{x_1-x_0}\quad(\forall x>x_1>x_0),
\]
根据单调有界原理，有极限
\[
  \lim_{x\to b-}\frac{f(x)-f(x_0)}{x-x_0}=A.
\]
从而
\[
  \lim_{x\to b-}f(x)
  =\lim_{x\to b-}\left[(x-x_0)\cdot\frac{f(x)-f(x_0)}{x-x_0}+f(x_0)\right]
  =A(b-x_0)+f(x_0)
\]
亦存在. 类似可证 $\lim\limits_{x\to a+}f(x)$ 存在.

关于凸函数的习题见本节习题 3.4.19 至 3.4.23. 关于凸函数的积分性质，见下章例 4.3.26 至例 4.3.30 以及习题 4.3.26 至习题 4.3.28.

\noindent\textbf{$\star$例 3.4.14}\quad 1）是否存在 $\mathbb R\to\mathbb R$ 的连续可导函数 $f$，满足：$f(x)>0$ 且 $f'(x)=f(f(x))$？（北京大学，华中师范大学）

2）若 $f(x)>0$ 改为 $f(x)\ge0$，结论如何？（华中师范大学）

\noindent\textbf{解}\quad 1）不存在！下面证之.

\noindent\textbf{证法 I}\quad $f'(x)=f(f(x))$，亦即 $\dfrac{d}{dx}f(x)=f(f(x))$，也即 $\dfrac{d}{dx}(\cdot)=f(\cdot)$，故
\[
  f''(x)=f(f(x))=\frac{d}{dx}\left(\frac{d}{dx}(x)\right)=0. \tag{1}
\]
但另一方面，$f(x)>0$（$f$ 在每点上都为正），知 $f'(x)=f(f(x))=f(y)>0\ (y=f(x))$. 与式 (1) 矛盾. 说明 $f(x)>0$ 和 $f'(x)=f(f(x))$ 不可能同时满足，故所求的 $f$ 不存在.

\noindent\textbf{证法 II}\quad （反证法）假设 $\forall x\in\mathbb R$，$f(x)>0$ 且 $f'(x)=f(f(x))>0$，则 $f(x)$ 严格递增. 于是利用 Lagrange 定理，存在 $\xi\in(-1,0)$ 使得
\[
  f(0)-f(-1)=f'(\xi)=f(f(\xi))>f(0)\quad(\text{因为 }f(\xi)>0).
\]
于是，$f(-1)<0$，与条件 $f(x)>0\ (\forall x\in\mathbb R)$ 矛盾.

\noindent\textbf{证法 III}\quad 由已知条件可知：$f(x)$ 有二阶（乃至任意阶的）导数，且 $f''(x)>0$，因此 $f(x)$ 为严格凸函数（曲线严格在切线之上方）. 特别地，$f(0)$ 在如下切线的上方（该切线是 $f(x)$ 在点 $x=-1$ 处的切线）：
\[
  y=f(-1)+f'(-1)\cdot[x-(-1)],
\]
即
\[
  f(0)>f(-1)+f'(-1)\cdot[0-(-1)]
  =f(-1)+f(f(-1)),
\]
\[
  f(0)-f(-1)>f(f(-1))>f(0)\quad(\text{因为 }f(-1)>0).
\]
由此得 $f(-1)<0$. 与“处处 $f(x)>0$”矛盾！

\noindent\textbf{证法 IV}\quad （未学积分的读者，此段暂缓阅读.）由 $f'(x)=f(f(x))>0$ 可知 $f$ 严格 $\nearrow$，故 $f'(x)=f(f(x))\overset{f(x)>0}{>}f(0)$，即 $f'(x)>f(0)$. 两端同时积分，
\[
  \int_{-1}^0 f'(x)\,dx>\int_{-1}^0 f(0)\,dx,
\]
得 $f(0)-f(-1)>f(0)$. 故 $f(-1)<0$. 这与“处处 $f(x)>0$”矛盾.

2）若 1）中的条件 $f(x)>0$ 改为 $f(x)\ge0$，则问题只有唯一解，即 $f(x)\equiv0$.

事实上，若 $f$ 为零函数（$f(x)\equiv0$），则 $f$ 满足一切条件，显然它是问题的解，剩下只需证明：若 $f(x)\not\equiv0,\ f(x)\ge0$，则问题无解.

因此时至少存在一点使得 $f(x)>0$，由于连续保号性，至少有个小区间，其中 $f(x)>0$. 再利用 1）中的证 I 即可找出矛盾.

\noindent\textbf{$\star\star$例 3.4.15}\quad 已知函数 $f(x)$ 在 $(-\infty,+\infty)$ 内有二阶连续导数，满足：

1）$\lim\limits_{x\to\infty}(f(x)-|x|)=0$；

2）存在 $x_0$，使得 $f(x_0)\le0$.

证明：$f''(x)$ 在 $(-\infty,+\infty)$ 内变号.（北京大学）

\noindent\textbf{证}\quad $1^\circ$ $\lim\limits_{x\to\infty}(f(x)-|x|)=0$ 表明：$x\to+\infty$ 时，$f(x)\to+\infty$，且以 $y=x$ 为渐近线；$x\to-\infty$ 时，$f(x)\to-\infty$，且以 $y=-x$ 为渐近线.

$2^\circ$（反证法）若 $f''(x)$ 不变号，据 $1^\circ$，只能 $f''(x)\ge0,\ f'(x)\nearrow$，故 $f(x)$ 为凸函数（凸向下）. 已知：$\exists x_0$ 使得 $f(x_0)\le0$. 下面分三种情况分别讨论：1）$x_0>0$；2）$x_0<0$；3）$x_0=0$.

1）若 $\exists x_0>0$，使得 $f(x_0)\le0$，那么 $f(x_0)\le0<x_0$. 此式表明：在点 $x=x_0$ 处，曲线 $y=f(x)$ 位于渐近线 $y=x$ 的下方，且竖直距离为
\[
  \varepsilon_0=x_0-f(x_0)>0. \tag{1}
\]
因为 $\lim\limits_{x\to+\infty}(f(x)-x)=0$，$\exists\Delta>0$，当 $x_1>\Delta$ 时，有
\[
  |f(x_1)-x_1|<\varepsilon_0\overset{\text{式 (1)}}{=}x_0-f(x_0). \tag{2}
\]
于是
\[
  x_1-f(x_1)\le |f(x_1)-x_1|<\varepsilon_0\overset{\text{式 (1)}}{=}x_0-f(x_0),
\]
即 $x_1-x_0<f(x_1)-f(x_0)$. 故 $\exists\xi_1\in(x_0,x_1)$ 使得
\[
  f'(\xi_1)\overset{\text{Lagrange 定理}}{=}\frac{f(x_1)-f(x_0)}{x_1-x_0}>1.
\]
记 $k=f'(\xi_1)$，则 $y=k(x-\xi_1)+f(\xi_1)$ 是 $f(x)$ 在 $x=\xi_1$ 处的切线. 根据凸函数性质，曲线 $f(x)$ 总在切线上方，故
\[
  f(x)\ge k(x-\xi_1)+f(\xi_1). \tag{3}
\]
由此得
\[
  f(x)-x\ge (k-1)x-k\xi_1+f(\xi_1)\to+\infty\quad(x\to+\infty),
\]
与已知条件 1）矛盾.

（注\quad 式 (3) 也可利用 Taylor 公式得到：因 $f''(\eta)\ge0,\ f'(\xi_1)=k$，故
\[
  f(x)=f(\xi_1)+f'(\xi_1)(x-\xi_1)+\frac{f''(\eta)}2(x-\xi_1)^2\ge k(x-\xi_1)+f(\xi_1).
\]
）

2）若 $\exists x_0<0$，使得 $f(x_0)\le0$，类似可证：当 $x\to-\infty$ 时，也与条件 1）矛盾.

3）若 $\exists x_0=0$，使得 $f(x_0)\le0$.

i）当 $f(x_0)<0$ 时，记 $\varepsilon_0=x_0-f(x_0)>0$，用 1）的方法，推出与条件 1）矛盾.

ii）如果 $f(x_0)=f(0)=0$: $f(x)=|x|$ 不会发生，否则在 $x=0$ 处 $f(x)$ 不可导（与题设矛盾）. 若 $f(x)\ne |x|$，则 $\exists x_1\ne0$ 使得 $f(x_1)>x_1$，或 $f(x_1)<x_1$.

若 $f(x_1)<x_1$，则回到已讨论的情况 1）和 2）.

若 $f(x_1)>x_1>0$，因为 $f(0)=0$，可在 $[0,x_1]$ 上对 $f(x)$ 应用 Lagrange 定理，找出 $\xi_1:f'(\xi_1)>1$，如上用 Taylor 公式，导出矛盾.

总之，$f''(x)$ 不变号是错的. 证毕.

\noindent\textbf{注}\quad $f''(x)$ 变号意味着 $f$ 有拐点. 因此本题实际上给出了函数 $f(x)$ 在 $(-\infty,+\infty)$ 上存在拐点的一个充分条件.

\noindent\textbf{$\star\star$例 3.4.16}\quad 设 $f(x)$ 和 $g(x)$ 在区间 $(a,b)$ 内有定义，且 $\forall x,x_0\in(a,b)$ 有
\[
  f(x)-f(x_0)\ge g(x_0)(x-x_0), \tag{1}
\]
试证：任何 $x_0\in(a,b)$，$f(x)$ 在 $x_0$ 处连续，并有左、右单侧导数.（南开大学）

\noindent\textbf{证 I}\quad 若记 $k=g(x_0)$ 则式 (1) 可变为
\[
  f(x)\ge k(x-x_0)+f(x_0)\quad(\forall x,x_0\in(a,b)).
\]
由本节定理 4 知，$f(x)$ 是 $(a,b)$ 上的凸函数. 又由本节定理 3 的推论 4 知，凸函数 $f(x)$ 在每点左、右单侧导数都存在，且 $f'_-(x)\le f'_+(x)$（虽然可以不相等，但能断言单侧导数存在）. 对 $\forall x\in(a,b)$，$f(x)$ 左、右都连续，因此 $f(x)$ 在 $(a,b)$ 内处处连续.

\noindent\textbf{证 II}\quad 由式 (1) 得
\[
  g(x_3)\ge \frac{f(x_3)-f(x_2)}{x_3-x_2}
  \ge g(x_2)\ge \frac{f(x_2)-f(x_1)}{x_2-x_1}
  \ge g(x_1)\quad(\forall x_1<x_2<x_3).
\]
用分数的和比不等式：$\dfrac ab>\dfrac cd\Rightarrow\dfrac ab>\dfrac{a+c}{b+d}>\dfrac cd$，得
\[
  g(x_3)\ge \frac{f(x_3)-f(x_2)}{x_3-x_2}
  \ge \frac{f(x_3)-f(x_1)}{x_3-x_1}
  \ge \frac{f(x_2)-f(x_1)}{x_2-x_1}
  \ge g(x_1).
\]
其中右边两个不等式表明：当 $x_3\to x_1+$ 时，$\dfrac{f(x_3)-f(x_1)}{x_3-x_1}$ 递减有下界. 故 $f'_+(x_1)$ 存在（$\forall x_1\in(a,b)$）.

同理（用左边两不等式）可证 $f'_-(x_3)$ 存在. 注意：$x_1$ 为任意点，知 $(a,b)$ 里的每点两侧导数都存在，因而 $f$ 在每点左、右都连续，故 $f$ 在 $(a,b)$ 内连续. 证毕.

\subsection*{单元练习 3.4}

\noindent\textbf{3.4.1}\quad 1）设 $b>a>e$，证明：$a^b>b^a$；（数学一）

2）比较 $\pi^e$ 与 $e^\pi$ 的大小.（复旦大学）

\noindent\textbf{提示}\quad 1）$\left(\dfrac{\ln x}{x}\right)'<0\Rightarrow\dfrac{\ln x}{x}\searrow\Rightarrow\dfrac{\ln a}{a}>\dfrac{\ln b}{b}$.

\noindent\textbf{3.4.2}\quad 设 $0<b\le a$，证明：
\[
  \frac{a-b}{a}\le\ln\frac ab\le\frac{a-b}{b}.
\]
（兰州大学，四川大学，华中科技大学等）

\noindent\textbf{3.4.3}\quad 证明：$2^n\ge1+n\sqrt{2^{n-1}}\ (n\ge1\text{ 为自然数})$.（北京邮电大学）

\noindent\textbf{提示}\quad 只要证 $f(x)=2^x-1-x\cdot\dfrac{2^{x/2}}{\sqrt2}\ge0$（当 $x\ge1$ 时）. 而 $f(1)=0,\ f'(x)>0$.

\noindent\textbf{3.4.4}\quad 设 $f(x)$ 定义在 $[0,c]$ 上，$f'(x)$ 存在且单调下降，$f(0)=0$，请用 Lagrange 定理证明：对于 $0\le a\le b\le a+b\le c$，恒有 $f(a+b)\le f(a)+f(b)$.（复旦大学）

\noindent\textbf{提示}\quad 原式等价于
\[
  \frac{f(a+b)-f(b)}{(a+b)-b}\le\frac{f(a)-f(0)}{a-0}.
\]

\noindent\textbf{3.4.5}\quad 试证：当 $x>0$ 时，$(x^2-1)\ln x\ge(x-1)^2$.（数学一）

\noindent\textbf{提示}\quad $x\ne1$ 时，
\[
  (x^2-1)\ln x-(x-1)^2
  =(x-1)^2\left[(x+1)\frac{\ln x-\ln1}{x-1}-1\right]\ge0
  \quad\left(\frac{x+1}{\xi}>1\right).
\]

\noindent\textbf{3.4.6}\quad 设在 $[0,1]$ 上 $f''(x)>0$，则 $f'(0),f'(1),f(1)-f(0)$ 的大小顺序是（\quad）.（数学一）

（A）$f(1)>f'(0)>f(1)-f(0)$ \qquad
（B）$f'(1)>f(1)-f(0)>f'(0)$

（C）$f(1)-f(0)>f'(1)>f'(0)$ \qquad
（D）$f'(1)>f(0)-f(1)>f'(0)$

\noindent\textbf{提示}\quad $f'\nearrow\Rightarrow f(1)-f(0)=f'(\xi)
\begin{cases}
<f'(1),\\
>f'(0),
\end{cases}\quad0<\xi<1.$

\noindent\textbf{3.4.7}\quad 已知在 $x>-1$ 里定义的可微函数 $f(x)$ 满足 $f'(x)+f(x)-\dfrac1{x+1}\int_0^x f(t)\,dt=0$ 和 $f(0)=1$.

1）求 $f'(x)$；

2）证明：$f(x)$ 在 $x\ge0$ 满足 $e^{-x}\le f(x)\le1$.（大连理工大学）

\noindent\textbf{提示}\quad 1）将原式求导并与原式联立可得出关于 $y=f'(x)$ 的微分方程 $y'+\dfrac{x+2}{x+1}y=0$.

2）左边不等式可考虑 $F(x)=f(x)-e^{-x}
\begin{cases}
=0, & \text{当 }x=0\text{ 时},\\
\nearrow, & \text{当 }x\ge0\text{ 时}.
\end{cases}$

\noindent\textbf{3.4.8}\quad 已知 $x<0$，求证：
\[
  \frac1x+\frac1{\ln(1-x)}<1.
\]
（中国地质大学）

\noindent\textbf{提示}\quad 宜令 $x=-t$，原式等价于 $f(t)=t\ln(1+t)+\ln(1+t)-t>0$，而 $f(0)=f'(0)=0$，$f''(t)=\dfrac1{1+t}>0$.

\noindent\textbf{3.4.9}\quad 证明：
\[
  \frac{e^a-e^b}{a-b}<\frac{e^a+e^b}2\quad(a\ne b).
\]
（国外赛题）

\noindent\textbf{提示}\quad 设 $a<b$，只需证明
\[
  f(x)=e^x-e^a-\frac{e^a+e^b}2(x-a)<0,
\]
而 $f(a)=f'(a)=0,\ f''(x)<0$（当 $x>a$ 时）.

\noindent\textbf{$\star$3.4.10}\quad 证明：对自然数 $n$，有
\[
  0<\frac{e}{\left(1+\dfrac1n\right)^n}-1<\frac1{2n}.
\]

\noindent\textbf{提示}\quad 左边不等式由 $\left(1+\dfrac1n\right)^n\nearrow e$ 自明.

右边不等式等价于
\[
  \ln\left(1+\frac1n\right)+\frac1n\ln\left(1+\frac1{2n}\right)-\frac1n>0.
\]
只要证 $f(x)=\ln(1+x)+x\ln\left(1+\dfrac12x\right)-x>0$，而 $f(0)=f'(0)=0,\ f''(x)>0$（当 $x>0$ 时）.

\noindent\textbf{3.4.11}\quad $x>1,\ r>1$，证明：
\[
  x^r>1+r(x-1)+\frac12r(r-1)\left(\frac{x-1}{x}\right)^2.
\]

\noindent\textbf{提示}\quad 对函数 $f(x)=x^r=[1+(x-1)]^r$ 用 Taylor 公式展开至 1 次项，并用 Lagrange 余项.

\noindent\textbf{再提示}\quad
\[
  x^r=1+r(x-1)+\frac12r(r-1)\frac{\xi^r}{\xi^2}(x-1)^2\quad(1<\xi<x).
\]

\noindent\textbf{$\star$3.4.12}\quad 设 $g(x)$ 在 $[a,b]$ 内连续，在 $(a,b)$ 内二阶可导，且 $|g''(x)|\ge m>0\ (m\text{ 为常数})$，又 $g(a)=g(b)=0$. 证明：
\[
  \max_{a\le x\le b}|g(x)|\ge\frac m8(b-a)^2.
\]
（北京师范大学）

\noindent\textbf{提示}\quad 可以看出最大值必在内部某点达到，记此点为 $x_0$，对 $f(b)$（或 $f(a)$）在 $x=x_0$ 处应用 Taylor 公式. 注意 $|g(x_0)|=\max\limits_{a\le x\le b}|g(x)|$.

\noindent\textbf{再提示}\quad
\[
  g(x)=g(x_0)+g'(x_0)(x-x_0)+\frac12g''(\xi)(x-x_0)^2.
\]
令 $x=a$ 或 $b$ 得
\[
  0=g(x_0)+\frac12g''(\xi)(x-x_0)^2.
\]
取离 $x_0$ 较远的端点，知
\[
  |g(x_0)|\ge\frac12|g''(\xi)|\left(\frac{b-a}2\right)^2
  \ge\frac m8(b-a)^2.
\]
此即欲证之式.

\noindent\textbf{$\star$3.4.13}\quad 证明：
\[
  \left(\frac{\sin x}{x}\right)^3\ge \cos x\quad\left(0<|x|<\frac\pi2\right).
\]
（国外赛题）

\noindent\textbf{提示}\quad 宜将 $x^3$ 单独作一项，如变形为
\[
  f(x)\equiv\sin^3x\cdot(\cos x)^{-1}-x^3\ge0
\]
（可设 $0<x<\dfrac\pi2$，因只需证明 $x>0$ 的情况）. 然后证 $f(0)=f'(0)=f''(0)=f'''(0)=0,\ f^{(4)}(x)>0$.

\noindent\textbf{注意}\quad 此题计算虽很烦琐，但很典型，更能体现 Taylor 公式的意义.

\noindent\textbf{小结}\quad 以上各题主要练习用单调性、中值定理、Taylor 公式证明不等式.

\noindent\textbf{$\star$3.4.14}\quad 设 $0<x<y<1$ 或 $1<x<y$，证明：
\[
  \frac yx>\frac{y^x}{x^y}.
\]
（中国科学院）

\noindent\textbf{提示}\quad 原式 $\Leftrightarrow x^{y-1}>y^{x-1}\Leftrightarrow\dfrac{y-1}{\ln y}>\dfrac{x-1}{\ln x}$. 因此只要证 $f(x)\equiv\dfrac{x-1}{\ln x}\nearrow$ 或 $f'(x)\ge0$.

\noindent\textbf{再提示}\quad
\[
  f'(x)=\frac{x\ln x-(x-1)}{x\ln^2x},
\]
只需证明 $g(x)=x\ln x-(x-1)>0$. 但是
\[
g'(x)=
\begin{cases}
<0, & x<1,\\
0, & x=1,\\
>0, & x>1,
\end{cases}
\]
故 $g(x)\ge \min g(x)=g(1)=0$.

\noindent\textbf{$\star$3.4.15}\quad 若 $p>1$，证明：对于 $[0,1]$ 内任一 $x$，有
\[
  x^p+(1-x)^p\ge\frac1{2^{p-1}}.
\]
（南京邮电大学）

\noindent\textbf{提示}\quad 利用极值法：$f(x)=x^p+(1-x)^p$，问 $f_{\min}=?$

\noindent\textbf{再提示}\quad
\[
f'(x)
\begin{cases}
<0, & x<\dfrac12,\\
=0, & x=\dfrac12,\\
>0, & x>\dfrac12,
\end{cases}
\quad
f_{\min}=f\left(\frac12\right)=\frac1{2^{p-1}}.
\]

\noindent\textbf{$\star$3.4.16}\quad 设 $n$ 为自然数 $0<x<1$，证明：
\[
  x^n(1-x)<\frac1{en}.
\]
（江西师范大学）

\noindent\textbf{提示}\quad
\[
  f(x)\equiv x^n(1-x)\le\max f=f\left(\frac n{n+1}\right)
  =\frac1{\left(1+\dfrac1n\right)^{n+1}}\cdot\frac1n<\frac1{en}.
\]

\noindent\textbf{再提示}\quad $g(x)\equiv\left(1+\dfrac1x\right)^{x+1}$，$g'(x)<0\Rightarrow g(n)=\left(1+\dfrac1n\right)^{n+1}\searrow e\Rightarrow\left(1+\dfrac1n\right)^{n+1}>e\Rightarrow$
\[
  \frac1{\left(1+\dfrac1n\right)^{n+1}}<\frac1e.
\]

\noindent\textbf{$\star$3.4.17}\quad 设 $0<x<1$，试证：
\[
  \sum_{i=1}^\infty x^i(1-x)^{2i}\le\frac4{23}.
\]
（中国科学院）

\noindent\textbf{提示}\quad 可求通项的最大值.

\noindent\textbf{再提示}\quad $f_i(x)=x^i(1-x)^{2i}$，令
\[
  f_i'(x)=ix^{i-1}(1-x)^{2i-1}(1-3x)=0,
\]
得 $(0,1)$ 内唯一可疑点 $x=\dfrac13$，$f_i\left(\dfrac13\right)>f_i(0)=f_i(1)=0$，故 $\max\limits_{0<x<1}f_i(x)=f_i\left(\dfrac13\right)$.

\noindent\textbf{$\star$3.4.18}\quad 求出使得下列不等式对所有自然数 $n$ 都成立的最大的数 $\alpha$ 及最小的数 $\beta$：
\[
  \left(1+\frac1n\right)^{n+\alpha}\le e\le
  \left(1+\frac1n\right)^{n+\beta}.
\]
（中国科学院，北京师范大学）

\noindent\textbf{提示}\quad 见例 3.4.7.

\noindent\textbf{小结}\quad 以上五题主要练习用极值方法证明不等式，以及用单调极限方法获得不等式.

\noindent\textbf{凸函数}

\noindent\textbf{3.4.19}\quad 证明：1）两凸函数之和仍为凸函数；

2）两递增非负凸函数之积仍为凸函数.

\noindent\textbf{提示}\quad 直接用定义 1 可得.

2）欲证
\[
  (f\cdot g)(\lambda x_1+(1-\lambda)x_2)
  \le \lambda(f\cdot g)(x_1)+(1-\lambda)(f\cdot g)(x_2), \tag{1}
\]
而
\[
\begin{aligned}
\text{式 (1) 左端}
&=f(\lambda x_1+(1-\lambda)x_2)\cdot g(\lambda x_1+(1-\lambda)x_2)\\
&\le[\lambda f(x_1)+(1-\lambda)f(x_2)]
  [\lambda g(x_1)+(1-\lambda)g(x_2)]\\
&=\lambda^2f(x_1)g(x_1)+(1-\lambda)^2f(x_2)g(x_2)\\
&\quad+\lambda(1-\lambda)[f(x_1)g(x_2)+f(x_2)g(x_1)].
\end{aligned}\tag{2}
\]
\noindent\textbf{再提示}\quad 式 (1) 右端 $-$ 式 (2) 右端
\[
  =\lambda(1-\lambda)[f(x_1)-f(x_2)][g(x_1)-g(x_2)]\ge0.
\]

\noindent\textbf{$\star$3.4.20}\quad 设 $0<\alpha<1,\ x,y\ge0$，证明：
\[
  x^\alpha y^{1-\alpha}\le \alpha x+(1-\alpha)y.
\]
（华中科技大学）

\noindent\textbf{提示}\quad $x=0$ 或 $y=0$ 时结论自明；$x,y>0$ 时，
\[
  \text{原式}\Leftrightarrow \alpha\ln x+(1-\alpha)\ln y\le\ln[\alpha x+(1-\alpha)y]\Leftrightarrow \ln x\text{ 为凹函数}.
\]
因 $(\ln x)''=-\dfrac1{x^2}<0$，利用定理 5 的推论之对偶结论，结果自明.

\noindent\textbf{3.4.21}\quad 设 $f(x)$ 在 $[a,b]$ 上连续，且 $\forall x_1,x_2\in[a,b],0\le\lambda\le1$，有
\[
  f[\lambda x_1+(1-\lambda)x_2]\ge\lambda f(x_1)+(1-\lambda)f(x_2),
\]
试证：对任何 $T\in(0,b-a)$，必存在 $x_0\in(a,b)$，使 $x_0+T\in[a,b]$，
\[
  \frac{f(x_0+T)-f(x_0)}T=\frac{f(b)-f(a)}{b-a}. \tag{1}
\]
即在 $[a,b]$ 上曲线 $y=f(x)$ 有任意长度（不超过端点弦）平行端点弦的弦.（广西大学）

\noindent\textbf{提示}\quad $f$ 为凹函数，利用定理 3 的 ii）和 iii）之对偶结论（注意凸函数改为凹函数时，不等号反向），可知函数
\[
  g(x)=\frac{f(x+T)-f(x)}{(x+T)-x}-\frac{f(b)-f(a)}{b-a}
  \begin{cases}
  \ge0, & \text{当 }x=a\text{ 时},\\
  \le0, & \text{当 }x=b-T\text{ 时}.
  \end{cases}
\]
故由连续的介值性可得：$\exists x_0\in(a,b)$ 满足式 (1).

\noindent\textbf{$\star$3.4.22}\quad 设 $f(x)$ 在 $[a,b]$ 上满足 $f''(x)>0$，试证：对于 $[a,b]$ 上任意两个不同的点 $x_1,x_2$，有
\[
  \frac12[f(x_1)+f(x_2)]>f\left(\frac{x_1+x_2}2\right). \tag{1}
\]
（陕西师范大学，天津大学等）

\noindent\textbf{提示}\quad $f''(x)>0\Rightarrow f(x)$ 严格凸 $\Rightarrow$ 式 (1) 成立.

或不妨设 $x_1<x_2$，记 $x_0=\dfrac{x_1+x_2}2,\ h=x_2-x_0=x_0-x_1$，则
\[
  f(x_i)=f(x_0)+f'(x_0)(-1)^ih+\frac12f''(\xi_i)h^2\quad(i=1,2),
\]
两式相加可得 (1).

\noindent\textbf{3.4.23}\quad 设 $f(x)$ 是区间 $I$ 上的严格凹函数，即
\[
  f(\lambda x_0+(1-\lambda)x_1)>\lambda f(x_0)+(1-\lambda)f(x_1),\ \forall x_0,x_1\in I,\ \forall\lambda\in(0,1).
\]
试证：若 $f$ 有极大值 $f(x_0)$，则 $f(x_0)$ 必为 $f$ 在 $I$ 上的严格最大值，即 $\forall x\in I$，有 $f(x)<f(x_0)$. 因而 $f$ 的极大值若有必唯一.

\noindent\textbf{提示}\quad 可用反证法.

\noindent\textbf{再提示}\quad 若 $\exists x_1\in I$ 使 $f(x_1)\ge f(x_0)$，则 $\forall\lambda\in(0,1)$，有
\[
  f(\lambda x_0+(1-\lambda)x_1)>
  \lambda f(x_0)+(1-\lambda)f(x_1)\ge f(x_0).
\]
于是在 $x_0$ 的任意 $\delta$ 邻域 $U(x_0,\delta)$ 里 $(0<\delta<|x_1-x_0|)$，只要令 $\lambda:1-\lambda<\dfrac{\delta}{|x_1-x_0|}$，取 $x=\lambda x_0+(1-\lambda)x_1$，则 $x\in U(x_0,\delta)$，但 $f(x)=f(\lambda x_0+(1-\lambda)x_1)>f(x_0)$，与 $f(x_0)$ 为极大值矛盾.

\section*{$\star$§3.5\quad 导数的综合应用}

\subsection*{一、极值问题}

\noindent\textbf{要点}\quad 函数 $f(x)$ 在某点 $x_0$ 有极大（小）值，意指在 $x_0$ 的某邻域里恒有 $f(x)\le f(x_0)$（$f(x)\ge f(x_0)$）（将 $\le(\ge)$ 改为 $< (>)$，则称为严格极值）.

求极值的方法步骤：

1）求可疑点. 可疑点包括：i）稳定点（亦称为驻点或逗留点，皆指一阶导数等于零的点）；ii）导数不存在的点；iii）区间端点.

2）对可疑点进行判断. 基本方法是

i）直接利用定义判断；

ii）利用实际背景来判断；

iii）查看一阶导数的符号，当 $x$ 从左向右穿越可疑点 $x_0$ 时，

若 $f'(x)$ 由“正”变为“负”，则 $f(x_0)$ 为严格极大值；

若 $f'(x)$ 由“负”变为“正”，则 $f(x_0)$ 为严格极小值；

若 $f'(x)$ 不变号，则 $f(x_0)$ 不是极值；

iv）若 $f'(x_0)=0,\ f''(x_0)
\begin{cases}
>0, & \text{则 }f(x_0)\text{ 为严格极小值},\\
<0, & \text{则 }f(x_0)\text{ 为严格极大值};
\end{cases}$

v）$f^{(k)}(x_0)=0\ (k=1,2,\cdots,n-1),\ f^{(n)}(x_0)\ne0$，

若 $n$ 为偶数，则 $f(x_0)$ 为极值：
\[
\begin{cases}
f^{(n)}(x_0)>0\text{ 时为严格极小值},\\
f^{(n)}(x_0)<0\text{ 时为严格极大值};
\end{cases}
\]

若 $n$ 为奇数，则 $f(x_0)$ 不是极值.

所谓最值，指最大、最小值. 它要求极值定义里的不等式在整个定义域里统统成立，也可以说最值是整体极值. 相对而言，前面讲的极值是局部极值. 显然内部最值必为极值，反之未必. 求最值时，有时为了省事，在求出可疑点之后，不判断极大、极小值，可将所有可疑点的值都拿来比较，其中最大、最小者就是整体最大、最小值.

\noindent\textbf{例 3.5.1}\quad 讨论在指定点处函数 $f(x)$ 的极值：

1）若 $\lim\limits_{x\to a}\dfrac{f(x)-f(a)}{(x-a)^2}=-1$，在点 $x=a$ 处；（数学一）

2）若 $f(0)=0,\ f(x)$ 在 $x=0$ 的某邻域内连续，$\lim\limits_{x\to0}\dfrac{f(x)}{1-\cos x}=2$，在 $x=0$ 处；（数学一）

3）设 $f(x)$ 有二阶连续导数，且 $f'(0)=0,\ \lim\limits_{x\to0}\dfrac{f''(x)}{|x|}=1$，在 $x=0$ 处.（数学一）

\noindent\textbf{提示}\quad 利用极限之保号性及极值的定义，知 1）中 $f(a)$ 为极大值，2）中 $f(0)$ 为极小值.

3）$\exists\delta>0$，在 $U_0(0,\delta)$ 内 $\dfrac{f''(x)}{|x|}>0$，从而 $f''(x)>0,\ f'(x)$ 严格 $\nearrow$. 而 $f'(0)=0$，故 $x$ 从左向右穿过 0 点时，$f'(x)$ 由负变正，$f(0)$ 为 $f$ 的极小值.

\noindent\textbf{练习}\quad 求函数 $f(x)=|x(x^2-1)|$ 的极值以及 $[-2,2]$ 上的最大、最小值.

\noindent\textbf{解}\quad
\[
  f(x)=[\operatorname{sgn}(x(x^2-1))]\cdot(x(x^2-1)),
\]
\[
  f'(x)=[\operatorname{sgn}(x(x^2-1))]\cdot(3x^2-1)\quad(x\ne0,\pm1\text{ 时}).
\]
令 $f'(x)=0$，得 $x=\pm\dfrac{\sqrt3}{3}$.
\[
  f''\left(\frac{\sqrt3}{3}\right)
  =[\operatorname{sgn}(x(x^2-1))]\bigg|_{x=\frac{\sqrt3}{3}}\cdot6x\bigg|_{x=\frac{\sqrt3}{3}}
  =-2\sqrt3<0.
\]
故 $f$ 在 $\dfrac{\sqrt3}{3}$ 处取极大值.

因 $f$ 为偶函数，在 $-\dfrac{\sqrt3}{3}$ 处亦为极大值.

当 $x=0,\pm1$ 时，$f(0)=f(1)=f(-1)=0\le f(x)\ (\forall x)$，故为最小值，自然也是极小值.

$f(2)=f(-2)=6\ge f(x)\ (\forall x\in[-2,2])$，故 $f(x)$ 在 $[-2,2]$ 上最大值为 6，最小值为 0.

\noindent\textbf{例 3.5.2}\quad 1）设 $f(x)$ 是 $\mathbb R$ 上的连续函数，试证：若 $\mathbb R$ 的每点都为 $f(x)$ 的极值点，则 $f(x)$ 必为常值函数；（华东师范大学）

2）设 $f(x)$ 是 $\mathbb R$ 上的实值函数，试证：$f(x)$ 在 $\mathbb R$ 上的极值点最多是可数集；（华东师范大学）

3）证明：连续函数（无常值区间）严格极大、严格极小值必然相间出现，删去连续性可以不成立；

4）下面的论断是否正确，说明理由：

（a）函数在有限区间 $[a,b]$ 上的严格极值点，最多只有有限个；

（b）严格极值点必是孤立的.

\noindent\textbf{证}\quad 1）（反证法）非常值函数不可能处处有极值. 事实上，若 $f(x)\ne$ 常数，$\exists a,b$ 使得 $f(a)\ne f(b)$，不妨设 $f(a)<f(b)$（对 $f(a)>f(b)$ 类似可证）. 利用 $f(x)$ 连续的介值性，容易找出另外两点 $a_1,b_1$，使得 $a<a_1<b_1<b,\ b_1-a_1<\dfrac12(b-a)$，且 $f(a)<f(a_1)<f(b_1)<f(b)$.（这不难做到：可先在 $a,b$ 之间找出一点 $\eta:a<\eta<b$ 使得 $f(a)<f(\eta)<f(b)$，再在 $(a,\eta),(\eta,b)$ 中分别找出 $a_1$ 和 $b_1$，使得 $f(a)<f(a_1)<f(\eta)<f(b_1)<f(b)$.）对 $a_1,b_1$ 重复此项工作，如此不断地做下去，我们就可得到一个区间套：
\[
  [a,b]\supset[a_1,b_1]\supset[a_2,b_2]\supset\cdots\supset[a_n,b_n]\supset\cdots,\quad
  0\le b_n-a_n\le\frac1{2^n}(b-a).
\]
利用区间套定理，存在唯一公共点 $\xi\in(a_n,b_n)\ (n=1,2,\cdots)$，这时，$a_n\nearrow\xi$，且 $f(a_n)\nearrow f(\xi)$；$b_n\searrow\xi$，且 $f(b_n)\searrow f(\xi)$.（当 $n$ 充分大时）在点 $\xi$ 的任何邻域里，既有比 $f(\xi)$ 严格大的 $f(b_n)$，又有比 $f(\xi)$ 严格小的 $f(a_n)$. 至此，我们利用介值性，在 $f(x)$ 异值的两点间，总能找到非极值点 $\xi\in\mathbb R$. “处处皆为极限点”的论断不成立，除非 $f(x)$ 是常值函数.

2）$f(x)$ 在每一个区间 $[n,n+1]$ 上的极值点最多只有有限个 $(n=0,\pm1,\pm2,\cdots)$. 这是因为：若 $x_0$ 是 $f(x)$ 的极大值点，则可找出一个尽可能大的邻域 $(\alpha_0,\beta_0)$，使得此点是 $f(x)$ 在该邻域内的（严格）最大值点. 为此只需令
\[
  \alpha_0=\inf\{t\mid t<x_0,\ f(t)<f(x_0),\text{ 但 }t\text{ 不是 }f\text{ 的极值点}\},
\]
\[
  \beta_0=\sup\{t\mid t>x_0,\ f(t)<f(x_0),\text{ 但 }t\text{ 不是 }f\text{ 的极值点}\}.
\]
此时 $\alpha_0$ 和 $\beta_0$ 必是 $f(x)$ 的另外两个极值点. 类似地，对 $f$ 的极小值点也可找出尽可能大的邻域 $(\alpha,\beta)$，使得 $f(x)$ 是邻域里的（严格）最小值.

所有这些开邻域组成了 $[n,n+1]$ 的开覆盖，应用有限覆盖定理可以找出有限子覆盖，可见每个 $[n,n+1]$ 上极值点最多只有有限个 $(n=0,\pm1,\pm2,\cdots)$.

注意到 $\mathbb R=[0,1]\cup[-1,0]\cup[1,2]\cup[-2,-1]\cup\cdots$，知 $f(x)$ 在 $\mathbb R$ 上的极值点最多只是可数集.

3）（极大与极小相间问题）设 $a$ 和 $b$ 是严格极大值点，（不论间距多么小）在 $a,b$ 之间必有严格极小值点. 这是因为

若 $f$ 连续（且无常值区间），则 $f$ 在 $[a,b]$ 上应有（严格）最小值点 $\xi$. 既然 $a,b$ 是 $f$ 的极大值点，那么最小值点 $\xi$ 只能出现在 $(a,b)$ 内部：$a<\xi<b$，内点是最小值点，当然是极小值点. 此即“任意两个严格极大值点 $a,b$ 之间必有严格极小值点”. 同理可证：两个严格极小值点之间必有严格极大值点.

若 $f(x)$ 不连续，结论可以不成立. 例如：下面的 $f(x)$，只有极小值点 $\left(x=n\cdot\dfrac\pi2\right)$，无极大值点：
\[
  f(x)=
  \begin{cases}
  |\sin x|, & x\ne n\pi+\dfrac\pi2,\\
  0, & x=n\pi+\dfrac\pi2
  \end{cases}
  \quad(n=0,\pm1,\pm2,\cdots).
\]

4）两个论断都不成立，例如在 $[-1,1]$ 上，定义函数：
\[
  f(x)=
  \begin{cases}
  x\sin\dfrac1x+2|x|, & x\ne0,\\
  0, & x=0.
  \end{cases} \tag{1}
\]
此函数是偶函数，关于 $y$ 轴对称. 在 $[-1,1]$ 内其图像无限振荡，有无穷多个极值点，原点是最小值点，当然也是极值点. 在它的任一个小邻域内，总有无穷多个极值点，可见极值点未必是孤立的.

\noindent\textbf{例 3.5.3}\quad 设 $f(x)$ 是二阶连续可导的偶函数，且 $f''(0)\ne0$，问 $x=0$ 是不是极值点？为什么？（吉林大学）

\noindent\textbf{解}\quad 因 $f(x)$ 为偶函数，$f'(x)$ 必为奇函数，由 $f'(0)=-f'(0)$ 知 $f'(0)=0$. 故
\[
  f(x)=f(0)+\frac12f''(\xi)x^2\quad(\xi\text{ 在 }0\text{ 与 }x\text{ 之间}). \tag{1}
\]
又 $f''(0)\ne0$，根据极限保号性，在 $x=0$ 的充分小的邻域里 $f''(x)$（从而 $f''(\xi)$）保持跟 $f''(0)$ 同号，故
\[
  f(x)
  \begin{cases}
  >f(0), & f(0)\text{ 为极小值（当 }f''(0)>0\text{ 时）},\\
  <f(0), & f(0)\text{ 为极大值（当 }f''(0)<0\text{ 时）}.
  \end{cases}
\]

\noindent\textbf{注}\quad $1^\circ$ 类似可证：若 $f(x)$ 有连续的 $n$ 阶导数，
\[
  f^{(k)}(x_0)=0\quad(k=1,2,\cdots,n-1),\quad f^{(n)}(x_0)\ne0,
\]
则

i）当 $n$ 为偶数时，$f$ 有极值：$f^{(n)}(x_0)>0$ 时 $f(x_0)$ 为极小值；$f^{(n)}(x_0)<0$ 时 $f(x_0)$ 为极大值；

ii）当 $n$ 为奇数时，$x=x_0$ 处无极值.

$2^\circ$ 该例中，若将“$f(x)$ 有二阶连续导数”的条件放松为“$f(x)$ 在 $x=0$ 的某个邻域里可导，$x=0$ 处有二阶导数”，其余条件不变，结论仍然成立. 因为上面式 (1) 是带 Lagrange 余项的 Taylor 公式，若改用带 Peano 余项的 Taylor 公式：
\[
  f(x)=f(0)+\frac12f''(0)x^2+o(x^2)
  =f(0)+\frac12[f''(0)+o(1)]x^2,
\]
其中 $o(1)$ 为无穷小量（当 $x\to0$ 时）. 可见 $[f''(0)+o(1)]$ 的符号完全由 $f''(0)$ 决定（当 $x$ 与 0 充分接近时），从而后面的推理保持有效.

同理，注 $1^\circ$ 中“有连续的 $n$ 阶导数”的条件可以放松为“$f$ 在 $x_0$ 的某邻域里有 $n-1$ 阶导数，在 $x_0$ 处有 $n$ 阶导数”，其余条件不变，结论仍然成立.

$3^\circ$ 注意 $1^\circ$ 和 $2^\circ$ 所述极值存在的条件都只是充分的，不是必要的！例如
\[
  f(x)=
  \begin{cases}
  e^{-\frac1{x^2}}, & x\ne0,\\
  0, & x=0
  \end{cases}
\]
在 $x=0$ 处各阶导数皆为 0（见例 3.1.13），在 $x=0$ 处并不满足上述条件，但 $f(x)$ 却明显在 $x=0$ 处有极小值而且是最小值.

\noindent\textbf{练习 1}\quad 已知 $f(0)<0,\ f''(x)>0\ (-\infty<x<+\infty)$. 试证：函数 $\dfrac{f(x)}x$ 在 $(-\infty,0]$ 与 $[0,+\infty)$ 上分别都是严格单调递增的.（复旦大学）

\noindent\textbf{提示}\quad 记 $F(x)=xf'(x)-f(x)$，则 $F(0)=-f(0)>0$，且
\[
  F'(x)=xf''(x)
  \begin{cases}
  >0, & x>0,\\
  <0, & x<0.
  \end{cases}
\]
$F(0)=\min\limits_{x\in\mathbb R}F(x)>0\Rightarrow F(x)>0\ (x\ne0)\Rightarrow \left(\dfrac{f(x)}x\right)'=\dfrac{F(x)}{x^2}>0\Rightarrow \dfrac{f(x)}x$ 严增.

\noindent\textbf{练习 2}\quad 已知 $f(x)$ 在 $[a,b]$ 上有四阶导数，且有点 $\beta\in(a,b)$ 使得 $f''(\beta)=0,\ f^{(4)}(\beta)\ne0$，试证：存在 $(x_1,x_2)\subset(a,b)$，使得
\[
  f(x_1)-f(x_2)=f'(\beta)(x_1-x_2). \tag{1}
\]
（北京大学）

\noindent\textbf{提示}\quad 当 $x_1\ne x_2$ 时，式 (1) 等价于：$f$ 在 $\beta$ 点的切线斜率
\[
  f'(\beta)=\frac{f(x_1)-f(x_2)}{x_1-x_2}.
\]
$f$ 在 $\beta$ 点的切线：
\[
  f(x)=f'(\beta)(x-\beta)+f(\beta).
\]
作辅助函数
\[
  F(x)=f(x)-f'(\beta)(x-\beta)-f(\beta).
\]

\noindent\textbf{再提示}\quad 显然 $F(x)$ 也在 $[a,b]$ 上有四阶导数（如图 3.5.1）：
\[
  F(\beta)=F'(\beta)=0,\quad F''(\beta)=f''(\beta),
\]
\[
  F'''(\beta)=f'''(\beta)=0,\quad F^{(4)}(\beta)=f^{(4)}(\beta)\ne0.
\]
利用 Taylor 公式得
\[
  F(x)=\frac12F''(\beta)(x-\beta)^2+o((x-\beta)^3).
\]

\begin{center}
  \includegraphics[width=0.42\textwidth]{3-5-1.png}

  图 3.5.1
\end{center}

可见

i）若 $F''(\beta)=f''(\beta)>0$（而 $F(\beta)=F'(\beta)=0$），则 $F(x)$ 在 $x=\beta$ 处取严格极小值.（利用介值性）在点 $x=\beta$ 的左（右）邻域，分别存在 $x_1$（和 $x_2$），使得 $F(x_1)=F(x_2)$，即
\[
  f(x_1)-f'(\beta)(x_1-\beta)-f(\beta)
  =f(x_2)-f'(\beta)(x_2-\beta)-f(\beta),
\]
亦即 $f(x_1)-f(x_2)=f'(\beta)(x_1-x_2)$，式 (1) 获证.

ii）若 $F''(\beta)=f''(\beta)<0$（而 $F(\beta)=F'(\beta)=0$），则 $F(x)$ 在 $x=\beta$ 处取严格极大值，同理可得式 (1).

iii）若 $F''(\beta)=f''(\beta)=0$（加之：上面已有 $F'(\beta)=F'''(\beta)=0$），且 $f^{(4)}(\beta)\ne0$，所以不论 $F^{(4)}(\beta)=f^{(4)}(\beta)>0$（或 $<0$），知 $F(x)$ 必在 $x=\beta$ 处取极小值（或极大值），同理可得到式 (1).（证毕.）

\noindent\textbf{$\star$例 3.5.4}\quad 证明：函数
\[
  f(x)=\left(\frac2\pi-1\right)\ln x-\ln2+\ln(1+x)
\]
在 $(0,1)$ 内只有一个零点.（北京大学）

\noindent\textbf{证}\quad 令
\[
  f'(x)=\frac{\frac2\pi\left[x-\left(\frac\pi2-1\right)\right]}{x(1+x)}=0,
\]
得 $x=\dfrac\pi2-1$. 由
\[
f'(x)
\begin{cases}
>0, & \dfrac\pi2-1<x<1,\\
=0, & x=\dfrac\pi2-1,\\
<0, & 0<x<\dfrac\pi2-1,
\end{cases}
\]
知 $x=\dfrac\pi2-1$ 既是极小值点也是 $(0,1)$ 内最小值点，故 $f\left(\dfrac\pi2-1\right)<f(1-0)=0$. 在 $\left(\dfrac\pi2-1,1\right)$ 上 $f'(x)>0,\ f$ 严格 $\nearrow,\ f(0^+)=+\infty>0$；在 $\left(0,\dfrac\pi2-1\right)$ 上 $f'(x)<0,\ f$ 严格 $\searrow$. 因此，$f(x)$ 在 $\left(0,\dfrac\pi2-1\right)$ 内有唯一零点，在 $\left[\dfrac\pi2-1,1\right)$ 上无零点. 故 $f(x)$ 在 $(0,1)$ 只有一个零点.

\noindent\textbf{$\star$例 3.5.5}\quad 设
\[
  f(x)=1-x+\frac{x^2}2-\frac{x^3}3+\cdots+(-1)^n\frac{x^n}{n},
\]
证明：方程 $f(x)=0$ 当 $n$ 为奇数时，恰有一实根；当 $n$ 为偶数时，无实根.（昆明理工大学）

\noindent\textbf{证}\quad
\[
  f(x)=\sum_{i=1}^n(-1)^i\frac{x^i}{i}+1,
\]
\[
  f'(x)=\sum_{i=1}^n(-1)^ix^{i-1}
  =
  \begin{cases}
  \dfrac{-1+(-1)^nx^n}{1+x}, & \text{当 }x\ne-1\text{ 时},\\
  -n, & \text{当 }x=-1\text{ 时}.
  \end{cases}
\]

$1^\circ$ 当 $n=2k+1$ 时，$f'(x)<0,\ f$ 严格 $\searrow$，又 $f(-\infty)=+\infty,\ f(+\infty)=-\infty$，故 $f(x)$ 有唯一实根.

$2^\circ$ 当 $n=2k$ 时，
\[
f'(x)
\begin{cases}
>0, & \text{当 }x>1\text{ 时},\\
=0, & \text{当 }x=1\text{ 时},\\
<0, & \text{当 }x<1\text{ 时}.
\end{cases}
\]
故 $f(1)=\min\limits_{\mathbb R}f(x)$.
\[
  f(x)=\left.1+\sum_{i=1}^{2k}(-1)^i\frac{x^i}{i}\right|_{x=1}
  =(1-1)+\left(\frac12-\frac13\right)+\cdots+
  \left(\frac1{2k-2}-\frac1{2k-1}\right)+\frac1{2k}>0,
\]
故 $\forall x\in\mathbb R,\ f(x)\ge f(1)>0,\ f$ 无实根.

\noindent\textbf{$\star$例 3.5.6}\quad 已知小球半径为 $r$，求其外切圆锥的最小体积.（华中师范大学）

\noindent\textbf{解 I}\quad 如图 3.5.2，记圆锥底圆半径为 $R$，圆锥高为 $h$，圆锥中轴线与母线夹角为 $\theta$，则
\[
  \frac Rh=\tan\theta=\frac{r}{\sqrt{(h-r)^2-r^2}}
  =\frac r{\sqrt{h^2-2rh}},\quad R^2=\frac{r^2h}{h-2r}.
\]
故圆锥体积
\[
  V=\frac13\pi R^2\cdot h=\frac{\pi r^2}3\cdot\frac{h^2}{h-2r}.
\]
令
\[
  \left(\frac{h^2}{h-2r}\right)'_h=\frac{h(h-4r)}{(h-2r)^2}=0,
\]
得 $h=0,\ h=4r$. $0$ 不合题意，故 $h=4r$.
\[
  V'=\frac{\pi r^2}3\cdot\frac{h(h-4r)}{(h-2r)^2},
\]
在 $h=4r$ 邻域里从左到右由负变为正，故
\[
  V_{\min}=\left.\frac{\pi r^2}3\cdot\frac{h^2}{h-2r}\right|_{h=4r}
  =\frac83\pi r^3.
\]

\begin{center}
  \includegraphics[width=0.42\textwidth]{3-5-2.png}

  图 3.5.2
\end{center}

\noindent\textbf{解 II}\quad
\[
  h=\frac r{\sin\theta}+r=\frac{1+\sin\theta}{\sin\theta}\cdot r,\quad
  R=h\cdot\tan\theta,
\]
\[
  V=\frac13\pi(h\tan\theta)^2\cdot h
  =\frac13\pi r^3\frac{(\sin\theta+1)^3}{\sin\theta\cos^2\theta}
  \overset{\text{令 }\sin\theta=x}{=}
  \frac13\pi r^3\frac{(1+x)^2}{x(1-x)}.
\]
令
\[
  \left[\frac{(1+x)^2}{x(1-x)}\right]'=0,
\]
得 $x_1=\dfrac13,\ x_2=-1$（舍去）. 实际背景有最小值，只有一个可疑点，故
\[
  V_{\min}=\frac83\pi r^3\quad\left(\text{这时 }\sin\theta=\frac13\right).
\]

\noindent\textbf{注}\quad 求实际问题之极值：（1）要明确目标函数；（2）必须选择恰当的自变量，使之能最简捷地表达目标函数.

\subsection*{二、导数在几何中的应用（举例）}

\noindent\textbf{例 3.5.7}\quad 1）求曲线 $\Gamma:y=x^2-1\ (x>0)$ 上一点 $P$，作 $\Gamma$ 的切线，与坐标轴交于 $M,N$（如图 3.5.3），试求点 $P$ 坐标使 $\triangle OMN$ 的面积最小；（同济大学）

\noindent\textbf{$\star$2）}\quad 已知星形线 $L$ 的方程：
\[
  x^{2/3}+y^{2/3}=a^{2/3}\quad(a>0).
\]
求其切线与两条坐标轴围成的三角形面积的最大值.（中国科学院）

\noindent\textbf{解}\quad 1）$P(x,y)$ 处切线为（斜率为 $(x^2-1)'=2x$）
\[
  Y=2x(X-x)+y.
\]
令 $X=0$ 和 $Y=0$，可得
\[
  N=(0,y-2x^2),\quad M=\left(x-\frac y{2x},0\right).
\]
$\triangle OMN$ 之面积
\[
  S=\frac14\left(x^3+2x+\frac1x\right).
\]
令 $S_x'=0$，得 $x=\dfrac13\sqrt3$（负值被舍去）. 实际背景知存在最小值，又只有一个可疑点，故
\[
  S_{\min}=\frac49\sqrt3\quad\left(\text{此时 }P\text{ 的坐标为 }\left(\frac13\sqrt3,-\frac23\right)\right).
\]

\begin{center}
  \includegraphics[width=0.42\textwidth]{3-5-3.png}

  图 3.5.3
\end{center}

2）\noindent\textbf{提示}\quad 点 $(x_0,y_0)\in L$ 处，切线方程为
\[
  x_0^{-1/3}x+y_0^{-1/3}y=a^{2/3}. \tag{1}
\]
下面说明如何导出式 (1).

\noindent\textbf{方法 I}\quad 记 $F(x,y)=x^{2/3}+y^{2/3}-a^{2/3}$，因 $F'_x=\dfrac23x^{-1/3},F'_y=\dfrac23y^{-1/3}$，故 $L$ 在 $(x_0,y_0)$ 处法线的方向数为：$x_0^{-1/3},y_0^{-1/3}$.（方向数里公共因子 $\dfrac23$ 可以删除.）于是 $L$ 在 $(x_0,y_0)$ 处
的切线方程为
\[
  0=(F'_x,F'_y)\cdot(x-x_0,y-y_0)
  =F'_x\cdot(x-x_0)+F'_y\cdot(y-y_0),
\]
亦即
\[
  0=x_0^{-1/3}x+y_0^{-1/3}y-(x_0^{-1/3}x_0+y_0^{-1/3}y_0)
  =x_0^{-1/3}x+y_0^{-1/3}y-a^{2/3}.
\]
移项即得式 (1).

\noindent\textbf{方法 II}\quad 对 $F(x,y)=0$ 用隐式微分法得
\[
  y'(x_0)=-\left.\frac{F'_x}{F'_y}\right|_{(x_0,y_0)}
  =-\frac{x_0^{-1/3}}{y_0^{-1/3}}.
\]
代入点斜式方程：
\[
  y=y'\big|_{x=x_0}(x-x_0)+y_0,
\]
也可得切线方程 (1).

\noindent\textbf{解法 I}\quad 式 (1) 中，令 $y=0$，得切线在 $x$ 轴上的截距为 $a^{2/3}x_0^{1/3}$；同理可得：在 $y$ 轴上的截距为 $a^{2/3}y_0^{1/3}$. 切线与坐标轴所围三角形面积：
\[
  S_\triangle=\frac12a^{2/3}x_0^{1/3}\cdot a^{2/3}y_0^{1/3}
  =\frac12a^{4/3}x_0^{1/3}y_0^{1/3}. \tag{2}
\]
因此
\[
  \max S_\triangle=\frac12a^{4/3}\max x_0^{1/3}y_0^{1/3}
  \overset{x_0=y_0=\frac a{\sqrt8}}{=}\frac14a^2.
\]
（均值不等式：
\[
  x_0^{1/3}y_0^{1/3}\le\frac12(x_0^{2/3}+y_0^{2/3})=\frac12a^{2/3},
\]
当且仅当 $x_0=y_0$ 时，等号成立，此时乘积达到最大.）

\noindent\textbf{解法 II}\quad 利用星形线 $L$ 的参数方程：$x=a\cos^3\theta,\ y=a\sin^3\theta$，代入 (2)，所求的三角形面积为
\[
  S_\triangle=\frac12a^2\cos\theta\sin\theta
  =\frac14a^2\sin2\theta,
\]
故
\[
  \max S_\triangle=S_\triangle\big|_{\theta=\frac\pi4}=\frac14a^2.
\]

\noindent\textbf{解法 III}\quad （观察法）（仅适用于验算或抢答）如图 3.5.4，用 $M$ 表示曲线 $L$ 与第一象限分角线 $y=x$ 的交点，让切点 $P$ 沿曲线滑动，则切点 $P$ 滑到 $M=\left(\dfrac1{\sqrt8}a,\dfrac1{\sqrt8}a\right)$ 时，面积 $S_\triangle$ 达最大，此时 $x_0=y_0=\dfrac a{\sqrt8}$. 故
\[
  \max S_\triangle=S_\triangle\big|_M=2\left(\frac a{\sqrt8}\right)^2=\frac14a^2.
\]

\begin{center}
  \includegraphics[width=0.62\textwidth]{3-5-4.png}

  图 3.5.4
\end{center}

\noindent\textbf{$\star$例 3.5.8}\quad 求曲线
\[
  x=\frac{3at}{1+t^2},\quad y=\frac{3at^2}{1+t^2}
\]
在 $t=2$ 处的切线方程与法线方程.（中北大学）

\noindent\textbf{提示}\quad 切线方程为
\[
  y=\left.\left(\frac{y'_t}{x'_t}\right)\right|_{t=2}(x-x(2))+y(2),
\]
得
\[
  4x+3y-12a=0.
\]
法线方程为
\[
  y=-\left.\left(\frac{x'_t}{y'_t}\right)\right|_{t=2}(x-x(2))+y(2),
\]
得
\[
  3x-4y+6a=0.
\]

\noindent\textbf{例 3.5.9}\quad 求对数螺线 $\rho=e^\theta$ 在点 $(\rho,\theta)=\left(e^{\pi/2},\dfrac\pi2\right)$ 处切线的直角坐标方程.（数学一）

\noindent\textbf{提示}\quad 可先改写为参数方程：
\[
  \begin{cases}
  x=e^\theta\cos\theta,\\
  y=e^\theta\sin\theta,
  \end{cases}
\]
于是在 $\theta=\dfrac\pi2$ 处切线为
\[
  y=\left.\left(\frac{y'_\theta}{x'_\theta}\right)\right|_{\theta=\frac\pi2}
  \left(x-x\left(\frac\pi2\right)\right)+y\left(\frac\pi2\right),
\]
最后得
\[
  x+y=e^{\pi/2}.
\]

\noindent\textbf{例 3.5.10}\quad 设 $f(x)$ 有连续二阶导数. 已知曲线 $c:y=f(x)$ 在点 $M(x,f(x))$ 处之曲率圆
\[
  (x-a)^2+(y-b)^2=R^2 \tag{1}
\]
跟曲线 $c$ 在点 $M$ 相切（圆心落于凹向的一侧），在点 $M$ 该圆与曲线 $c$ 有相等的一、二阶导数，试求 $a,b,R$ 的表达式.

\noindent\textbf{解}\quad 将 $y$ 看成式 (1) 所确定的 $x$ 的函数. 在式 (1) 两端同时对 $x$ 求导，得
\[
  y'=\frac{a-x}{y-b}. \tag{2}
\]
将式 (2) 再对 $x$ 求导，得 $1+y'^2+(y-b)y''=0$，即
\[
  y''=-\frac{1+y'^2}{y-b}\overset{\text{式 (2)}}{=}
  -\frac{(y-b)^2+(x-a)^2}{(y-b)^3}. \tag{3}
\]
简记 $Z=x-a,\ Y=y-b$，注意到上面的 $y',y''$ 是曲率圆的导数，它们应跟曲线 $c$ 的导数 $f',f''$ 相等.

由式 (2) 得
\[
  Z=-Yf'(x).
\]
由式 (3) 得
\[
  \frac{Z^2+Y^2}{Y^3}=-f''(x).
\]
联立解出 $Z,Y$，最后得
\[
  a=x-\frac{f'(1+f'^2)}{f''},\quad
  b=y+\frac{1+f'^2}{f''},\quad
  R^2=Z^2+Y^2=\frac{(1+f'^2)^3}{f''^2}.
\]

\noindent\textbf{例 3.5.11}\quad 求曲线 $y=\tan x$ 在点 $\left(\dfrac\pi4,1\right)$ 处的曲率圆方程.（山东大学）

\noindent\textbf{提示}\quad 应用上例的结果，算出 $a,b,R$，代入 (1) 可得
\[
  \left(x-\frac{\pi-10}4\right)^2+\left(y-\frac94\right)^2=\frac{125}{16}.
\]

\subsection*{三、导数的实际应用（举例）}

\noindent\textbf{例 3.5.12}\quad 某船从点 $B$ 出发以恒定速度 $v$ 向东航行，观察者在点 $B$ 正南距离为 $s$ 的点 $A$ 进行观察，问视线跟随偏转到分别为多少度时，感觉船速是 $B$ 处船速的 $\dfrac34,\dfrac12,\dfrac14$？

\noindent\textbf{解}\quad 观察者的感觉船速是视线的偏转速度. 如图 3.5.5，用 $\theta$ 表示视线偏转角度，$t$ 表示时间，则 $\theta$ 是时间的函数，且
\[
  s\cdot\tan\theta=v\cdot t.
\]
求导知
\[
  \theta_t'=\frac vs\cos^2\theta.
\]
感觉船速是起始时速度的 $\dfrac34$，即
\[
  \frac vs\cos^2\theta=\frac34\theta'(0)=\frac34\frac vs\cos^20,
\]
解得 $\cos^2\theta=\dfrac34$，故 $\theta=30^\circ$.

类似可得当 $\theta=45^\circ$ 和 $60^\circ$ 时，感觉船速分别是起始时速度的 $\dfrac12$ 和 $\dfrac14$.

\begin{center}
  \includegraphics[width=0.42\textwidth]{3-5-5.png}

  图 3.5.5
\end{center}

\noindent\textbf{$\star$例 3.5.13}\quad 将长度为 $l$ 的均匀细棒放入内空半径为 $a$ 的半球面的杯中，已知 $2a<l<4a$，如不计摩擦力，问什么状况才是平衡位置？

\noindent\textbf{解}\quad 根据物理知识，均匀细棒的重心在中点，平衡位置重心最低.

如图 3.5.6，$AB$ 表示细棒，$M$ 为其中心（即重心），$CD$ 为杯口直径，$\angle CBD=90^\circ,\ ME\perp CD$. 设棒在杯内的部分 $CB$ 之长度为 $x$，于是问题化为：当 $x$ 为多少时，棒的重心最低，即 $EM$ 长度最大？记 $h=EM$ 之长度，$\theta=\angle DCB$. 则
\[
  \frac{h}{x-\frac l2}=\sin\theta=\frac{\sqrt{4a^2-x^2}}{2a},
\]
解得
\[
  h=\frac1{2a}\left(x-\frac l2\right)\sqrt{4a^2-x^2},
\]
求导得
\[
  h'_x=\frac{4a^2-2x^2+\frac l2x}{2a\sqrt{4a^2-x^2}}.
\]
令 $h'_x=0$，在 $h>0$ 范围只有唯一解
\[
  x=\frac l8+\sqrt{\frac{l^2}{64}+2a^2}.
\]
因平衡点客观存在，故可断言杯内长度为此数时重心最低，是平衡位置.

\begin{center}
  \includegraphics[width=0.42\textwidth]{3-5-6.png}

  图 3.5.6
\end{center}

\subsection*{四、导数在求极限中的应用}

见 §1.3 中例 1.3.12 和例 1.3.13.

\subsection*{单元练习 3.5}

\noindent\textbf{3.5.1}\quad 试确定 $a,b,c$，使 $y=x^3+ax^2+bx+c$ 在 $x=1$ 处有拐点，在 $x=0$ 处有极大值 1.（江南大学）

\hfill $\langle a=-3,\ b=0,\ c=1\rangle$

\noindent\textbf{提示}\quad 令 $y''\big|_{x=1}=0,\ y'\big|_{x=0}=0,\ y\big|_{x=0}=1$. 联立求解.

\noindent\textbf{3.5.2}\quad 设
\[
  F(x)=\int_0^x e^{-t}\cos t\,dt,
\]
试求 $F(x)$ 在 $[0,\pi]$ 上的极大值与极小值.（北方交通大学）

\hfill $\left\langle x=\dfrac\pi2\text{ 处取极大值 }F\left(\dfrac\pi2\right)=\dfrac{e^{-\pi/2}+1}{2},\ x=0\text{ 和 }x=\pi\text{ 处取极小值 }F(0)=0,\ F(\pi)=-\dfrac{e^{-\pi}+1}{2}\right\rangle$

\noindent\textbf{$\star$3.5.3}\quad 作函数
\[
  f(x)=|x+2|e^{-1/x}
\]
图.（清华大学）

\noindent\textbf{解}\quad $f(-\infty)=+\infty$.

当 $x<-2$ 时，$f(x)=-(x+2)e^{-1/x}$，这时
\[
  f'(x)=-\frac{e^{-1/x}(x^2+x+2)}{x^2}<0,
\]
$f(x)$ 严格 $\searrow,\ f''(x)>0$，凹向上. $\forall x\ne0$，有 $f(x)\ge0=f(-2)$，故在 $x=-2$ 处取最小值 0.
\[
  f'_-(-2)=\lim_{x\to-2-}\frac{-(x+2)e^{-1/x}-0}{x-(-2)}=-e^{1/2},
\]
\[
  f'_+(-2)=\lim_{x\to-2+}\frac{(x+2)e^{-1/x}-0}{x-(-2)}=e^{1/2}.
\]
故 $x=-2$ 处导数不存在，曲线在此点与 $x$ 轴相交，但不相切.

当 $x>-2$ 时，$f(x)=(x+2)e^{-1/x}\ (x\ne0)$，这时
\[
  f'(x)=\frac{e^{-1/x}(x^2+x+2)}{x^2}>0,
\]
$f$ 严格 $\nearrow$，
\[
  f''(x)=\frac{e^{-1/x}}{x^4}(2-3x).
\]
$x<\dfrac23$ 时 $f''(x)>0$，凹向上；$x>\dfrac23$ 时 $f''(x)<0$，曲线凹向下，故 $x=\dfrac23$ 处有拐点；$x=0$ 处，
\[
  f(0-0)=\lim_{x\to0-}(x+2)e^{-1/x}=+\infty.
\]
以上表明曲线在 $[-2,0)$ 从零单调上升（凹向上）趋向 $+\infty$，以 $y$ 轴为垂直渐近线，
\[
  f(0+0)=\lim_{x\to0+}(x+2)e^{-1/x}=0,\quad f(+\infty)=+\infty,
\]
\[
  \lim_{x\to+\infty}\frac{f(x)}x
  =\lim_{x\to+\infty}\frac{(x+2)e^{-1/x}}x=1,
\]
\[
  \lim_{x\to+\infty}[f(x)-x]
  =\lim_{x\to+\infty}[(x+2)e^{-1/x}-x]
  =\lim_{x\to+\infty}\frac{\left(1+\frac2x\right)e^{-1/x}-1}{\frac1x}
\]
\[
  \overset{t=\frac1x}{=}
  \lim_{t\to0+}\frac{(1+2t)e^{-t}-1}{t}
  \overset{\text{L'Hospital 法则}}{=}
  \lim_{t\to0+}[2e^{-t}-(1+2t)e^{-t}]=1,
\]
因此曲线在 $(0,+\infty)$ 内从 $0^+$ 单调上升（凹向上）至 $x=\dfrac23$ 拐为凹向下，继续上升趋向 $+\infty$，并以 $y=x+1$ 作为斜渐近线. 总之该曲线的图像如图 3.5.7.

\begin{center}
  \includegraphics[width=0.48\textwidth]{3-5-7.png}

  图 3.5.7
\end{center}

\noindent\textbf{注}\quad 该题几乎概括了用导数作图的全部内容，还连带考查了极限. 综合性很强，值得关注. 如果函数还有奇、偶性或周期性，应充分利用，以减少工作量.

\noindent\textbf{3.5.4}\quad 写出下列函数的渐近线：

1）曲线 $y=x\sin\dfrac1x\ (x>0)$；2）曲线
\[
  y=\frac{1+e^{-x^2}}{1-e^{-x^2}}.
\]
（数学一）

\noindent\textbf{提示}\quad 1）为偶函数，其图形关于 $y$ 轴对称，只有水平渐近线 $y=1$（$\lim\limits_{x\to\pm\infty}\dfrac{\sin\frac1x}{\frac1x}=1$）；

2）为偶函数，其图形关于 $y$ 轴对称，有一水平渐近线 $y=1$ 和一竖直渐近线 $x=0$（因 $\lim\limits_{x\to\pm\infty}\dfrac{1+e^{-x^2}}{1-e^{-x^2}}=1,\ \lim\limits_{x\to0}\dfrac{1+e^{-x^2}}{1-e^{-x^2}}=+\infty$）.

\noindent\textbf{3.5.5}\quad 已知一直线切曲线 $y=0.1x^3$ 于 $x=2$，且交此曲线于另一点，求此点坐标.（上海科技大学）

\hfill $\langle(-4,-6.4)\rangle$

\noindent\textbf{提示}\quad 切线为 $y=1.2(x-2)+0.8$，结合 $y=0.1x^3$ 求得交点为 $(-4,-6.4)$.

\noindent\textbf{$\star$3.5.6}\quad 试在一半径为 $R$ 的半圆内作一面积最大的矩形.（山东大学）

\noindent\textbf{提示}\quad 设半圆 $x^2+y^2=R^2\ (y\ge0)$.

\[
  S=2x\sqrt{R^2-x^2},
\]
令 $S'_x=0$ 得
\[
  x=\frac1{\sqrt2}R,\quad y=\frac1{\sqrt2}R.
\]

或
\[
  S=2xy=R^2\sin2\theta,
\]
令 $S'_\theta=0$，得 $\theta=\dfrac\pi4,\ x=\dfrac1{\sqrt2}R,\ y=\dfrac1{\sqrt2}R$.

最大面积矩形 $ABCD$ 为
\[
  A\left(\frac{\sqrt2}2R,\frac{\sqrt2}2R\right),\quad
  B\left(-\frac{\sqrt2}2R,\frac{\sqrt2}2R\right),\quad
  C\left(-\frac{\sqrt2}2R,0\right),\quad
  D\left(\frac{\sqrt2}2R,0\right).
\]

\noindent\textbf{3.5.7}\quad 函数 $f(x)$ 在 $(-\infty,+\infty)$ 内连续，其导函数 $f'(x)$ 的图像如图 3.5.8 所示，问函数 $f(x)$ 有几个极大、极小值点.（数学一）

\noindent\textbf{提示}\quad 图 3.5.8 是导函数的图像，从中可找出四个可疑点，并可根据导数符号的变化判断是极大值点还是极小值点.

\begin{center}
  \includegraphics[width=0.42\textwidth]{3-5-8.png}

  图 3.5.8
\end{center}

\noindent\textbf{3.5.8}\quad 设 $f(x)$ 是 $(-\infty,+\infty)$ 内定义的严格递增函数，$g(x)$ 是某区间 $I$ 上的函数，$x_0\in I$ 为内点（即 $\exists\delta>0$，使得 $U(x_0,\delta)\subset I$），试证：

1）$x=x_0$ 为 $g(x)$ 的极大（极小）值点 $\Leftrightarrow f(g(x))$ 亦以 $x=x_0$ 为极大（极小）值点；

2）函数 $g(x)$ 无极值 $\Rightarrow f(g(x))$ 亦无极值. $f$ 在 $\mathbb R$ 上严格 $\searrow$ 有类似结论.

\noindent\textbf{提示}\quad 因 $f$ 严格 $\nearrow$，$g(x)$ 与 $f(g(x))$ 同增、同减用定义易证.

\noindent\textbf{3.5.9}\quad 设 $g(x),h(x)$ 是某区间 $I$ 上的两函数，$g(x)\ne h(x)$，且 $h\ne0$，试证：只有如下两种可能性：

1）$\dfrac{g(x)}{h(x)}$ 无极值 $\Leftrightarrow \dfrac{g(x)-h(x)}{g(x)+h(x)}$ 亦无极值；

2）$\dfrac{g(x)}{h(x)}$ 与 $\dfrac{g(x)-h(x)}{g(x)+h(x)}$ 有相同的极大、极小值点.

\noindent\textbf{提示}\quad 设 $f(x)=\dfrac{x-1}{x+1}$，则 $f'(x)=\dfrac2{(x+1)^2}>0$，$f(x)$ 严格 $\nearrow$，且
\[
  \frac{g(x)-h(x)}{g(x)+h(x)}=f\left(\frac{g(x)}{h(x)}\right),
\]
然后利用上题.

\noindent\textbf{3.5.10}\quad 证明：函数
\[
  f(x)=\frac{(x-1)(x+2)}{(x+1)(x-2)}
\]
在 $(1,2)$ 内无极值.

\noindent\textbf{提示}\quad 宜用对数导数方法求导：对 $1<x<2$，
\[
  (\ln|f(x)|)'=\frac1{x-1}+\frac1{x+2}-\frac1{x+1}+\frac1{2-x}>0,
\]
$f(x)<0$，故
\[
  f'(x)=\frac{f'(x)}{f(x)}\cdot f(x)=(\ln|f(x)|)'\cdot f(x)<0,
\]
所以 $f$ 在 $(1,2)$ 内无极值.

\noindent\textbf{3.5.11}\quad 设函数 $f(x)$ 在区间 $I$ 上连续，且在 $I$ 上无恒等于常数的子区间，若 $f(x)$ 在 $I$ 上既有极大值又有极小值，试证：其极大、极小值只可能交错地出现，并且每个极大值必比与之相邻的极小值大.

\noindent\textbf{提示}\quad 可用反证法.

设 $x_1<x_2$ 为相邻两极值点，由 $f$ 的连续性，在 $[x_1,x_2]$ 上必有最大、最小值. 如果最大或最小值在 $(x_1,x_2)$ 内部达到，则得出矛盾.

\noindent\textbf{$\star$3.5.12}\quad 一个圆锥面如果沿某一母线剪开，展平，就会得到一个扇形（如图 3.5.9）. 反之，每个扇形可卷成圆锥面. 问半径为 $R$ 的扇形中心角为多大时，卷成的圆锥面体积最大？

\hfill $\left\langle\frac{2\sqrt6}{3}\pi\right\rangle$

\noindent\textbf{提示}\quad $2\pi r=R\alpha$. 体积
\[
  V=\frac13\pi r^2\sqrt{R^2-r^2}.
\]
令 $V'_r=0$，得 $r=\dfrac{\sqrt6}3R$. 故 $\alpha=\dfrac{2\sqrt6}3\pi$.

\begin{center}
  \includegraphics[width=0.36\textwidth]{3-5-9.png}

  图 3.5.9
\end{center}

\noindent\textbf{$\star$3.5.13}\quad 求椭圆
\[
  x^2+\frac{y^2}4=1
\]
在第一象限部分的切线，使它被坐标轴截下的线段最短.

\noindent\textbf{提示}\quad 椭圆上点 $(x,y)$ 处切线方程为
\[
  \frac{Xx}{1^2}+\frac{Yy}{2^2}=1\quad((X,Y)\text{ 为切线上的流动点}).
\]
令 $X=0$（和 $Y=0$），可得切线在 $y$ 轴（和 $x$ 轴）上的截距.

\noindent\textbf{再提示}\quad 截下的线段长为
\[
  l=\sqrt{\frac1{x^2}+\frac{16}{y^2}}\quad(\text{其中 }y^2=4(1-x^2)),
\]
取 $f(x)=l^2=\dfrac1{x^2}+\dfrac4{1-x^2}$ 作为目标函数，令 $f'(x)=0$，得 $x=\dfrac13\sqrt3,\ y=\sqrt{\dfrac83}$. 所求切线方程为
\[
  X+\frac{\sqrt2}{2}Y=\sqrt3.
\]

\end{document}
